% Supplementary Information for
% "Assessment of computational methods for diabetic retinopathy lesion segmentation in fundus images"
%
% Build: make supplement   (tables from data/results.csv via scripts/make_supplement.py, then latexmk)
% The preamble mirrors main.tex so that the table macros and the placeholder colour behave the same way.
\documentclass[10pt,a4paper]{article}

\usepackage[T1]{fontenc}
\usepackage[utf8]{inputenc}
\usepackage{mathptmx}
\usepackage[scaled=0.9]{helvet}
\usepackage{courier}
\usepackage[a4paper,margin=18mm,top=20mm,bottom=22mm]{geometry}
\usepackage{microtype}
\usepackage{amsmath,amssymb}
\usepackage{graphicx}
\usepackage[table]{xcolor}
\usepackage{booktabs,colortbl,multirow,makecell,array,longtable,pdflscape}
\usepackage{threeparttable}
\usepackage{adjustbox}
\usepackage{pifont}
\usepackage[font={small,sf},labelfont=bf,labelsep=space,justification=justified,singlelinecheck=false]{caption}
\usepackage{titlesec}
\usepackage{enumitem}
\usepackage[numbers,super,sort&compress]{natbib}
\usepackage[colorlinks=true,allcolors=linkblue]{hyperref}

\definecolor{linkblue}{HTML}{1C5CAB}
% Shared table macros. \input this after loading:
%   booktabs, colortbl, xcolor (with the "table" option), multirow, makecell,
%   threeparttable, adjustbox, pifont, amssymb, array.
%
% Draft and final. In draft, every number that is still simulated is typeset
% in \phcolor so it cannot be mistaken for a result. \retifinaltrue prints
% all numbers in black.
\newif\ifretifinal
\retifinalfalse

\definecolor{phcolor}{HTML}{5B7FA6}
\definecolor{sdgray}{HTML}{8A8A86}
\definecolor{rulegray}{HTML}{C9C8C0}
\definecolor{headgray}{HTML}{F4F3EF}
\definecolor{inkmuted}{HTML}{6B6B66}
\definecolor{famgeneral}{HTML}{2A78D6}
\definecolor{famauto}{HTML}{1BAF7A}
\definecolor{famfundus}{HTML}{EB6834}
\definecolor{famfoundation}{HTML}{4A3AA7}
\definecolor{tierfine}{HTML}{1D3F5E}
\definecolor{tiermixed}{HTML}{3F6F9C}
\definecolor{tiercoarse}{HTML}{7AA3C7}
\definecolor{tierregion}{HTML}{AABFD3}
\definecolor{deltaup}{HTML}{0B7A3E}
\definecolor{deltadown}{HTML}{B3261E}
\definecolor{barfill}{HTML}{8FB0CC}

% --- numbers
\newcommand{\ph}[1]{\ifretifinal#1\else{\color{phcolor}#1}\fi}          % simulated value
\newcommand{\best}[1]{\textbf{#1}}                                      % best in column
\newcommand{\second}[1]{\underline{#1}}                                 % second best
\newcommand{\sd}[1]{{\scriptsize\color{sdgray}\,$\pm$\,#1}}             % seed s.d.
\newcommand{\up}[1]{{\scriptsize\color{deltaup}$\blacktriangle$\,#1}}   % gain
\newcommand{\down}[1]{{\scriptsize\color{deltadown}$\blacktriangledown$\,#1}} % loss
\newcommand{\same}[1]{{\scriptsize\color{sdgray}#1}}
\newcommand{\na}{{\color{rulegray}--}}
\newcommand{\tbd}{{\color{deltadown}\scriptsize tbd}}

% --- marks
\newcommand{\yes}{{\color{deltaup}\ding{51}}}
\newcommand{\no}{{\color{rulegray}\ding{55}}}
\newcommand{\famchip}[1]{{\color{fam#1}\rule[0.05ex]{0.62em}{0.62em}}}       % family chip
\newcommand{\tier}[1]{{\color{tier#1}\rule[0.05ex]{0.62em}{0.62em}}}     % granularity chip
\newcommand{\tierbadge}[2]{{\setlength{\fboxsep}{1.6pt}\colorbox{tier#1}{\color{white}\scriptsize\sffamily\strut #2}}}
\newcommand{\rlevel}[1]{% reproducibility meter, #1 = 0..4
  \ifcase#1 $\circ\circ\circ\circ$\or $\bullet\circ\circ\circ$\or $\bullet\bullet\circ\circ$\or
  $\bullet\bullet\bullet\circ$\or $\bullet\bullet\bullet\bullet$\fi}
\newcommand{\inlinebar}[1]{{\color{barfill}\rule[0.1ex]{#1em}{0.55em}}}  % #1 = length in em

% --- structure
\newcommand{\grouprow}[2]{\rowcolor{headgray}\multicolumn{#1}{l}{\textbf{#2}}\\}
\newcommand{\colh}[1]{\textbf{#1}}
\newcommand{\tablefont}{\footnotesize\setlength{\tabcolsep}{4.2pt}\renewcommand{\arraystretch}{1.12}}
\newcommand{\phnote}{\ifretifinal\else\item[] {\color{phcolor}Values in this colour are simulated placeholders, not results.}\fi}

\ifdefined\retisegfinal\retifinaltrue\fi
\retifinaltrue   % 2026-10-26: no placeholder row is left in data/results.csv (README); every number printed is measured

\makeatletter
\newcommand{\rn}[1]{\@ifundefined{rn@#1}{\textbf{??#1??}}{\@nameuse{rn@#1}}}
\makeatother
% Generated by scripts/make_numbers.py. Do not edit by hand.
\makeatletter
\@namedef{rn@backbone.best}{HRNet-W48}
\@namedef{rn@backbone.best.ddr}{HRNet-W48}
\@namedef{rn@backbone.span}{0.05--0.15}
\@namedef{rn@bstd.DDR.EX.best}{0.65}
\@namedef{rn@bstd.DDR.EX.median}{0.60}
\@namedef{rn@bstd.DDR.HE.best}{0.56}
\@namedef{rn@bstd.DDR.HE.median}{0.46}
\@namedef{rn@bstd.DDR.MA.best}{0.19}
\@namedef{rn@bstd.DDR.MA.median}{0.15}
\@namedef{rn@bstd.DDR.SE.best}{0.53}
\@namedef{rn@bstd.DDR.SE.median}{0.39}
\@namedef{rn@bstd.DDR.best.model}{SAMed}
\@namedef{rn@bstd.DDR.best.value}{0.46}
\@namedef{rn@bstd.DDR.gap12}{0.010}
\@namedef{rn@bstd.DDR.sd.max}{0.027}
\@namedef{rn@bstd.DDR.sd.median}{0.009}
\@namedef{rn@bstd.DDR.second.model}{RETFound + UPerNet}
\@namedef{rn@bstd.DDR.second.value}{0.45}
\@namedef{rn@bstd.DDR.unet.SE}{0.38}
\@namedef{rn@bstd.DDR.worst.model}{FCT}
\@namedef{rn@bstd.DDR.worst.value}{0.18}
\@namedef{rn@bstd.IDRiD.EX.best}{0.89}
\@namedef{rn@bstd.IDRiD.EX.median}{0.87}
\@namedef{rn@bstd.IDRiD.HE.best}{0.74}
\@namedef{rn@bstd.IDRiD.HE.median}{0.66}
\@namedef{rn@bstd.IDRiD.MA.best}{0.50}
\@namedef{rn@bstd.IDRiD.MA.median}{0.44}
\@namedef{rn@bstd.IDRiD.SE.best}{0.78}
\@namedef{rn@bstd.IDRiD.SE.median}{0.68}
\@namedef{rn@bstd.IDRiD.best.model}{SAMed}
\@namedef{rn@bstd.IDRiD.best.value}{0.71}
\@namedef{rn@bstd.IDRiD.gap12}{0.000}
\@namedef{rn@bstd.IDRiD.sd.max}{0.084}
\@namedef{rn@bstd.IDRiD.sd.median}{0.010}
\@namedef{rn@bstd.IDRiD.second.model}{SAM-Adapter}
\@namedef{rn@bstd.IDRiD.second.value}{0.71}
\@namedef{rn@bstd.IDRiD.unet.SE}{0.68}
\@namedef{rn@bstd.IDRiD.worst.model}{FCT}
\@namedef{rn@bstd.IDRiD.worst.value}{0.30}
\@namedef{rn@cd.chi2}{97.9}
\@namedef{rn@cd.first.model}{SAMed}
\@namedef{rn@cd.first.rank}{1.63}
\@namedef{rn@cd.last.model}{FCT}
\@namedef{rn@cd.last.rank}{17.88}
\@namedef{rn@cd.ntop}{11}
\@namedef{rn@cd.p}{$<$ 0.001}
\@namedef{rn@cd.second.model}{HRNet-W48}
\@namedef{rn@cd.second.rank}{4.13}
\@namedef{rn@cd.third.model}{Swin-Unet}
\@namedef{rn@cd.third.rank}{5.25}
\@namedef{rn@cd.value}{9.90}
\@namedef{rn@ci.DDR.EX.width}{0.14}
\@namedef{rn@ci.DDR.HE.width}{0.19}
\@namedef{rn@ci.DDR.MA.width}{0.07}
\@namedef{rn@ci.DDR.SE.width}{0.24}
\@namedef{rn@ci.DDR.width.hi}{0.19}
\@namedef{rn@ci.DDR.width.lo}{0.07}
\@namedef{rn@ci.IDRiD.EX.width}{0.12}
\@namedef{rn@ci.IDRiD.HE.width}{0.11}
\@namedef{rn@ci.IDRiD.MA.width}{0.08}
\@namedef{rn@ci.IDRiD.SE.width}{0.33}
\@namedef{rn@ci.IDRiD.width.hi}{0.12}
\@namedef{rn@ci.IDRiD.width.lo}{0.08}
\@namedef{rn@ci.hrnet.DDR.HE}{0.477 (0.39--0.59)}
\@namedef{rn@ci.hrnet.DDR.MA}{0.155 (0.12--0.20)}
\@namedef{rn@ci.hrnet.DDR.SE}{0.550 (0.42--0.65)}
\@namedef{rn@ci.hrnet.IDRiD.HE}{0.684 (0.63--0.72)}
\@namedef{rn@ci.hrnet.IDRiD.MA}{0.446 (0.41--0.48)}
\@namedef{rn@ci.hrnet.IDRiD.SE}{0.783 (0.59--0.87)}
\@namedef{rn@ci.nnunet.DDR.HE}{0.501 (0.39--0.57)}
\@namedef{rn@ci.nnunet.DDR.MA}{0.154 (0.11--0.17)}
\@namedef{rn@ci.nnunet.DDR.SE}{0.458 (0.30--0.56)}
\@namedef{rn@ci.nnunet.IDRiD.HE}{0.671 (0.61--0.72)}
\@namedef{rn@ci.nnunet.IDRiD.MA}{0.499 (0.42--0.51)}
\@namedef{rn@ci.nnunet.IDRiD.SE}{0.641 (0.40--0.79)}
\@namedef{rn@ci.retfound.DDR.HE}{0.491 (0.39--0.60)}
\@namedef{rn@ci.retfound.DDR.MA}{0.189 (0.14--0.24)}
\@namedef{rn@ci.retfound.DDR.SE}{0.546 (0.38--0.66)}
\@namedef{rn@ci.retfound.IDRiD.HE}{0.718 (0.66--0.75)}
\@namedef{rn@ci.retfound.IDRiD.MA}{0.484 (0.45--0.52)}
\@namedef{rn@ci.retfound.IDRiD.SE}{0.734 (0.50--0.85)}
\@namedef{rn@ci.samed.DDR.HE}{0.574 (0.52--0.64)}
\@namedef{rn@ci.samed.DDR.MA}{0.173 (0.13--0.22)}
\@namedef{rn@ci.samed.DDR.SE}{0.460 (0.30--0.60)}
\@namedef{rn@ci.samed.IDRiD.HE}{0.736 (0.67--0.78)}
\@namedef{rn@ci.samed.IDRiD.MA}{0.449 (0.40--0.50)}
\@namedef{rn@ci.samed.IDRiD.SE}{0.790 (0.61--0.88)}
\@namedef{rn@clin.ddr.hi.auc}{0.75}
\@namedef{rn@clin.ddr.hi.auc.n}{12}
\@namedef{rn@clin.ddr.hi.maupr}{0.36}
\@namedef{rn@clin.ddr.hi.n}{203}
\@namedef{rn@clin.ddr.lo.auc}{0.54}
\@namedef{rn@clin.ddr.lo.auc.n}{12}
\@namedef{rn@clin.ddr.lo.maupr}{0.21}
\@namedef{rn@clin.ddr.lo.n}{20}
\@namedef{rn@clin.fgadr.hi.auc}{0.76}
\@namedef{rn@clin.fgadr.hi.auc.n}{19}
\@namedef{rn@clin.fgadr.hi.maupr}{0.22}
\@namedef{rn@clin.fgadr.hi.n}{459}
\@namedef{rn@clin.fgadr.lo.auc}{0.71}
\@namedef{rn@clin.fgadr.lo.auc.n}{19}
\@namedef{rn@clin.fgadr.lo.maupr}{0.04}
\@namedef{rn@clin.fgadr.lo.n}{94}
\@namedef{rn@clin.maples.hi.auc}{0.83}
\@namedef{rn@clin.maples.hi.auc.n}{19}
\@namedef{rn@clin.maples.hi.maupr}{0.40}
\@namedef{rn@clin.maples.hi.n}{4}
\@namedef{rn@clin.maples.lo.auc}{0.73}
\@namedef{rn@clin.maples.lo.auc.n}{19}
\@namedef{rn@clin.maples.lo.maupr}{0.33}
\@namedef{rn@clin.maples.lo.n}{56}
\@namedef{rn@clin.retles.hi.auc}{0.83}
\@namedef{rn@clin.retles.hi.auc.n}{19}
\@namedef{rn@clin.retles.hi.maupr}{0.17}
\@namedef{rn@clin.retles.hi.n}{328}
\@namedef{rn@clin.retles.lo.auc}{0.81}
\@namedef{rn@clin.retles.lo.auc.n}{19}
\@namedef{rn@clin.retles.lo.maupr}{0.07}
\@namedef{rn@clin.retles.lo.n}{150}
\@namedef{rn@cost.ddr.max}{72}
\@namedef{rn@cost.ddr.max.model}{FCT}
\@namedef{rn@cost.ddr.ratio}{3.9}
\@namedef{rn@cost.max}{20}
\@namedef{rn@cost.max.model}{FCT}
\@namedef{rn@cost.min}{1.5}
\@namedef{rn@cost.min.model}{HRDecoder}
\@namedef{rn@cross.fgadr.fine}{0.15}
\@namedef{rn@cross.max}{0.28}
\@namedef{rn@cross.mean}{0.16}
\@namedef{rn@e1.m2mrf.DDR.exact}{0.401}
\@namedef{rn@e1.m2mrf.DDR.fn}{0.492}
\@namedef{rn@e1.m2mrf.DDR.gap}{0.091}
\@namedef{rn@e1.m2mrf.DDR.paper}{0.489}
\@namedef{rn@e1.m2mrf.DDR.stored}{0.426}
\@namedef{rn@e1.m2mrf.IDRiD.exact}{0.639}
\@namedef{rn@e1.m2mrf.IDRiD.fn}{0.676}
\@namedef{rn@e1.m2mrf.IDRiD.gap}{0.036}
\@namedef{rn@e1.m2mrf.IDRiD.paper}{0.672}
\@namedef{rn@e1.m2mrf.IDRiD.stored}{0.647}
\@namedef{rn@e1r.HRDecoder.IDRiD.delta}{$-$0.041}
\@namedef{rn@e1r.HRDecoder.IDRiD.exact}{0.674}
\@namedef{rn@e1r.HRDecoder.IDRiD.fn}{0.672}
\@namedef{rn@e1r.M2MRF.IDRiD.delta}{$-$0.060}
\@namedef{rn@e1r.M2MRF.IDRiD.exact}{0.549}
\@namedef{rn@e1r.M2MRF.IDRiD.fn}{0.612}
\@namedef{rn@e1r.first.exact}{0.493}
\@namedef{rn@e1r.first.fn}{0.576}
\@namedef{rn@e4.fgadr}{+0.02}
\@namedef{rn@e4.fgadr383}{+0.01}
\@namedef{rn@e4.fgadronly}{$-$0.16}
\@namedef{rn@e4.ref}{0.56}
\@namedef{rn@e4.retles}{$-$0.04}
\@namedef{rn@e4.retlesonly}{$-$0.35}
\@namedef{rn@e7.aupr.1440}{0.42}
\@namedef{rn@e7.aupr.512}{0.38}
\@namedef{rn@e7.aupr.native}{0.38}
\@namedef{rn@e7.small.1440}{0.28}
\@namedef{rn@e7.small.512}{0.15}
\@namedef{rn@e8.ddr.10}{0.31}
\@namedef{rn@e8.ddr.100}{0.38}
\@namedef{rn@e8.ddr.50}{0.36}
\@namedef{rn@e8.ddr.all}{0.37}
\@namedef{rn@e8.tjdr.100}{0.50}
\@namedef{rn@e8.tjdr.fine}{0.56}
\@namedef{rn@e9.range}{0.14}
\@namedef{rn@family.auto.independent}{0.36}
\@namedef{rn@family.auto.internal}{0.54}
\@namedef{rn@family.foundation.independent}{0.38}
\@namedef{rn@family.foundation.internal}{0.55}
\@namedef{rn@family.fundus.independent}{0.32}
\@namedef{rn@family.fundus.internal}{0.47}
\@namedef{rn@family.general.independent}{0.30}
\@namedef{rn@family.general.internal}{0.47}
\@namedef{rn@fm.any.ref.nnunet.IDRiD}{0.78}
\@namedef{rn@fm.any.sam.DDR}{0.03}
\@namedef{rn@fm.any.sam.IDRiD}{0.07}
\@namedef{rn@fm.any.samtwo.DDR}{0.03}
\@namedef{rn@fm.any.samtwo.IDRiD}{0.08}
\@namedef{rn@fm.box.medsam.DDR}{0.65}
\@namedef{rn@fm.box.medsam.IDRiD}{0.66}
\@namedef{rn@fm.box.sam.DDR}{0.77}
\@namedef{rn@fm.box.sam.IDRiD}{0.86}
\@namedef{rn@fm.box.samtwo.DDR}{0.76}
\@namedef{rn@fm.box.samtwo.IDRiD}{0.85}
\@namedef{rn@fm.boxref.m2mrf.DDR}{0.67}
\@namedef{rn@fm.boxref.m2mrf.IDRiD}{0.81}
\@namedef{rn@fm.boxref.nnunet.DDR}{0.71}
\@namedef{rn@fm.boxref.nnunet.IDRiD}{0.86}
\@namedef{rn@fm.enc.c1.ind}{-0.026}
\@namedef{rn@fm.enc.c1.int}{-0.053}
\@namedef{rn@fm.enc.c1.int.abs}{0.053}
\@namedef{rn@fm.enc.c1.std.abs}{0.045}
\@namedef{rn@fm.enc.c2.ind}{0.006}
\@namedef{rn@fm.enc.c2.int}{-0.020}
\@namedef{rn@fm.enc.c2.int.abs}{0.020}
\@namedef{rn@fm.enc.c2.std.abs}{0.003}
\@namedef{rn@fm.enc.di.TJDR-std}{0.62}
\@namedef{rn@fm.enc.di.TJDR-uwf}{0.14}
\@namedef{rn@fm.enc.di.ind}{0.41}
\@namedef{rn@fm.enc.di.int}{0.51}
\@namedef{rn@fm.enc.dinov2.ind}{0.41}
\@namedef{rn@fm.enc.i.c1.ind}{-0.038}
\@namedef{rn@fm.enc.i.c1.int}{-0.021}
\@namedef{rn@fm.enc.i.c1.int.abs}{0.021}
\@namedef{rn@fm.enc.i.c2.ind}{-0.006}
\@namedef{rn@fm.enc.i.c2.int}{-0.028}
\@namedef{rn@fm.enc.i.c2.int.abs}{0.028}
\@namedef{rn@fm.enc.i.di.MAPLES-DR}{0.44}
\@namedef{rn@fm.enc.i.di.TJDR-std}{0.53}
\@namedef{rn@fm.enc.i.di.TJDR-uwf}{0.25}
\@namedef{rn@fm.enc.i.di.ind}{0.41}
\@namedef{rn@fm.enc.i.di.int}{0.74}
\@namedef{rn@fm.enc.i.mi.MAPLES-DR}{0.43}
\@namedef{rn@fm.enc.i.mi.TJDR-std}{0.54}
\@namedef{rn@fm.enc.i.mi.TJDR-uwf}{0.27}
\@namedef{rn@fm.enc.i.mi.ind}{0.41}
\@namedef{rn@fm.enc.i.mi.int}{0.73}
\@namedef{rn@fm.enc.i.rd.MAPLES-DR}{0.42}
\@namedef{rn@fm.enc.i.rd.TJDR-std}{0.56}
\@namedef{rn@fm.enc.i.rd.TJDR-uwf}{0.22}
\@namedef{rn@fm.enc.i.rd.ind}{0.40}
\@namedef{rn@fm.enc.i.rd.int}{0.71}
\@namedef{rn@fm.enc.i.rm.MAPLES-DR}{0.41}
\@namedef{rn@fm.enc.i.rm.TJDR-std}{0.46}
\@namedef{rn@fm.enc.i.rm.TJDR-uwf}{0.25}
\@namedef{rn@fm.enc.i.rm.ind}{0.37}
\@namedef{rn@fm.enc.i.rm.int}{0.71}
\@namedef{rn@fm.enc.imagenet.ind}{0.41}
\@namedef{rn@fm.enc.imagenet.int}{0.52}
\@namedef{rn@fm.enc.mi.TJDR-std}{0.63}
\@namedef{rn@fm.enc.mi.TJDR-uwf}{0.13}
\@namedef{rn@fm.enc.mi.ind}{0.41}
\@namedef{rn@fm.enc.mi.int}{0.52}
\@namedef{rn@fm.enc.rd.TJDR-std}{0.62}
\@namedef{rn@fm.enc.rd.TJDR-uwf}{0.15}
\@namedef{rn@fm.enc.rd.ind}{0.42}
\@namedef{rn@fm.enc.rd.int}{0.49}
\@namedef{rn@fm.enc.retfound.ind}{0.42}
\@namedef{rn@fm.enc.retfound.int}{0.49}
\@namedef{rn@fm.enc.rm.TJDR-std}{0.59}
\@namedef{rn@fm.enc.rm.TJDR-uwf}{0.11}
\@namedef{rn@fm.enc.rm.ind}{0.39}
\@namedef{rn@fm.enc.rm.int}{0.47}
\@namedef{rn@fm.enc.sets}{the standard-field TJDR set, the ultra-widefield TJDR set and MAPLES-DR}
\@namedef{rn@fm.oracle.all}{0.85}
\@namedef{rn@fm.oracle.five}{0.35}
\@namedef{rn@fm.oracle.one}{0.16}
\@namedef{rn@fm.oracle.twenty}{0.55}
\@namedef{rn@fm.peft.model}{SAMed}
\@namedef{rn@fm.peft.value}{0.46}
\@namedef{rn@fm.point.all.DDR}{0.02}
\@namedef{rn@fm.point.all.IDRiD}{0.12}
\@namedef{rn@fm.ref.nnunet}{0.43}
\@namedef{rn@fm.sam3.EX.IDRiD}{0.41}
\@namedef{rn@fm.sam3.EX.max}{0.41}
\@namedef{rn@fm.sam3.EX.min}{0.03}
\@namedef{rn@fm.sam3.HE.IDRiD}{0.29}
\@namedef{rn@fm.sam3.HE.max}{0.29}
\@namedef{rn@fm.sam3.HE.min}{0.02}
\@namedef{rn@fm.sam3.MA.IDRiD}{0.01}
\@namedef{rn@fm.sam3.MA.max}{0.01}
\@namedef{rn@fm.sam3.MA.min}{0.00}
\@namedef{rn@fm.sam3.SE.IDRiD}{0.21}
\@namedef{rn@fm.sam3.SE.max}{0.21}
\@namedef{rn@fm.sam3.SE.min}{0.00}
\@namedef{rn@fm.wording.ratio}{27.9}
\@namedef{rn@fm.zero.model}{SAM 3}
\@namedef{rn@fm.zero.ub.model}{SAM 2}
\@namedef{rn@fm.zero.ub.value}{0.03}
\@namedef{rn@fm.zero.value}{0.40}
\@namedef{rn@funnel.code}{6}
\@namedef{rn@funnel.dr}{21}
\@namedef{rn@funnel.full}{2}
\@namedef{rn@funnel.train}{4}
\@namedef{rn@human.expert.MA.hi}{0.022}
\@namedef{rn@human.expert.MA.lo}{0.001}
\@namedef{rn@human.expert.auc.hi}{0.93}
\@namedef{rn@human.expert.auc.lo}{0.87}
\@namedef{rn@human.expert.maupr.hi}{0.28}
\@namedef{rn@human.expert.maupr.lo}{0.15}
\@namedef{rn@human.expert.maupr.median}{0.21}
\@namedef{rn@human.models.DDR.auc.median}{0.84}
\@namedef{rn@human.models.DDR.auc.nreach}{7}
\@namedef{rn@human.models.DDR.maupr.median}{0.28}
\@namedef{rn@human.models.IDRiD.auc.median}{0.83}
\@namedef{rn@human.models.IDRiD.auc.nreach}{5}
\@namedef{rn@human.models.IDRiD.maupr.median}{0.25}
\@namedef{rn@human.models.MA.median}{0.039}
\@namedef{rn@human.models.auc.hi}{0.94}
\@namedef{rn@human.models.auc.lo}{0.71}
\@namedef{rn@human.models.below}{3}
\@namedef{rn@human.models.maupr.hi}{0.32}
\@namedef{rn@human.models.maupr.lo}{0.09}
\@namedef{rn@human.models.pairs}{38}
\@namedef{rn@impl.max}{0.060}
\@namedef{rn@impl.mean}{0.005}
\@namedef{rn@impl.moved}{10}
\@namedef{rn@impl.tau}{0.88}
\@namedef{rn@indep.DDR.DiaRetDB1.first.intrank}{10}
\@namedef{rn@indep.DDR.DiaRetDB1.first.model}{DeepLabv3}
\@namedef{rn@indep.DDR.DiaRetDB1.first.value}{0.348}
\@namedef{rn@indep.DDR.DiaRetDB1.mean}{0.251}
\@namedef{rn@indep.DDR.DiaRetDB1.pct}{$-$32}
\@namedef{rn@indep.DDR.DiaRetDB1.rho}{0.78}
\@namedef{rn@indep.DDR.FGADR.drop}{0.16}
\@namedef{rn@indep.DDR.FGADR.first.intrank}{1}
\@namedef{rn@indep.DDR.FGADR.first.model}{SAMed}
\@namedef{rn@indep.DDR.FGADR.first.value}{0.288}
\@namedef{rn@indep.DDR.FGADR.gain}{-0.16}
\@namedef{rn@indep.DDR.FGADR.mean}{0.211}
\@namedef{rn@indep.DDR.FGADR.nlower}{19}
\@namedef{rn@indep.DDR.FGADR.p}{$<$ 0.001}
\@namedef{rn@indep.DDR.FGADR.pct}{$-$43}
\@namedef{rn@indep.DDR.FGADR.rho}{0.79}
\@namedef{rn@indep.DDR.IDRiD.first.intrank}{3}
\@namedef{rn@indep.DDR.IDRiD.first.model}{HRNet-W48}
\@namedef{rn@indep.DDR.IDRiD.first.value}{0.643}
\@namedef{rn@indep.DDR.IDRiD.mean}{0.550}
\@namedef{rn@indep.DDR.IDRiD.pct}{+49}
\@namedef{rn@indep.DDR.IDRiD.rho}{0.92}
\@namedef{rn@indep.DDR.MAPLES-DR.drop}{-0.00}
\@namedef{rn@indep.DDR.MAPLES-DR.first.intrank}{1}
\@namedef{rn@indep.DDR.MAPLES-DR.first.model}{SAMed}
\@namedef{rn@indep.DDR.MAPLES-DR.first.value}{0.470}
\@namedef{rn@indep.DDR.MAPLES-DR.gain}{0.00}
\@namedef{rn@indep.DDR.MAPLES-DR.mean}{0.372}
\@namedef{rn@indep.DDR.MAPLES-DR.nlower}{6}
\@namedef{rn@indep.DDR.MAPLES-DR.p}{= 0.595}
\@namedef{rn@indep.DDR.MAPLES-DR.pct}{+1}
\@namedef{rn@indep.DDR.MAPLES-DR.rho}{0.82}
\@namedef{rn@indep.DDR.Retinal-Lesions.drop}{0.18}
\@namedef{rn@indep.DDR.Retinal-Lesions.first.intrank}{10}
\@namedef{rn@indep.DDR.Retinal-Lesions.first.model}{DeepLabv3}
\@namedef{rn@indep.DDR.Retinal-Lesions.first.value}{0.282}
\@namedef{rn@indep.DDR.Retinal-Lesions.gain}{-0.18}
\@namedef{rn@indep.DDR.Retinal-Lesions.mean}{0.193}
\@namedef{rn@indep.DDR.Retinal-Lesions.nlower}{19}
\@namedef{rn@indep.DDR.Retinal-Lesions.p}{$<$ 0.001}
\@namedef{rn@indep.DDR.Retinal-Lesions.pct}{$-$48}
\@namedef{rn@indep.DDR.Retinal-Lesions.rho}{0.76}
\@namedef{rn@indep.DDR.TJDR-std.drop}{-0.12}
\@namedef{rn@indep.DDR.TJDR-std.first.intrank}{3}
\@namedef{rn@indep.DDR.TJDR-std.first.model}{HRNet-W48}
\@namedef{rn@indep.DDR.TJDR-std.first.value}{0.616}
\@namedef{rn@indep.DDR.TJDR-std.gain}{0.12}
\@namedef{rn@indep.DDR.TJDR-std.mean}{0.490}
\@namedef{rn@indep.DDR.TJDR-std.nlower}{0}
\@namedef{rn@indep.DDR.TJDR-std.p}{$<$ 0.001}
\@namedef{rn@indep.DDR.TJDR-std.pct}{+33}
\@namedef{rn@indep.DDR.TJDR-std.rho}{0.88}
\@namedef{rn@indep.DDR.TJDR-uwf.drop}{0.20}
\@namedef{rn@indep.DDR.TJDR-uwf.first.intrank}{9}
\@namedef{rn@indep.DDR.TJDR-uwf.first.model}{MedSAM (full FT)}
\@namedef{rn@indep.DDR.TJDR-uwf.first.value}{0.243}
\@namedef{rn@indep.DDR.TJDR-uwf.gain}{-0.20}
\@namedef{rn@indep.DDR.TJDR-uwf.mean}{0.173}
\@namedef{rn@indep.DDR.TJDR-uwf.nlower}{19}
\@namedef{rn@indep.DDR.TJDR-uwf.p}{$<$ 0.001}
\@namedef{rn@indep.DDR.TJDR-uwf.pct}{$-$53}
\@namedef{rn@indep.DDR.TJDR-uwf.rho}{0.01}
\@namedef{rn@indep.DDR.fine.first.model}{SAMed}
\@namedef{rn@indep.DDR.fine.first.value}{0.426}
\@namedef{rn@indep.DDR.fine.hrnet}{0.412}
\@namedef{rn@indep.DDR.int.mean}{0.369}
\@namedef{rn@indep.DDR.ndrop}{14}
\@namedef{rn@indep.DDR.rho}{0.80}
\@namedef{rn@indep.DDR.sametop}{the same model}
\@namedef{rn@indep.DDR.top.model}{SAMed}
\@namedef{rn@indep.DDR.top.value}{0.43}
\@namedef{rn@indep.IDRiD.DDR.first.intrank}{1}
\@namedef{rn@indep.IDRiD.DDR.first.model}{SAMed}
\@namedef{rn@indep.IDRiD.DDR.first.value}{0.368}
\@namedef{rn@indep.IDRiD.DDR.mean}{0.286}
\@namedef{rn@indep.IDRiD.DDR.pct}{$-$55}
\@namedef{rn@indep.IDRiD.DDR.rho}{0.75}
\@namedef{rn@indep.IDRiD.DiaRetDB1.first.intrank}{18}
\@namedef{rn@indep.IDRiD.DiaRetDB1.first.model}{DeepLabv3}
\@namedef{rn@indep.IDRiD.DiaRetDB1.first.value}{0.311}
\@namedef{rn@indep.IDRiD.DiaRetDB1.mean}{0.239}
\@namedef{rn@indep.IDRiD.DiaRetDB1.pct}{$-$62}
\@namedef{rn@indep.IDRiD.DiaRetDB1.rho}{0.48}
\@namedef{rn@indep.IDRiD.FGADR.drop}{0.43}
\@namedef{rn@indep.IDRiD.FGADR.first.intrank}{1}
\@namedef{rn@indep.IDRiD.FGADR.first.model}{SAMed}
\@namedef{rn@indep.IDRiD.FGADR.first.value}{0.280}
\@namedef{rn@indep.IDRiD.FGADR.gain}{-0.43}
\@namedef{rn@indep.IDRiD.FGADR.mean}{0.196}
\@namedef{rn@indep.IDRiD.FGADR.nlower}{19}
\@namedef{rn@indep.IDRiD.FGADR.p}{$<$ 0.001}
\@namedef{rn@indep.IDRiD.FGADR.pct}{$-$69}
\@namedef{rn@indep.IDRiD.FGADR.rho}{0.53}
\@namedef{rn@indep.IDRiD.MAPLES-DR.drop}{0.28}
\@namedef{rn@indep.IDRiD.MAPLES-DR.first.intrank}{10}
\@namedef{rn@indep.IDRiD.MAPLES-DR.first.model}{TC-Net}
\@namedef{rn@indep.IDRiD.MAPLES-DR.first.value}{0.412}
\@namedef{rn@indep.IDRiD.MAPLES-DR.gain}{-0.28}
\@namedef{rn@indep.IDRiD.MAPLES-DR.mean}{0.350}
\@namedef{rn@indep.IDRiD.MAPLES-DR.nlower}{19}
\@namedef{rn@indep.IDRiD.MAPLES-DR.p}{$<$ 0.001}
\@namedef{rn@indep.IDRiD.MAPLES-DR.pct}{$-$44}
\@namedef{rn@indep.IDRiD.MAPLES-DR.rho}{0.71}
\@namedef{rn@indep.IDRiD.Retinal-Lesions.drop}{0.46}
\@namedef{rn@indep.IDRiD.Retinal-Lesions.first.intrank}{18}
\@namedef{rn@indep.IDRiD.Retinal-Lesions.first.model}{DeepLabv3}
\@namedef{rn@indep.IDRiD.Retinal-Lesions.first.value}{0.260}
\@namedef{rn@indep.IDRiD.Retinal-Lesions.gain}{-0.46}
\@namedef{rn@indep.IDRiD.Retinal-Lesions.mean}{0.172}
\@namedef{rn@indep.IDRiD.Retinal-Lesions.nlower}{19}
\@namedef{rn@indep.IDRiD.Retinal-Lesions.p}{$<$ 0.001}
\@namedef{rn@indep.IDRiD.Retinal-Lesions.pct}{$-$73}
\@namedef{rn@indep.IDRiD.Retinal-Lesions.rho}{0.52}
\@namedef{rn@indep.IDRiD.TJDR-std.drop}{0.21}
\@namedef{rn@indep.IDRiD.TJDR-std.first.intrank}{1}
\@namedef{rn@indep.IDRiD.TJDR-std.first.model}{SAMed}
\@namedef{rn@indep.IDRiD.TJDR-std.first.value}{0.534}
\@namedef{rn@indep.IDRiD.TJDR-std.gain}{-0.21}
\@namedef{rn@indep.IDRiD.TJDR-std.mean}{0.424}
\@namedef{rn@indep.IDRiD.TJDR-std.nlower}{19}
\@namedef{rn@indep.IDRiD.TJDR-std.p}{$<$ 0.001}
\@namedef{rn@indep.IDRiD.TJDR-std.pct}{$-$33}
\@namedef{rn@indep.IDRiD.TJDR-std.rho}{0.51}
\@namedef{rn@indep.IDRiD.TJDR-uwf.drop}{0.42}
\@namedef{rn@indep.IDRiD.TJDR-uwf.first.intrank}{14}
\@namedef{rn@indep.IDRiD.TJDR-uwf.first.model}{MedSAM (full FT)}
\@namedef{rn@indep.IDRiD.TJDR-uwf.first.value}{0.295}
\@namedef{rn@indep.IDRiD.TJDR-uwf.gain}{-0.42}
\@namedef{rn@indep.IDRiD.TJDR-uwf.mean}{0.209}
\@namedef{rn@indep.IDRiD.TJDR-uwf.nlower}{19}
\@namedef{rn@indep.IDRiD.TJDR-uwf.p}{$<$ 0.001}
\@namedef{rn@indep.IDRiD.TJDR-uwf.pct}{$-$67}
\@namedef{rn@indep.IDRiD.TJDR-uwf.rho}{0.35}
\@namedef{rn@indep.IDRiD.fine.first.model}{SAMed}
\@namedef{rn@indep.IDRiD.fine.first.value}{0.394}
\@namedef{rn@indep.IDRiD.fine.hrnet}{0.367}
\@namedef{rn@indep.IDRiD.int.mean}{0.630}
\@namedef{rn@indep.IDRiD.ndrop}{19}
\@namedef{rn@indep.IDRiD.rho}{0.65}
\@namedef{rn@indep.IDRiD.sametop}{the same model}
\@namedef{rn@indep.IDRiD.top.model}{SAMed}
\@namedef{rn@indep.IDRiD.top.value}{0.39}
\@namedef{rn@indep.pairs.DiaRetDB1.ddrbetter}{13}
\@namedef{rn@indep.pairs.FGADR.ddrbetter}{14}
\@namedef{rn@indep.pairs.MAPLES-DR.ddrbetter}{13}
\@namedef{rn@indep.pairs.Retinal-Lesions.ddrbetter}{15}
\@namedef{rn@indep.pairs.TJDR-std.ddrbetter}{15}
\@namedef{rn@indep.pairs.TJDR-uwf.ddrbetter}{4}
\@namedef{rn@indep.pairs.ddrbetter}{74}
\@namedef{rn@indep.pairs.n}{114}
\@namedef{rn@indep.pairs.std.hi}{15}
\@namedef{rn@indep.pairs.std.lo}{13}
\@namedef{rn@lit.base.copied1}{22}
\@namedef{rn@lit.base.copied2}{26}
\@namedef{rn@lit.base.n}{1,538}
\@namedef{rn@lit.base.rerun}{28}
\@namedef{rn@lit.base.unverified}{23}
\@namedef{rn@lit.code.yes}{8}
\@namedef{rn@lit.ddr.n}{24}
\@namedef{rn@lit.ddr.notstated}{2}
\@namedef{rn@lit.ddr.official}{20}
\@namedef{rn@lit.ddr.own}{2}
\@namedef{rn@lit.deeplab.hi}{0.67}
\@namedef{rn@lit.deeplab.lo}{0.48}
\@namedef{rn@lit.deeplab.n}{11}
\@namedef{rn@lit.fov.notstated}{30}
\@namedef{rn@lit.hed.hi}{0.64}
\@namedef{rn@lit.hed.lo}{0.37}
\@namedef{rn@lit.hed.n}{7}
\@namedef{rn@lit.idrid.n}{30}
\@namedef{rn@lit.idrid.notstated}{3}
\@namedef{rn@lit.idrid.official}{27}
\@namedef{rn@lit.idrid.own}{0}
\@namedef{rn@lit.impl.exact}{2}
\@namedef{rn@lit.impl.fixed}{6}
\@namedef{rn@lit.impl.n}{27}
\@namedef{rn@lit.impl.notstated}{18}
\@namedef{rn@lit.impl.perimage}{1}
\@namedef{rn@lit.lseg.hi}{0.66}
\@namedef{rn@lit.lseg.lo}{0.40}
\@namedef{rn@lit.lseg.n}{14}
\@namedef{rn@lit.m2mrf.hi}{0.67}
\@namedef{rn@lit.m2mrf.lo}{0.58}
\@namedef{rn@lit.m2mrf.n}{6}
\@namedef{rn@lit.n}{30}
\@namedef{rn@lit.n.methods}{29}
\@namedef{rn@lit.pooled.stated}{5}
\@namedef{rn@lit.published.multi}{7}
\@namedef{rn@lit.rerun.all}{5}
\@namedef{rn@lit.rerun.copied}{3}
\@namedef{rn@lit.rerun.notstated}{10}
\@namedef{rn@lit.rerun.some}{12}
\@namedef{rn@lit.unet.hi}{0.65}
\@namedef{rn@lit.unet.lo}{0.35}
\@namedef{rn@lit.unet.n}{6}
\@namedef{rn@n.images}{5,121}
\@namedef{rn@n.independent}{6}
\@namedef{rn@n.main}{19}
\@namedef{rn@n.models}{39}
\@namedef{rn@n.ranked}{19}
\@namedef{rn@n.sources}{8}
\@namedef{rn@n.tier1}{22}
\@namedef{rn@n.tier2}{6}
\@namedef{rn@n.tier3}{11}
\@namedef{rn@params.nnunet}{59.2}
\@namedef{rn@params.resenc}{105.9}
\@namedef{rn@params.resenc.DDR}{140.0}
\@namedef{rn@params.resenc.IDRiD}{105.9}
\@namedef{rn@recall.DDR.large.best}{0.70}
\@namedef{rn@recall.DDR.large.median}{0.66}
\@namedef{rn@recall.DDR.small.best}{0.31}
\@namedef{rn@recall.DDR.small.median}{0.25}
\@namedef{rn@ref.EX.delta.mean}{-0.001}
\@namedef{rn@ref.EX.delta.min}{-0.002}
\@namedef{rn@ref.HE.delta.mean}{-0.000}
\@namedef{rn@ref.HE.delta.min}{-0.001}
\@namedef{rn@ref.MA.delta.mean}{-0.002}
\@namedef{rn@ref.MA.delta.min}{-0.004}
\@namedef{rn@ref.SE.delta.mean}{-0.000}
\@namedef{rn@ref.SE.delta.min}{-0.001}
\@namedef{rn@ref.delta.max}{0.000}
\@namedef{rn@ref.delta.mean}{-0.001}
\@namedef{rn@ref.delta.min}{-0.002}
\@namedef{rn@ref.geometry.absmax}{0.006}
\@namedef{rn@ref.geometry.tau}{0.95}
\@namedef{rn@ref.n}{19}
\@namedef{rn@ref.tau}{1.00}
\@namedef{rn@tau.published}{-0.11}
\@namedef{rn@tau.published.n}{10}
\@namedef{rn@tau.published.p}{0.73}
\@namedef{rn@vanish.DDR.1440.EX}{0.1}
\@namedef{rn@vanish.DDR.1440.MA}{0.0}
\@namedef{rn@vanish.DDR.224.EX}{35.3}
\@namedef{rn@vanish.DDR.224.MA}{34.3}
\@namedef{rn@vanish.DDR.384.EX}{16.9}
\@namedef{rn@vanish.DDR.384.MA}{11.5}
\@namedef{rn@vanish.IDRiD.1440.EX}{0.0}
\@namedef{rn@vanish.IDRiD.1440.MA}{0.0}
\@namedef{rn@vanish.IDRiD.224.EX}{25.8}
\@namedef{rn@vanish.IDRiD.224.MA}{16.0}
\@namedef{rn@vanish.IDRiD.384.EX}{5.7}
\@namedef{rn@vanish.IDRiD.384.MA}{1.1}
\makeatother


\titleformat{\section}{\sffamily\bfseries\Large}{}{0pt}{}
\titleformat{\subsection}{\sffamily\bfseries\normalsize}{}{0pt}{}
\titleformat{\paragraph}[runin]{\sffamily\bfseries\small}{}{0pt}{}[.]
\titlespacing*{\section}{0pt}{14pt}{4pt}
\titlespacing*{\subsection}{0pt}{10pt}{2pt}
\setlength{\parskip}{0pt}
\setlength{\parindent}{1.2em}
\captionsetup[table]{name=Supplementary Table}
\renewcommand{\tablename}{Supplementary Table}
\newcommand{\mAUPR}{mAUPR}

\begin{document}

\begin{center}
{\sffamily\bfseries\LARGE Supplementary Information}\\[6pt]
{\sffamily\large Assessment of computational methods for diabetic retinopathy lesion segmentation in fundus images}\\[4pt]
{\small Author One, Author Two, \ldots, Corresponding Author}
\end{center}

\ifretifinal\else
\begin{center}\small
\fbox{\parbox{0.9\textwidth}{\textbf{Draft.} Numbers printed in {\color{phcolor}this colour} are placeholders and not results; every other number is measured. Tables are generated from the same data table as the main text (\texttt{data/results.csv}) by \texttt{scripts/make\_supplement.py}.}}
\end{center}
\fi

\tableofcontents

\section{Supplementary Notes}

\subsection*{Supplementary Note 1: the registered model variants}
One entry per registered variant: family, tier (1, main comparisons; 2, as time allows; 3, paper audit only), reproducibility level (R0 no code; R1 model definition only; R2 inference code and weights; R3 training code; R4 training code, weights, splits and environment), track (A, released weights; B, unified retraining of the official code; C, re-implementation from the paper), the code used, and what was changed or found when the code was run. The registry is \texttt{bench/configs/models.yaml} in the repository; the notes are those recorded at the code audit and updated as runs completed.
\paragraph{General segmentation networks.}
\begin{itemize}[leftmargin=*,itemsep=1pt]
\item \textbf{U-Net} (tier 1, reproducibility R3, track B); no author repo; use a reference PyTorch implementation.
\item \textbf{DeepLabv3} (tier 1, reproducibility R3, track B); code: \texttt{official\_code/DeepLabv3plus}; registered name is DeepLabv3. torchvision deeplabv3\_resnet101 is ASPP plus bilinear upsampling, not the DeepLabv3+ decoder. The TensorFlow repo is the paper code and is not the B1 network..
\item \textbf{HRNet} (tier 1, reproducibility R3, track B); code: \texttt{official\_code/HRNet-Semantic-Segmentation}.
\item \textbf{nnU-Net} (tier 1, reproducibility R4, track B); code: \texttt{official\_code/nnUNet}; shallow clone 2026-10-04; train entry nnunetv2/run/run\_training.py.
\item \textbf{Swin-Unet} (tier 1, reproducibility R3, track B); code: \texttt{official\_code/Swin-Unet}.
\item \textbf{FCT} (tier 1, reproducibility R3, track B); code: \texttt{official\_code/FCT}; official PyTorch port in PyTorch/main.py, not the TensorFlow entry. B1 adaptation widens the stem to 3 channels, assigns the four head channels to MA HE EX SE, and sets forward\_size 384..
\item \textbf{H2Former} (tier 1, reproducibility R3, track B); code: \texttt{official\_code/H2Former}; general medical segmentation model used on several organs. Not a DR-specific method. B1 uses idrid\_train.py for both IDRiD and DDR..
\item \textbf{UNet++} (tier 2, reproducibility R3, track B); code: \texttt{official\_code/Nested-UNet}.
\item \textbf{UPerNet} (tier 2, reproducibility R3, track B); code: \texttt{official\_code/UPerNet}.
\end{itemize}
\paragraph{Fundus-specific: multi-scale CNNs.}
\begin{itemize}[leftmargin=*,itemsep=1pt]
\item \textbf{M2MRF} (tier 1, reproducibility R4, track AB); code: \texttt{official\_code/M2MRF}; mmseg 0.8 / torch 1.6; weights on Google Drive; README: paper reports mean of 3 runs, code reports best.
\item \textbf{HACDR-Net} (tier 1, reproducibility R3, track B); code: \texttt{official\_code/HACDR-Net}; mmseg 0.30 / mmcv-full 1.6.2; no released weights found.
\item \textbf{HRDecoder} (tier 1, reproducibility R4, track AB); code: \texttt{official\_code/HRDecoder}; added 2026-10-14 (docs/paper/18): MICCAI 2024, mmseg 0.16 / mmcv 1.7, HRNet-W48 backbone; weights on Google Drive (IDRiD 1tRBNqpbFwnNBGQBMw8UYL8sjJzWlI3O9, DDR 1CBTHOtxToHQ13H4SZIUBZTwX0zwmIduh); runs in the retiseg-hacdr environment.
\item \textbf{WSRFNet} (tier 1, reproducibility R1, track AB); code: \texttt{official\_code/WSRFNet}; checked 2026-10-03: model.py + test.py, no train.py; weights by email to xuanli@stu.hit.edu.cn. 2026-10-14: tier 1, the training loop and the paper's loss are in bench/models/wsrfnet.py (docs/paper/18).
\item \textbf{L-Seg} (tier 2, reproducibility R1, track C); code: \texttt{official\_code/L-Seg}; Caffe prototxt and loss layer.
\item \textbf{SAA} (tier 2, reproducibility R0, track C); code: \texttt{official\_code/SAA}; official repo empty.
\item \textbf{CARNet} (tier 3, reproducibility R0).
\item \textbf{PMCNet} (tier 1, reproducibility R3, track B); code: \texttt{official\_code/PMCNet}; train.py present; TF 1.13/Keras 2.3; no released weights; bundled data is one sample image. 2026-10-14: tier 1, rewritten layer by layer in PyTorch (bench/models/pmcnet.py, docs/paper/18); the released dataGenerator one-hot-encodes 3 classes against the 5-class head.
\item \textbf{MLNet} (tier 3, reproducibility R0).
\item \textbf{LezioSeg} (tier 3, reproducibility R0).
\item \textbf{HEDNet-cGAN} (tier 3, reproducibility R0).
\item \textbf{Guo2024} (tier 3, reproducibility R0); Google Drive link of unknown content.
\item \textbf{CLC-Net} (tier 1, reproducibility R0, track C); code: \texttt{reproductions/Wang2023\_CLC-Net}; no author code; re-implemented from the paper by the user (assumptions A1-A22 in the reproduction), B-std through bench/models/clcnet.py (docs/paper/18).
\end{itemize}
\paragraph{Fundus-specific: relational and transformer models.}
\begin{itemize}[leftmargin=*,itemsep=1pt]
\item \textbf{RTNet} (tier 2, reproducibility R0, track C); image + vessel input.
\item \textbf{VTA} (tier 3, reproducibility R0).
\item \textbf{TC-Net} (tier 1, reproducibility R0, track C); code: \texttt{reproductions/Zhang2023\_TC-Net}; no author code; re-implemented from the paper by the user (assumptions in the reproduction), B-std through bench/models/tcnet.py (docs/paper/18).
\end{itemize}
\paragraph{Fundus-specific: semi-supervised.}
\begin{itemize}[leftmargin=*,itemsep=1pt]
\item \textbf{SSMD-UNet} (tier 3, reproducibility R0); author split 70/5/25; EyePACS pretraining overlaps Retinal-Lesions source.
\item \textbf{Zhou2019-CL} (tier 3, reproducibility R0).
\end{itemize}
\paragraph{Foundation models.}
\begin{itemize}[leftmargin=*,itemsep=1pt]
\item \textbf{SAM} (tier 1, reproducibility R2, track B); code: \texttt{official\_code/segment-anything}.
\item \textbf{SAM2} (tier 1, reproducibility R3, track B); code: \texttt{official\_code/sam2}.
\item \textbf{SAM3} (tier 1, reproducibility R3, track B); code: \texttt{official\_code/sam3}; shallow clone 2026-10-04; train entry sam3/train/train.py; checkpoints are gated on Hugging Face.
\item \textbf{MedSAM} (tier 1, reproducibility R3, track B); code: \texttt{official\_code/MedSAM}.
\item \textbf{SAMed} (tier 1, reproducibility R3, track B); code: \texttt{official\_code/SAMed}.
\item \textbf{SAM-Adapter} (tier 1, reproducibility R3, track B); code: \texttt{official\_code/SAM-Adapter}.
\item \textbf{SAM2-UNet} (tier 1, reproducibility R3, track B); code: \texttt{official\_code/SAM2-UNet}.
\item \textbf{RETFound} (tier 1, reproducibility R3, track B); code: \texttt{official\_code/RETFound}; rmaphoh/RETFound\_MAE, shallow clone 2026-10-04; main\_finetune.py is classification. UPerNet head is still to be added. Weights on Hugging Face.
\item \textbf{PerSAM} (tier 2, reproducibility R2, track B); code: \texttt{official\_code/PerSAM}.
\item \textbf{TP-DRSeg} (tier 3, reproducibility R0); code: \texttt{official\_code/TP-DRSeg}; paper link 404; mirror has README only.
\item \textbf{SAM2-DR-Chen2025} (tier 3, reproducibility R0).
\item \textbf{Ophthalmology-SAM} (tier 3, reproducibility R0).
\end{itemize}


\subsection*{Supplementary Note 2: the training budget, the selection rule and the disclosed revisions}
\paragraph{Budget} The training amount of the unified protocol is counted in canvas pixels: one equivalent epoch draws as many patch pixels as the canvases of the training split contain. The number of equivalent epochs was frozen from the validation curves of two models (U-Net and M2MRF) on each source before the first main-table run: 1,000 on IDRiD (44 training images) and 400 on DDR (383). The rule for the freeze was written down first: the smoothed validation curve of both probe models had to reach its maximum before 90\% of the schedule, otherwise the amount would be doubled once and the probe repeated. Both probes met the rule at these values; a later two-fold probe on three cells (HRNet-W48 on DDR, HRDecoder on both sources; main text, case study) is reported as a sensitivity analysis and did not change the budget.
\paragraph{Selection} A run is validated 20 times at equal intervals on the validation split of its training source and the checkpoint with the highest validation \mAUPR{} inside the field of view is scored; the last step is kept for sensitivity. The test split is never read during training. nnU-Net keeps its own schedule and selection.
\paragraph{Revisions after a test number had been read} One: the reference U-Net's source recipe, a binary cross-entropy plus a sigmoid Dice term computed per image with smoothing 1, drove the rare-class channels to zero on DDR. The revision (Dice pooled over the batch per class, smoothing $10^{-5}$) was decided from the training and validation curves of a diagnostic run, after the registered recipe's test numbers had been read; both versions are reported (Supplementary Table~6). Before any foundation model was trained the same rule was registered for SAM2-UNet and MedSAM, whose per-image terms share the mechanism, and it fired for both on DDR; their registered runs are reported next to the revised ones. No budget, augmentation or selection rule was changed after a test number had been read.
\paragraph{Scope changes} The source-recipe track (each model trained with its source recipe on the shared data) was stopped for compute on 2026-10-11 before any of its six started runs finished, and no number from it is reported. The author-recipe reruns were reduced to M2MRF (IDRiD) and HRDecoder (IDRiD); two further M2MRF seeds were added on 2026-10-25 when cards were free. The learning-rate sweep ran on DDR for HRNet-W48 and M2MRF only. The encoder comparison's IDRiD column (three cells) was added after the DDR result of the first pair had been read. The reference model's two cells and their ensemble were registered on 2026-10-23, before any merged run.

\subsection*{Supplementary Note 3: evaluation protocol}
\paragraph{Pixel level} For each class, the positive and negative pixels of all test images inside the field of view (and outside the ignore region of that class) are pooled, and AUPR is computed exactly over the 16-bit probability levels of the stored maps; this equals the average precision over the pooled pixels. The 11-threshold AUPR of the evaluation code released with M2MRF is computed on the same maps and reported next to it. Dice and intersection over union are given at the threshold that maximises the mean Dice on the validation split of the training source (val-opt) and at 0.5.
\paragraph{Lesion level} Lesions are 8-connected components. A ground-truth lesion is detected if any predicted positive pixel overlaps it; a predicted component is a true positive if it overlaps any ground-truth pixel of the class. Lesion size is the equivalent diameter of the component relative to that of the field of view; the small/medium/large boundaries, 0.005 and 0.011, are the medians of microaneurysms and hemorrhages on the training images of IDRiD and DDR, frozen before training.
\paragraph{Intervals and tests} 95\% intervals are percentile intervals of 2,000 bootstrap replicates over test images; replicates pool per-image histograms coarsened to 4,096 bins. Two models are compared by a paired bootstrap with shared image indices and by a two-sided Wilcoxon signed-rank test on per-image Dice with Holm correction. Rankings across the eight test-set blocks use the Friedman test and the Nemenyi critical difference; the DiaRetDB1 block uses the official image-level AUC on DiaRetDB1.
\paragraph{Ultra-widefield images} The TJDR ultra-widefield set is predicted at native resolution with sliding windows for every model trained on a source with a known field of view; for cells whose training sources fix no single field of view (the factor cells trained on TJDR standard-field images, on FGADR or on merged sources), only the version resized to the 1,440-pixel canvas is run and stands for the result on that set.

\clearpage
\section{Supplementary Tables}

\begin{landscape}
{\scriptsize\setlength{\tabcolsep}{3.2pt}\setlength{\LTcapwidth}{\textwidth}\renewcommand{\arraystretch}{1.08}
\begin{longtable}{p{34mm} c p{20mm} p{24mm} p{14mm} p{14mm} p{52mm} p{18mm} p{14mm} p{16mm} c}
\caption{\textbf{Literature audit: the experimental setup stated by each method paper in the ledger.} Split, input resolution, AUPR implementation and the provenance of baseline numbers as extracted from the paper's text and tables (literature/meta/paper\_audit.tsv in the repository holds the verbatim evidence for every cell).}\label{stab:literature}\\
\toprule
Paper & Year & Venue & Datasets & IDRiD split & DDR split & Input resolution & AUPR impl. & Pooling & Baselines & $n$ \\
\midrule
\endfirsthead
\multicolumn{11}{l}{\textit{continued}}\\
\toprule
Paper & Year & Venue & Datasets & IDRiD split & DDR split & Input resolution & AUPR impl. & Pooling & Baselines & $n$ \\
\midrule
\endhead
\midrule
\multicolumn{11}{r}{\textit{continued on the next page}}\\
\endfoot
\bottomrule
\endlastfoot
Guo2019 L-Seg & 2019 & Neurocomputing & IDRiD;e-ophtha;DDR & official & official & IDRiD and e-ophtha: 'resize each image and… & 11-threshold & not stated & rerun all & 4 \\
Zhou2019 collaborative semi s… & 2019 & CVPR & IDRiD & official & NA & 640x640 (Fig. 2 shows 640*640*3 input / 640… & not stated & not stated & not stated & 5 \\
Xiao2020 HEDNet cGAN & 2020 & arXiv (2007.138… & IDRiD & official & NA & original 4288 x 2848; training: 'we randoml… & exact & pooled ('… & rerun all & 2 \\
Bo2022 SAA & 2022 & ISBI & IDRiD;DDR & official & official & IDRiD: 'resized images and their correspond… & not stated & not stated & not stated & 6 \\
Guo2022 CARNet & 2022 & Complex \& Intel… & IDRiD;e-ophtha;DDR & official & official & 'all images are uniformly resized to 1024 x… & NA & not stated & rerun all & 8 \\
Hasan2022 DRG-Net & 2022 & arXiv (2212.146… & IDRiD;FGADR & official & NA & 'For all experiments, the sizes of the inpu… & not stated & not stated & copied & 8 \\
He2022 PMCNet & 2022 & arXiv (2205.157… & DDR;IDRiD;e-ophtha & official & official & IDRiD and E-ophtha: 'resized each image and… & not stated & not stated & rerun some & 11 \\
Huang2022 RTNet & 2022 & IEEE Transactio… & IDRiD;DDR & official & official & training: large images random resized and c… & exact & not stated & rerun some & 7 \\
Ali2023 LezioSeg & 2023 & Electronics (MD… & IDRiD;e-ophtha;DDR & official & official & training: each image split into 4 non-overl… & per-image mean & per-image & not stated & 11 \\
Geetha2023 RILBP-YNet & 2023 & Biomedical Sign… & IDRiD;e-ophtha (plu… & official & NA & black border cropped by contour analysis, t… & 11-threshold & not stated & copied & 11 \\
Liu2023 M2MRF & 2023 & Pattern Recogni… & DDR;IDRiD & official & official & DDR: downsample so long side = 1024, zero-p… & 11-threshold & pooled & rerun some & 12 \\
Qiu2023 Learnable Ophthalmolo… & 2023 & arXiv & IDRiD & not stated & NA & not stated & NA & not stated & not stated & 1 \\
Ullah2023 SSMD-UNet & 2023 & Sci Rep & FGADR;IDRiD & official & NA & 'bicubic interpolation to resize the images… & not stated & not stated & rerun some & 9 \\
Wang2023 CLC-Net & 2023 & Neurocomputing & IDRiD;e-ophtha;DDR & official & official & no global resize: Otsu to remove non-tissue… & not stated & not stated & rerun some & 8 \\
Wang2023 VTA hyperbolic & 2023 & Scientific Repo… & IDRiD;DDR & official & official & not stated (ViT patch size 16; model pre-tr… & not stated & not stated & rerun all & 4 \\
Yin2023 DB-UNet & 2023 & IEEE Access & IDRiD;DDR;e-ophtha… & official & official & IDRiD: ROI 2560x3840 downscaled by t (t=2 f… & 11-threshold & not stated & rerun some & 8 \\
Zhang2023 TC-Net & 2023 & Comput Biol Med & IDRiD;DDR;ISIC-2017 & not stated & own & all images resized to 512x512 & NA & not stated & not stated & 7 \\
Bian2024 MLNet & 2024 & Algorithms (MDP… & DDR;DiaRetDB1;IDRiD & official & official & black borders cropped, images uniformly res… & not stated & not stated & rerun some & 7 \\
Fu2024 LFRC-Net & 2024 & Applied Science… & IDRiD;DDR;Local ult… & official & official & IDRiD: cropping, zero-padding and resizing… & not stated & not stated & rerun some & 12 \\
Guo2024 deep multiscale & 2024 & SSRN preprint (… & IDRiD;DDR;DiaRetDB1… & official & official & three input scales 384x384, 512x512, 1024x1… & not stated & pooled & not stated & 5 \\
Li2024 TP-DRSeg & 2024 & MICCAI & IDRiD;DDR & not stated & not stated & resized to 1024x1024 ('We resize the input… & not stated & not stated & not stated & 8 \\
Li2024 WSRFNet & 2024 & IJCAI & IDRiD;DDR & official & official & resized to 1440x960 (IDRiD) and 1024x1024 (… & 11-threshold & pooled & rerun some & 8 \\
Playout2024 cross dataset gen… & 2024 & arXiv (ISBI-sty… & IDRiD;DDR;FGADR;Ret… & official & official & cropped to remove black background, resized… & 11-threshold & not stated & copied & 3 \\
Xu2024 HACDR-Net & 2024 & AAAI & IDRiD;DDR & official & official & IDRiD resized to 1440x960 (following M2MRF)… & not stated & not stated & not stated & 13 \\
Xu2024 MSAG & 2024 & Scientific Repo… & IDRiD;DDR & official & not stated & IDRiD: 4288x2848 cut into 54 blocks of 512x… & not stated & not stated & not stated & 9 \\
Xu2024 PGAFT & 2024 & Sci Rep & IDRiD;DDR & official & official & 'both the larger color fundus images and va… & not stated & not stated & rerun all & 11 \\
Chen2025 SAM2 DR & 2025 & ISBI & IDRiD;DDR & official & official & resized to 1024x1024 ('Input images are res… & not stated & not stated & not stated & 6 \\
Guo2026 TVFNet & 2026 & Neural Comput A… & IDRiD;DDR;FGADR;e-o… & official & official & 'the fundus images were uniformly resized t… & not stated & not stated & rerun some & 18 \\
Qi2026 TEDR-YNet & 2026 & IEEE Access & IDRiD;DDR;TJDR & official & own & images from all datasets were resized to 51… & not stated & not stated & rerun some & 10 \\
Wang2026 CSIF & 2026 & Applied Soft Co… & DDR;IDRiD & official & official & DDR: long edge scaled to 1024 then zero-pad… & not stated & not stated & rerun some & 11 \\
\end{longtable}
}
\vspace{-6pt}\noindent{\scriptsize ``not stated'' means the paper gives no information. Baselines: how the comparison numbers were obtained (rerun, copied, mixed, not stated). $n$, number of compared methods.}
\end{landscape}
\clearpage

{\scriptsize\setlength{\tabcolsep}{3.2pt}\setlength{\LTcapwidth}{\textwidth}\renewcommand{\arraystretch}{1.08}
\begin{longtable}{l l c c c c c c c}
\caption{\textbf{One method, many reported values.} Every method that appears with a mean AUPR on IDRiD or DDR in at least two papers of the ledger: the number of papers, the lowest and highest value, their spread, and how many of the values were produced by the reporting paper (rerun), carried from an earlier paper (copied) or cite a source outside the ledger (unverifiable).}\label{stab:reported}\\
\toprule
Method & Dataset & Papers & Min & Max & Spread & Rerun & Copied & Unverif. \\
\midrule
\endfirsthead
\multicolumn{9}{l}{\textit{continued}}\\
\toprule
Method & Dataset & Papers & Min & Max & Spread & Rerun & Copied & Unverif. \\
\midrule
\endhead
\midrule
\multicolumn{9}{r}{\textit{continued on the next page}}\\
\endfoot
\bottomrule
\endlastfoot
AIMIA & IDRiD & 3 & .437 & .437 & .000 & 0 & 0 & 0 \\
Attention U-Net & DDR & 2 & .198 & .389 & .191 & 2 & 0 & 0 \\
Attention U-Net & IDRiD & 2 & .418 & .549 & .132 & 2 & 0 & 0 \\
CASENet & DDR & 3 & .193 & .193 & .000 & 1 & 0 & 0 \\
CASENet & IDRiD & 4 & .482 & .482 & .000 & 1 & 0 & 0 \\
CE-Net & DDR & 2 & .309 & .504 & .195 & 2 & 0 & 0 \\
CE-Net & IDRiD & 2 & .470 & .635 & .165 & 2 & 0 & 0 \\
CIL & IDRiD & 3 & .535 & .535 & .000 & 0 & 0 & 0 \\
DNL & DDR & 2 & .401 & .401 & .000 & 1 & 0 & 0 \\
DNL & IDRiD & 2 & .591 & .591 & .000 & 1 & 0 & 0 \\
Deep-multiscale (Guo2024) & DDR & 3 & .349 & .353 & .004 & 0 & 0 & 0 \\
Deep-multiscale (Guo2024) & IDRiD & 2 & .710 & .710 & .000 & 0 & 0 & 0 \\
DeepLabv3+ & DDR & 10 & .290 & .452 & .162 & 6 & 0 & 0 \\
DeepLabv3+ & IDRiD & 11 & .483 & .666 & .183 & 6 & 0 & 0 \\
DenseUNet & DDR & 2 & .336 & .483 & .147 & 1 & 0 & 0 \\
DenseUNet & IDRiD & 2 & .564 & .651 & .087 & 1 & 0 & 0 \\
EAD-Net & IDRiD & 2 & .549 & .549 & .000 & 0 & 0 & 0 \\
FCRN & DDR & 2 & .096 & .096 & .000 & 1 & 0 & 0 \\
FCRN & IDRiD & 3 & .455 & .455 & .000 & 1 & 0 & 0 \\
FCT & DDR & 2 & .250 & .250 & .000 & 0 & 0 & 0 \\
FCT & IDRiD & 2 & .431 & .431 & .000 & 0 & 0 & 0 \\
HED & DDR & 6 & .191 & .430 & .239 & 2 & 0 & 0 \\
HED & IDRiD & 7 & .366 & .639 & .274 & 2 & 0 & 0 \\
HRNet & DDR & 4 & .452 & .483 & .031 & 3 & 0 & 0 \\
HRNet & IDRiD & 4 & .649 & .689 & .040 & 3 & 0 & 0 \\
HRNet+OCR & DDR & 2 & .452 & .474 & .022 & 0 & 0 & 0 \\
HRNet+OCR & IDRiD & 2 & .649 & .660 & .011 & 0 & 0 & 0 \\
L-Seg & DDR & 13 & .321 & .321 & .000 & 0 & 0 & 0 \\
L-Seg & IDRiD & 14 & .399 & .661 & .262 & 0 & 0 & 0 \\
M2MRF & DDR & 9 & .489 & .499 & .010 & 0 & 0 & 0 \\
M2MRF & IDRiD & 9 & .583 & .672 & .089 & 0 & 0 & 0 \\
MedLabs & IDRiD & 3 & .440 & .440 & .000 & 0 & 0 & 0 \\
MedSAM & DDR & 2 & .193 & .193 & .000 & 0 & 0 & 0 \\
MedSAM & IDRiD & 2 & .266 & .266 & .000 & 0 & 0 & 0 \\
PMCNet & DDR & 4 & .364 & .364 & .000 & 0 & 0 & 0 \\
PMCNet & IDRiD & 3 & .681 & .681 & .000 & 0 & 0 & 0 \\
PSPNet & DDR & 4 & .177 & .392 & .215 & 3 & 0 & 0 \\
PSPNet & IDRiD & 4 & .337 & .587 & .250 & 3 & 0 & 0 \\
RILBP-YNet & IDRiD & 3 & .659 & .659 & .000 & 0 & 0 & 0 \\
RTNet & DDR & 4 & .336 & .336 & .000 & 0 & 0 & 0 \\
RTNet & IDRiD & 3 & .421 & .708 & .287 & 0 & 0 & 0 \\
SAA & DDR & 4 & .366 & .410 & .044 & 0 & 0 & 0 \\
SAA & IDRiD & 5 & .648 & .674 & .026 & 0 & 0 & 0 \\
SAM & DDR & 2 & .003 & .003 & .000 & 0 & 0 & 0 \\
SAM & IDRiD & 2 & .006 & .006 & .000 & 0 & 0 & 0 \\
SAM-Adapter & DDR & 2 & .299 & .299 & .000 & 0 & 0 & 0 \\
SAM-Adapter & IDRiD & 2 & .346 & .346 & .000 & 0 & 0 & 0 \\
SAMed & DDR & 2 & .300 & .300 & .000 & 0 & 0 & 0 \\
SAMed & IDRiD & 2 & .354 & .354 & .000 & 0 & 0 & 0 \\
SDNU & IDRiD & 3 & .477 & .477 & .000 & 0 & 0 & 0 \\
SOONER & IDRiD & 3 & .554 & .554 & .000 & 0 & 0 & 0 \\
Swin (UPerNet) & DDR & 5 & .468 & .486 & .018 & 2 & 0 & 0 \\
Swin (UPerNet) & IDRiD & 4 & .649 & .690 & .042 & 2 & 0 & 0 \\
Twins-SVT-B & DDR & 2 & .461 & .461 & .000 & 1 & 0 & 0 \\
Twins-SVT-B & IDRiD & 2 & .638 & .638 & .000 & 1 & 0 & 0 \\
U-Net & DDR & 7 & .190 & .429 & .239 & 2 & 0 & 0 \\
U-Net & IDRiD & 6 & .347 & .652 & .305 & 2 & 0 & 0 \\
U-Net++ & DDR & 4 & .192 & .409 & .217 & 3 & 0 & 0 \\
U-Net++ & IDRiD & 3 & .391 & .572 & .181 & 3 & 0 & 0 \\
VRT & IDRiD & 8 & .647 & .647 & .000 & 0 & 0 & 0 \\
iFLYTEK-MIG & IDRiD & 7 & .648 & .648 & .000 & 0 & 0 & 0 \\
\end{longtable}
}
\vspace{-6pt}\noindent{\scriptsize Values as printed, normalised to [0, 1]. The full extraction, one row per paper, baseline, dataset, class and metric, is literature/meta/reported\_results.tsv.}

\clearpage

{\scriptsize\setlength{\tabcolsep}{3.2pt}\setlength{\LTcapwidth}{\textwidth}\renewcommand{\arraystretch}{1.08}
\begin{longtable}{l l r c c c c c c c}
\caption{\textbf{Training amount, batch and selection of every main-table run (seed 0).} Steps are the shared pixel budget expressed in the model's patch size (Methods); the loader batch is the largest that fits a 40\,GB card and is accumulated to the recipe's effective batch; the selected checkpoint is the validation \mAUPR{} maximum over 20 evaluations. nnU-Net follows its own 1,000-epoch schedule and planning.}\label{stab:runs}\\
\toprule
Model & Source & Steps & Loader batch & Accum. & Eff. batch & Patch & LR & Best val.\ \mAUPR{} & GPU h \\
\midrule
\endfirsthead
\multicolumn{10}{l}{\textit{continued}}\\
\toprule
Model & Source & Steps & Loader batch & Accum. & Eff. batch & Patch & LR & Best val.\ \mAUPR{} & GPU h \\
\midrule
\endhead
\midrule
\multicolumn{10}{r}{\textit{continued on the next page}}\\
\endfoot
\bottomrule
\endlastfoot
U-Net & IDRiD & 21,000 & 4 & 1 & 4 & 960$\times$960 & 0.0001 & 0.582 & 1.5 \\
U-Net & DDR & 83,600 & 4 & 1 & 4 & 960$\times$960 & 0.0001 & 0.556 & 5.4 \\
DeepLabv3 & IDRiD & 6,000 & 1 & 16 & 16 & 960$\times$960 & 0.007 & 0.428 & 3.6 \\
DeepLabv3 & DDR & 21,200 & 1 & 16 & 16 & 960$\times$960 & 0.007 & 0.436 & 12.8 \\
HRNet-W48 & IDRiD & 7,000 & 2 & 6 & 12 & 960$\times$960 & 0.01 & 0.618 & 2.6 \\
HRNet-W48 & DDR & 28,000 & 2 & 6 & 12 & 960$\times$960 & 0.01 & 0.617 & 10.1 \\
Swin-Unet & IDRiD & 63,000 & 24 & 1 & 24 & 224$\times$224 & 0.01 & 0.644 & 2.6 \\
Swin-Unet & DDR & 255,600 & 24 & 1 & 24 & 224$\times$224 & 0.01 & 0.620 & 10.4 \\
FCT & IDRiD & 129,000 & 4 & 1 & 4 & 384$\times$384 & 0.001 & 0.257 & 19.7 \\
FCT & DDR & 521,600 & 4 & 1 & 4 & 384$\times$384 & 0.001 & 0.175 & 72.0 \\
nnU-Net & IDRiD & \na & \na & \na & \na & \na & \na & \na & \na \\
nnU-Net & DDR & \na & \na & \na & \na & \na & \na & \na & \na \\
nnU-Net ResEnc-M & IDRiD & \na & \na & \na & \na & \na & \na & \na & \na \\
nnU-Net ResEnc-M & DDR & \na & \na & \na & \na & \na & \na & \na & \na \\
M2MRF & IDRiD & 21,000 & 4 & 1 & 4 & 960$\times$960 & 0.01 & 0.541 & 3.0 \\
M2MRF & DDR & 83,600 & 4 & 1 & 4 & 960$\times$960 & 0.01 & 0.558 & 11.6 \\
HACDR-Net & IDRiD & 83,000 & 1 & 1 & 1 & 960$\times$960 & 6e-05 & 0.518 & 5.9 \\
HACDR-Net & DDR & 334,000 & 1 & 1 & 1 & 960$\times$960 & 6e-05 & 0.383 & 27.8 \\
HRDecoder & IDRiD & 21,000 & 4 & 1 & 4 & 960$\times$960 & 0.01 & 0.562 & 1.5 \\
HRDecoder & DDR & 83,600 & 4 & 1 & 4 & 960$\times$960 & 0.01 & 0.539 & 5.9 \\
PMCNet & IDRiD & 83,000 & 1 & 1 & 1 & 960$\times$960 & 0.0001 & 0.544 & 1.8 \\
PMCNet & DDR & 334,000 & 1 & 1 & 1 & 960$\times$960 & 0.0001 & 0.512 & 7.1 \\
WSRFNet & IDRiD & 21,000 & 4 & 1 & 4 & 960$\times$960 & 0.01 & 0.511 & 2.7 \\
WSRFNet & DDR & 83,600 & 4 & 1 & 4 & 960$\times$960 & 0.01 & 0.524 & 10.5 \\
CLC-Net & IDRiD & 73,000 & 16 & 1 & 16 & 512$\times$512 & 0.0003 & 0.580 & 9.9 \\
CLC-Net & DDR & 293,600 & 16 & 1 & 16 & 512$\times$512 & 0.0003 & 0.539 & 39.2 \\
TC-Net & IDRiD & 73,000 & 4 & 1 & 4 & 512$\times$512 & 0.0005 & 0.570 & 6.0 \\
TC-Net & DDR & 293,600 & 4 & 1 & 4 & 512$\times$512 & 0.0005 & 0.472 & 24.7 \\
SAMed & IDRiD & 13,000 & 24 & 1 & 24 & 512$\times$512 & 0.005 & 0.637 & 2.1 \\
SAMed & DDR & 49,200 & 24 & 1 & 24 & 512$\times$512 & 0.005 & 0.632 & 7.4 \\
SAM-Adapter & IDRiD & 37,000 & 2 & 1 & 2 & 1024$\times$1024 & 0.0002 & 0.617 & 6.5 \\
SAM-Adapter & DDR & 146,800 & 2 & 1 & 2 & 1024$\times$1024 & 0.0002 & 0.589 & 26.3 \\
SAM2-UNet & IDRiD & 7,000 & 4 & 3 & 12 & 960$\times$960 & 0.001 & 0.629 & 4.1 \\
SAM2-UNet & DDR & 28,000 & 4 & 3 & 12 & 960$\times$960 & 0.001 & 0.540 & 16.2 \\
MedSAM (full FT) & IDRiD & 37,000 & 2 & 1 & 2 & 1024$\times$1024 & 0.0001 & 0.631 & 2.8 \\
MedSAM (full FT) & DDR & 146,800 & 2 & 1 & 2 & 1024$\times$1024 & 0.0001 & 0.509 & 11.1 \\
RETFound + UPerNet & IDRiD & 63,000 & 24 & 1 & 24 & 224$\times$224 & 0.00046875 & 0.645 & 6.1 \\
RETFound + UPerNet & DDR & 255,600 & 24 & 1 & 24 & 224$\times$224 & 0.00046875 & 0.643 & 24.2 \\
\end{longtable}
}
\vspace{-6pt}\noindent{\scriptsize Every run's full recipe.json (deviations from the source recipe included) is in the repository's run records.}

\clearpage
\begin{landscape}
{\scriptsize\setlength{\tabcolsep}{3.2pt}\setlength{\LTcapwidth}{\textwidth}\renewcommand{\arraystretch}{1.08}
\begin{longtable}{l l c c c c c c c c c}
\caption{\textbf{Per-class AUPR of every main-table model on every test set (seed 0), with 95\% bootstrap intervals.} Exact pooled AUPR inside the field of view; the interval is the percentile interval of 2,000 bootstrap replicates over test images. Both training sources; the internal test is the row whose test set is the training source.}\label{stab:perclass}\\
\toprule
Model & Test set & MA & 95\% CI & HE & 95\% CI & EX & 95\% CI & SE & 95\% CI & \mAUPR{} \\
\midrule
\endfirsthead
\multicolumn{11}{l}{\textit{continued}}\\
\toprule
Model & Test set & MA & 95\% CI & HE & 95\% CI & EX & 95\% CI & SE & 95\% CI & \mAUPR{} \\
\midrule
\endhead
\midrule
\multicolumn{11}{r}{\textit{continued on the next page}}\\
\endfoot
\bottomrule
\endlastfoot
\multicolumn{11}{l}{\textbf{Trained on IDRiD}}\\
U-Net & IDRiD & .406 & \scriptsize .33--.44 & .622 & \scriptsize .57--.67 & .896 & \scriptsize .83--.92 & .707 & \scriptsize .45--.74 & .658 \\
U-Net & DDR & .110 & \scriptsize .08--.14 & .377 & \scriptsize .31--.41 & .477 & \scriptsize .39--.53 & .199 & \scriptsize .05--.23 & .291 \\
U-Net & TJDR std & .255 & \scriptsize .13--.30 & .626 & \scriptsize .51--.71 & .560 & \scriptsize .43--.65 & .235 & \scriptsize .05--.31 & .419 \\
U-Net & TJDR UWF & .079 & \scriptsize .04--.13 & .411 & \scriptsize .24--.53 & .164 & \scriptsize .08--.27 & .180 & \scriptsize .07--.31 & .208 \\
U-Net & MAPLES & .264 & \scriptsize .16--.29 & .328 & \scriptsize .16--.48 & .617 & \scriptsize .32--.73 & .035 & \scriptsize .00--.12 & .311 \\
U-Net & FGADR & .010 & \scriptsize .01--.01 & .200 & \scriptsize .16--.24 & .255 & \scriptsize .21--.29 & .107 & \scriptsize .07--.13 & .143 \\
U-Net & Retinal-L. & .044 & \scriptsize .03--.04 & .186 & \scriptsize .13--.18 & .208 & \scriptsize .13--.21 & .104 & \scriptsize .06--.12 & .136 \\
U-Net & DiaRetDB1 & .036 & \scriptsize .02--.06 & .269 & \scriptsize .17--.33 & .285 & \scriptsize .18--.28 & .112 & \scriptsize .03--.19 & .176 \\
DeepLabv3 & IDRiD & .090 & \scriptsize .07--.12 & .643 & \scriptsize .54--.69 & .734 & \scriptsize .62--.79 & .662 & \scriptsize .49--.76 & .532 \\
DeepLabv3 & DDR & .042 & \scriptsize .03--.06 & .327 & \scriptsize .26--.40 & .511 & \scriptsize .43--.59 & .149 & \scriptsize .07--.24 & .257 \\
DeepLabv3 & TJDR std & .133 & \scriptsize .07--.21 & .605 & \scriptsize .49--.69 & .585 & \scriptsize .32--.75 & .307 & \scriptsize .12--.52 & .407 \\
DeepLabv3 & TJDR UWF & .058 & \scriptsize .02--.14 & .517 & \scriptsize .31--.63 & .320 & \scriptsize .18--.47 & .213 & \scriptsize .10--.39 & .277 \\
DeepLabv3 & MAPLES & .160 & \scriptsize .12--.21 & .237 & \scriptsize .08--.43 & .661 & \scriptsize .41--.77 & .087 & \scriptsize .01--.21 & .286 \\
DeepLabv3 & FGADR & .016 & \scriptsize .01--.02 & .334 & \scriptsize .30--.37 & .335 & \scriptsize .27--.40 & .183 & \scriptsize .14--.23 & .217 \\
DeepLabv3 & Retinal-L. & .093 & \scriptsize .08--.10 & .349 & \scriptsize .31--.38 & .406 & \scriptsize .31--.50 & .195 & \scriptsize .15--.24 & .261 \\
DeepLabv3 & DiaRetDB1 & .010 & \scriptsize .01--.02 & .510 & \scriptsize .41--.59 & .559 & \scriptsize .46--.61 & .209 & \scriptsize .03--.41 & .322 \\
HRNet-W48 & IDRiD & .446 & \scriptsize .41--.48 & .684 & \scriptsize .63--.72 & .884 & \scriptsize .82--.91 & .783 & \scriptsize .59--.87 & .699 \\
HRNet-W48 & DDR & .146 & \scriptsize .11--.19 & .438 & \scriptsize .38--.51 & .507 & \scriptsize .42--.58 & .254 & \scriptsize .12--.38 & .336 \\
HRNet-W48 & TJDR std & .323 & \scriptsize .20--.42 & .706 & \scriptsize .60--.79 & .561 & \scriptsize .40--.69 & .283 & \scriptsize .11--.52 & .468 \\
HRNet-W48 & TJDR UWF & .111 & \scriptsize .05--.20 & .480 & \scriptsize .28--.64 & .132 & \scriptsize .07--.25 & .296 & \scriptsize .13--.49 & .255 \\
HRNet-W48 & MAPLES & .394 & \scriptsize .32--.47 & .492 & \scriptsize .30--.62 & .646 & \scriptsize .38--.78 & .078 & \scriptsize .01--.24 & .403 \\
HRNet-W48 & FGADR & .025 & \scriptsize .02--.03 & .375 & \scriptsize .33--.42 & .362 & \scriptsize .30--.42 & .234 & \scriptsize .19--.29 & .249 \\
HRNet-W48 & Retinal-L. & .070 & \scriptsize .06--.08 & .259 & \scriptsize .22--.30 & .262 & \scriptsize .20--.32 & .198 & \scriptsize .14--.25 & .197 \\
HRNet-W48 & DiaRetDB1 & .039 & \scriptsize .02--.06 & .456 & \scriptsize .36--.54 & .382 & \scriptsize .27--.45 & .212 & \scriptsize .06--.34 & .272 \\
Swin-Unet & IDRiD & .478 & \scriptsize .44--.51 & .709 & \scriptsize .65--.75 & .894 & \scriptsize .81--.93 & .685 & \scriptsize .42--.82 & .692 \\
Swin-Unet & DDR & .132 & \scriptsize .10--.17 & .400 & \scriptsize .33--.48 & .559 & \scriptsize .49--.62 & .216 & \scriptsize .10--.34 & .327 \\
Swin-Unet & TJDR std & .341 & \scriptsize .23--.43 & .495 & \scriptsize .38--.59 & .711 & \scriptsize .60--.80 & .347 & \scriptsize .16--.54 & .474 \\
Swin-Unet & TJDR UWF & .090 & \scriptsize .04--.16 & .471 & \scriptsize .29--.57 & .145 & \scriptsize .07--.27 & .189 & \scriptsize .07--.35 & .224 \\
Swin-Unet & MAPLES & .378 & \scriptsize .31--.44 & .454 & \scriptsize .26--.60 & .662 & \scriptsize .38--.79 & .120 & \scriptsize .03--.24 & .404 \\
Swin-Unet & FGADR & .018 & \scriptsize .01--.02 & .300 & \scriptsize .26--.34 & .303 & \scriptsize .25--.35 & .173 & \scriptsize .12--.23 & .198 \\
Swin-Unet & Retinal-L. & .083 & \scriptsize .07--.09 & .289 & \scriptsize .24--.32 & .359 & \scriptsize .27--.42 & .197 & \scriptsize .15--.24 & .232 \\
Swin-Unet & DiaRetDB1 & .058 & \scriptsize .03--.08 & .459 & \scriptsize .37--.53 & .458 & \scriptsize .34--.51 & .237 & \scriptsize .07--.38 & .303 \\
FCT & IDRiD & .077 & \scriptsize .06--.10 & .163 & \scriptsize .10--.23 & .684 & \scriptsize .57--.74 & .110 & \scriptsize .03--.26 & .259 \\
FCT & DDR & .008 & \scriptsize .00--.01 & .067 & \scriptsize .05--.09 & .224 & \scriptsize .13--.32 & .004 & \scriptsize .00--.01 & .076 \\
FCT & TJDR std & .000 & \scriptsize .00--.00 & .006 & \scriptsize .00--.01 & .004 & \scriptsize .00--.01 & .004 & \scriptsize .00--.01 & .004 \\
FCT & TJDR UWF & .001 & \scriptsize .00--.00 & .020 & \scriptsize .01--.04 & .021 & \scriptsize .01--.05 & .006 & \scriptsize .00--.01 & .012 \\
FCT & MAPLES & .020 & \scriptsize .01--.03 & .018 & \scriptsize .00--.05 & .465 & \scriptsize .15--.64 & .003 & \scriptsize .00--.01 & .126 \\
FCT & FGADR & .005 & \scriptsize .00--.01 & .061 & \scriptsize .05--.07 & .112 & \scriptsize .08--.15 & .009 & \scriptsize .01--.01 & .046 \\
FCT & Retinal-L. & .011 & \scriptsize .01--.01 & .037 & \scriptsize .03--.04 & .049 & \scriptsize .03--.07 & .007 & \scriptsize .01--.01 & .026 \\
FCT & DiaRetDB1 & .004 & \scriptsize .00--.01 & .102 & \scriptsize .04--.17 & .275 & \scriptsize .15--.37 & .002 & \scriptsize .00--.01 & .096 \\
nnU-Net & IDRiD & .499 & \scriptsize .42--.51 & .671 & \scriptsize .61--.72 & .876 & \scriptsize .79--.91 & .641 & \scriptsize .40--.79 & .671 \\
nnU-Net & DDR & .122 & \scriptsize .08--.14 & .384 & \scriptsize .24--.45 & .481 & \scriptsize .37--.54 & .200 & \scriptsize .09--.27 & .297 \\
nnU-Net & TJDR std & .210 & \scriptsize .13--.25 & .615 & \scriptsize .48--.70 & .490 & \scriptsize .32--.60 & .421 & \scriptsize .21--.64 & .434 \\
nnU-Net & TJDR UWF & .088 & \scriptsize .03--.15 & .419 & \scriptsize .20--.59 & .169 & \scriptsize .08--.25 & .144 & \scriptsize .06--.26 & .205 \\
nnU-Net & MAPLES & .314 & \scriptsize .19--.29 & .455 & \scriptsize .26--.56 & .682 & \scriptsize .37--.78 & .151 & \scriptsize .05--.28 & .401 \\
nnU-Net & FGADR & .018 & \scriptsize .01--.02 & .370 & \scriptsize .29--.38 & .360 & \scriptsize .29--.38 & .201 & \scriptsize .14--.24 & .237 \\
nnU-Net & Retinal-L. & .055 & \scriptsize .04--.04 & .283 & \scriptsize .19--.23 & .289 & \scriptsize .18--.27 & .146 & \scriptsize .09--.15 & .193 \\
nnU-Net & DiaRetDB1 & .035 & \scriptsize .02--.05 & .437 & \scriptsize .30--.45 & .428 & \scriptsize .26--.36 & .171 & \scriptsize .05--.29 & .268 \\
nnU-Net ResEnc-M & IDRiD & .449 & \scriptsize .36--.44 & .663 & \scriptsize .58--.68 & .878 & \scriptsize .79--.90 & .533 & \scriptsize .29--.68 & .631 \\
nnU-Net ResEnc-M & DDR & .113 & \scriptsize .08--.12 & .337 & \scriptsize .21--.38 & .476 & \scriptsize .36--.53 & .165 & \scriptsize .07--.24 & .273 \\
nnU-Net ResEnc-M & TJDR std & .217 & \scriptsize .12--.25 & .555 & \scriptsize .41--.61 & .526 & \scriptsize .36--.61 & .385 & \scriptsize .15--.58 & .421 \\
nnU-Net ResEnc-M & TJDR UWF & .071 & \scriptsize .02--.13 & .421 & \scriptsize .18--.60 & .134 & \scriptsize .05--.22 & .143 & \scriptsize .06--.26 & .192 \\
nnU-Net ResEnc-M & MAPLES & .269 & \scriptsize .17--.26 & .472 & \scriptsize .25--.57 & .677 & \scriptsize .34--.74 & .017 & \scriptsize .00--.08 & .359 \\
nnU-Net ResEnc-M & FGADR & .017 & \scriptsize .01--.02 & .322 & \scriptsize .24--.33 & .332 & \scriptsize .26--.36 & .148 & \scriptsize .10--.18 & .205 \\
nnU-Net ResEnc-M & Retinal-L. & .043 & \scriptsize .03--.04 & .192 & \scriptsize .14--.17 & .227 & \scriptsize .15--.23 & .086 & \scriptsize .05--.11 & .137 \\
nnU-Net ResEnc-M & DiaRetDB1 & .035 & \scriptsize .02--.04 & .319 & \scriptsize .20--.33 & .321 & \scriptsize .21--.29 & .232 & \scriptsize .05--.31 & .227 \\
M2MRF & IDRiD & .443 & \scriptsize .38--.46 & .615 & \scriptsize .55--.64 & .868 & \scriptsize .77--.89 & .675 & \scriptsize .47--.76 & .650 \\
M2MRF & DDR & .106 & \scriptsize .08--.13 & .294 & \scriptsize .20--.36 & .410 & \scriptsize .31--.46 & .127 & \scriptsize .06--.20 & .234 \\
M2MRF & TJDR std & .259 & \scriptsize .17--.31 & .482 & \scriptsize .33--.56 & .608 & \scriptsize .49--.69 & .286 & \scriptsize .07--.48 & .409 \\
M2MRF & TJDR UWF & .038 & \scriptsize .01--.07 & .360 & \scriptsize .21--.44 & .035 & \scriptsize .01--.10 & .205 & \scriptsize .08--.29 & .160 \\
M2MRF & MAPLES & .253 & \scriptsize .18--.29 & .400 & \scriptsize .21--.51 & .547 & \scriptsize .32--.65 & .044 & \scriptsize .00--.10 & .311 \\
M2MRF & FGADR & .007 & \scriptsize .01--.01 & .082 & \scriptsize .06--.09 & .172 & \scriptsize .11--.22 & .083 & \scriptsize .05--.11 & .086 \\
M2MRF & Retinal-L. & .044 & \scriptsize .04--.05 & .172 & \scriptsize .12--.17 & .185 & \scriptsize .11--.21 & .089 & \scriptsize .05--.11 & .122 \\
M2MRF & DiaRetDB1 & .038 & \scriptsize .02--.05 & .354 & \scriptsize .23--.38 & .233 & \scriptsize .16--.23 & .165 & \scriptsize .03--.28 & .198 \\
HACDR-Net & IDRiD & .359 & \scriptsize .31--.40 & .502 & \scriptsize .41--.61 & .827 & \scriptsize .74--.88 & .560 & \scriptsize .34--.68 & .562 \\
HACDR-Net & DDR & .091 & \scriptsize .05--.14 & .244 & \scriptsize .17--.31 & .266 & \scriptsize .15--.45 & .060 & \scriptsize .01--.17 & .165 \\
HACDR-Net & TJDR std & .378 & \scriptsize .19--.52 & .475 & \scriptsize .33--.59 & .641 & \scriptsize .55--.72 & .529 & \scriptsize .16--.76 & .506 \\
HACDR-Net & TJDR UWF & .071 & \scriptsize .03--.19 & .267 & \scriptsize .15--.39 & .103 & \scriptsize .03--.23 & .058 & \scriptsize .01--.16 & .125 \\
HACDR-Net & MAPLES & .280 & \scriptsize .21--.35 & .177 & \scriptsize .04--.41 & .579 & \scriptsize .27--.73 & .035 & \scriptsize .00--.15 & .268 \\
HACDR-Net & FGADR & .015 & \scriptsize .01--.02 & .069 & \scriptsize .05--.09 & .044 & \scriptsize .03--.06 & .008 & \scriptsize .00--.02 & .034 \\
HACDR-Net & Retinal-L. & .040 & \scriptsize .03--.04 & .064 & \scriptsize .04--.08 & .183 & \scriptsize .12--.21 & .022 & \scriptsize .01--.04 & .077 \\
HACDR-Net & DiaRetDB1 & .044 & \scriptsize .02--.07 & .290 & \scriptsize .17--.42 & .383 & \scriptsize .22--.43 & .072 & \scriptsize .01--.17 & .197 \\
HRDecoder & IDRiD & .410 & \scriptsize .36--.43 & .629 & \scriptsize .56--.66 & .823 & \scriptsize .70--.85 & .728 & \scriptsize .48--.81 & .647 \\
HRDecoder & DDR & .110 & \scriptsize .08--.13 & .388 & \scriptsize .30--.44 & .530 & \scriptsize .44--.58 & .168 & \scriptsize .06--.22 & .299 \\
HRDecoder & TJDR std & .236 & \scriptsize .14--.29 & .423 & \scriptsize .26--.54 & .593 & \scriptsize .43--.67 & .318 & \scriptsize .13--.46 & .393 \\
HRDecoder & TJDR UWF & .107 & \scriptsize .04--.18 & .362 & \scriptsize .15--.46 & .157 & \scriptsize .08--.23 & .131 & \scriptsize .05--.31 & .189 \\
HRDecoder & MAPLES & .326 & \scriptsize .24--.37 & .464 & \scriptsize .26--.59 & .597 & \scriptsize .33--.70 & .121 & \scriptsize .03--.23 & .377 \\
HRDecoder & FGADR & .020 & \scriptsize .02--.02 & .323 & \scriptsize .27--.35 & .293 & \scriptsize .22--.34 & .179 & \scriptsize .13--.22 & .204 \\
HRDecoder & Retinal-L. & .053 & \scriptsize .04--.05 & .180 & \scriptsize .14--.18 & .182 & \scriptsize .12--.20 & .116 & \scriptsize .07--.14 & .133 \\
HRDecoder & DiaRetDB1 & .052 & \scriptsize .03--.07 & .340 & \scriptsize .24--.37 & .218 & \scriptsize .15--.23 & .217 & \scriptsize .06--.33 & .207 \\
PMCNet & IDRiD & .418 & \scriptsize .36--.47 & .610 & \scriptsize .56--.67 & .883 & \scriptsize .80--.92 & .698 & \scriptsize .41--.82 & .652 \\
PMCNet & DDR & .133 & \scriptsize .10--.17 & .324 & \scriptsize .22--.43 & .464 & \scriptsize .34--.57 & .135 & \scriptsize .04--.28 & .264 \\
PMCNet & TJDR std & .372 & \scriptsize .22--.48 & .600 & \scriptsize .44--.71 & .506 & \scriptsize .39--.62 & .248 & \scriptsize .05--.43 & .432 \\
PMCNet & TJDR UWF & .125 & \scriptsize .04--.21 & .431 & \scriptsize .23--.58 & .210 & \scriptsize .11--.35 & .254 & \scriptsize .13--.40 & .255 \\
PMCNet & MAPLES & .344 & \scriptsize .25--.43 & .334 & \scriptsize .14--.50 & .616 & \scriptsize .35--.75 & .052 & \scriptsize .00--.16 & .337 \\
PMCNet & FGADR & .023 & \scriptsize .02--.03 & .335 & \scriptsize .29--.38 & .355 & \scriptsize .31--.40 & .133 & \scriptsize .09--.18 & .212 \\
PMCNet & Retinal-L. & .076 & \scriptsize .06--.07 & .213 & \scriptsize .16--.24 & .305 & \scriptsize .21--.32 & .129 & \scriptsize .08--.16 & .181 \\
PMCNet & DiaRetDB1 & .033 & \scriptsize .02--.05 & .325 & \scriptsize .22--.41 & .396 & \scriptsize .27--.40 & .153 & \scriptsize .01--.27 & .227 \\
WSRFNet & IDRiD & .388 & \scriptsize .33--.40 & .588 & \scriptsize .51--.61 & .849 & \scriptsize .70--.88 & .502 & \scriptsize .18--.65 & .582 \\
WSRFNet & DDR & .087 & \scriptsize .06--.10 & .287 & \scriptsize .19--.34 & .473 & \scriptsize .38--.49 & .186 & \scriptsize .09--.27 & .258 \\
WSRFNet & TJDR std & .198 & \scriptsize .12--.24 & .426 & \scriptsize .27--.51 & .539 & \scriptsize .40--.59 & .220 & \scriptsize .06--.43 & .346 \\
WSRFNet & TJDR UWF & .061 & \scriptsize .03--.09 & .379 & \scriptsize .22--.48 & .116 & \scriptsize .04--.20 & .230 & \scriptsize .09--.32 & .196 \\
WSRFNet & MAPLES & .245 & \scriptsize .18--.27 & .411 & \scriptsize .22--.51 & .544 & \scriptsize .29--.62 & .032 & \scriptsize .00--.10 & .308 \\
WSRFNet & FGADR & .010 & \scriptsize .01--.01 & .300 & \scriptsize .25--.31 & .314 & \scriptsize .25--.34 & .162 & \scriptsize .11--.19 & .197 \\
WSRFNet & Retinal-L. & .043 & \scriptsize .04--.04 & .141 & \scriptsize .11--.15 & .172 & \scriptsize .11--.20 & .148 & \scriptsize .10--.16 & .126 \\
WSRFNet & DiaRetDB1 & .038 & \scriptsize .02--.05 & .264 & \scriptsize .19--.31 & .221 & \scriptsize .16--.23 & .120 & \scriptsize .02--.20 & .161 \\
CLC-Net & IDRiD & .437 & \scriptsize .37--.46 & .683 & \scriptsize .59--.71 & .889 & \scriptsize .81--.92 & .634 & \scriptsize .40--.73 & .661 \\
CLC-Net & DDR & .115 & \scriptsize .08--.14 & .439 & \scriptsize .36--.48 & .539 & \scriptsize .45--.57 & .154 & \scriptsize .06--.23 & .312 \\
CLC-Net & TJDR std & .295 & \scriptsize .16--.38 & .631 & \scriptsize .50--.71 & .618 & \scriptsize .45--.73 & .292 & \scriptsize .12--.43 & .459 \\
CLC-Net & TJDR UWF & .091 & \scriptsize .03--.16 & .466 & \scriptsize .25--.57 & .170 & \scriptsize .08--.26 & .105 & \scriptsize .04--.18 & .208 \\
CLC-Net & MAPLES & .336 & \scriptsize .23--.37 & .477 & \scriptsize .28--.59 & .689 & \scriptsize .41--.78 & .065 & \scriptsize .01--.16 & .392 \\
CLC-Net & FGADR & .007 & \scriptsize .01--.01 & .318 & \scriptsize .27--.35 & .288 & \scriptsize .24--.31 & .086 & \scriptsize .05--.11 & .175 \\
CLC-Net & Retinal-L. & .060 & \scriptsize .04--.05 & .246 & \scriptsize .19--.26 & .348 & \scriptsize .20--.32 & .105 & \scriptsize .07--.12 & .190 \\
CLC-Net & DiaRetDB1 & .044 & \scriptsize .02--.06 & .400 & \scriptsize .28--.48 & .444 & \scriptsize .25--.36 & .140 & \scriptsize .03--.23 & .257 \\
TC-Net & IDRiD & .355 & \scriptsize .31--.40 & .650 & \scriptsize .59--.70 & .854 & \scriptsize .78--.89 & .754 & \scriptsize .57--.85 & .653 \\
TC-Net & DDR & .119 & \scriptsize .09--.15 & .404 & \scriptsize .32--.49 & .572 & \scriptsize .50--.63 & .153 & \scriptsize .07--.25 & .312 \\
TC-Net & TJDR std & .233 & \scriptsize .16--.31 & .605 & \scriptsize .50--.69 & .691 & \scriptsize .58--.77 & .398 & \scriptsize .15--.61 & .482 \\
TC-Net & TJDR UWF & .124 & \scriptsize .07--.22 & .489 & \scriptsize .27--.61 & .017 & \scriptsize .01--.04 & .223 & \scriptsize .12--.34 & .213 \\
TC-Net & MAPLES & .372 & \scriptsize .29--.46 & .453 & \scriptsize .27--.58 & .682 & \scriptsize .39--.81 & .143 & \scriptsize .03--.30 & .412 \\
TC-Net & FGADR & .022 & \scriptsize .02--.02 & .366 & \scriptsize .32--.41 & .318 & \scriptsize .26--.37 & .144 & \scriptsize .11--.19 & .212 \\
TC-Net & Retinal-L. & .080 & \scriptsize .07--.08 & .270 & \scriptsize .22--.30 & .237 & \scriptsize .16--.29 & .114 & \scriptsize .08--.15 & .175 \\
TC-Net & DiaRetDB1 & .039 & \scriptsize .02--.06 & .433 & \scriptsize .34--.51 & .426 & \scriptsize .29--.48 & .218 & \scriptsize .09--.32 & .279 \\
SAMed & IDRiD & .449 & \scriptsize .40--.50 & .736 & \scriptsize .67--.78 & .892 & \scriptsize .82--.92 & .790 & \scriptsize .61--.88 & .717 \\
SAMed & DDR & .128 & \scriptsize .10--.16 & .443 & \scriptsize .37--.52 & .609 & \scriptsize .54--.67 & .339 & \scriptsize .18--.48 & .380 \\
SAMed & TJDR std & .301 & \scriptsize .20--.38 & .669 & \scriptsize .53--.76 & .706 & \scriptsize .61--.79 & .559 & \scriptsize .25--.74 & .559 \\
SAMed & TJDR UWF & .115 & \scriptsize .04--.19 & .470 & \scriptsize .24--.66 & .226 & \scriptsize .13--.36 & .158 & \scriptsize .05--.33 & .242 \\
SAMed & MAPLES & .351 & \scriptsize .26--.44 & .471 & \scriptsize .30--.60 & .684 & \scriptsize .40--.80 & .136 & \scriptsize .04--.30 & .410 \\
SAMed & FGADR & .025 & \scriptsize .02--.03 & .418 & \scriptsize .38--.46 & .409 & \scriptsize .36--.46 & .282 & \scriptsize .23--.34 & .283 \\
SAMed & Retinal-L. & .083 & \scriptsize .07--.09 & .354 & \scriptsize .31--.38 & .362 & \scriptsize .28--.42 & .219 & \scriptsize .16--.26 & .255 \\
SAMed & DiaRetDB1 & .055 & \scriptsize .03--.08 & .473 & \scriptsize .39--.54 & .455 & \scriptsize .35--.50 & .276 & \scriptsize .08--.45 & .315 \\
SAM-Adapter & IDRiD & .451 & \scriptsize .39--.46 & .740 & \scriptsize .66--.78 & .882 & \scriptsize .79--.92 & .750 & \scriptsize .54--.82 & .706 \\
SAM-Adapter & DDR & .117 & \scriptsize .08--.14 & .403 & \scriptsize .30--.46 & .560 & \scriptsize .47--.59 & .157 & \scriptsize .08--.23 & .309 \\
SAM-Adapter & TJDR std & .256 & \scriptsize .14--.33 & .581 & \scriptsize .43--.69 & .563 & \scriptsize .40--.63 & .351 & \scriptsize .16--.55 & .438 \\
SAM-Adapter & TJDR UWF & .067 & \scriptsize .02--.10 & .447 & \scriptsize .21--.58 & .125 & \scriptsize .05--.21 & .196 & \scriptsize .09--.32 & .209 \\
SAM-Adapter & MAPLES & .292 & \scriptsize .19--.32 & .395 & \scriptsize .21--.50 & .605 & \scriptsize .31--.71 & .144 & \scriptsize .05--.27 & .359 \\
SAM-Adapter & FGADR & .018 & \scriptsize .01--.02 & .355 & \scriptsize .29--.39 & .343 & \scriptsize .27--.39 & .166 & \scriptsize .12--.22 & .221 \\
SAM-Adapter & Retinal-L. & .046 & \scriptsize .04--.04 & .246 & \scriptsize .17--.23 & .317 & \scriptsize .17--.32 & .142 & \scriptsize .10--.16 & .188 \\
SAM-Adapter & DiaRetDB1 & .037 & \scriptsize .02--.06 & .449 & \scriptsize .30--.49 & .365 & \scriptsize .19--.32 & .130 & \scriptsize .03--.26 & .245 \\
SAM2-UNet & IDRiD & .439 & \scriptsize .37--.45 & .670 & \scriptsize .59--.72 & .871 & \scriptsize .78--.90 & .678 & \scriptsize .43--.77 & .664 \\
SAM2-UNet & DDR & .112 & \scriptsize .08--.13 & .448 & \scriptsize .38--.50 & .516 & \scriptsize .42--.56 & .239 & \scriptsize .05--.17 & .329 \\
SAM2-UNet & TJDR std & .289 & \scriptsize .16--.35 & .680 & \scriptsize .58--.75 & .643 & \scriptsize .53--.69 & .510 & \scriptsize .14--.52 & .530 \\
SAM2-UNet & TJDR UWF & .091 & \scriptsize .03--.14 & .482 & \scriptsize .30--.62 & .121 & \scriptsize .04--.23 & .230 & \scriptsize .06--.31 & .231 \\
SAM2-UNet & MAPLES & .319 & \scriptsize .23--.37 & .404 & \scriptsize .21--.53 & .644 & \scriptsize .38--.75 & .097 & \scriptsize .01--.21 & .366 \\
SAM2-UNet & FGADR & .016 & \scriptsize .01--.02 & .378 & \scriptsize .32--.40 & .377 & \scriptsize .30--.40 & .242 & \scriptsize .16--.24 & .253 \\
SAM2-UNet & Retinal-L. & .063 & \scriptsize .05--.06 & .249 & \scriptsize .18--.25 & .274 & \scriptsize .18--.29 & .159 & \scriptsize .08--.15 & .186 \\
SAM2-UNet & DiaRetDB1 & .037 & \scriptsize .02--.05 & .327 & \scriptsize .24--.39 & .319 & \scriptsize .23--.30 & .197 & \scriptsize .04--.30 & .220 \\
MedSAM (full FT) & IDRiD & .373 & \scriptsize .33--.41 & .665 & \scriptsize .57--.73 & .855 & \scriptsize .77--.90 & .674 & \scriptsize .43--.83 & .642 \\
MedSAM (full FT) & DDR & .144 & \scriptsize .11--.18 & .400 & \scriptsize .35--.46 & .555 & \scriptsize .49--.62 & .155 & \scriptsize .08--.25 & .314 \\
MedSAM (full FT) & TJDR std & .224 & \scriptsize .17--.27 & .593 & \scriptsize .42--.71 & .600 & \scriptsize .47--.70 & .395 & \scriptsize .19--.57 & .453 \\
MedSAM (full FT) & TJDR UWF & .119 & \scriptsize .05--.24 & .443 & \scriptsize .24--.57 & .213 & \scriptsize .12--.34 & .403 & \scriptsize .25--.55 & .295 \\
MedSAM (full FT) & MAPLES & .331 & \scriptsize .27--.39 & .364 & \scriptsize .17--.52 & .650 & \scriptsize .38--.77 & .132 & \scriptsize .05--.25 & .369 \\
MedSAM (full FT) & FGADR & .024 & \scriptsize .02--.03 & .367 & \scriptsize .33--.41 & .394 & \scriptsize .35--.44 & .254 & \scriptsize .20--.31 & .260 \\
MedSAM (full FT) & Retinal-L. & .078 & \scriptsize .07--.08 & .260 & \scriptsize .22--.29 & .394 & \scriptsize .33--.45 & .125 & \scriptsize .10--.15 & .214 \\
MedSAM (full FT) & DiaRetDB1 & .043 & \scriptsize .02--.07 & .417 & \scriptsize .30--.52 & .535 & \scriptsize .42--.60 & .131 & \scriptsize .03--.24 & .282 \\
RETFound + UPerNet & IDRiD & .484 & \scriptsize .45--.52 & .718 & \scriptsize .66--.75 & .895 & \scriptsize .82--.93 & .734 & \scriptsize .50--.85 & .708 \\
RETFound + UPerNet & DDR & .152 & \scriptsize .12--.19 & .413 & \scriptsize .32--.50 & .572 & \scriptsize .50--.63 & .278 & \scriptsize .13--.43 & .354 \\
RETFound + UPerNet & TJDR std & .393 & \scriptsize .25--.49 & .582 & \scriptsize .46--.67 & .776 & \scriptsize .68--.86 & .070 & \scriptsize .02--.15 & .455 \\
RETFound + UPerNet & TJDR UWF & .099 & \scriptsize .04--.18 & .494 & \scriptsize .28--.63 & .106 & \scriptsize .05--.27 & .311 & \scriptsize .15--.51 & .253 \\
RETFound + UPerNet & MAPLES & .375 & \scriptsize .29--.45 & .503 & \scriptsize .30--.64 & .680 & \scriptsize .38--.81 & .092 & \scriptsize .03--.23 & .413 \\
RETFound + UPerNet & FGADR & .023 & \scriptsize .02--.03 & .332 & \scriptsize .29--.38 & .357 & \scriptsize .29--.42 & .215 & \scriptsize .16--.28 & .232 \\
RETFound + UPerNet & Retinal-L. & .081 & \scriptsize .06--.08 & .288 & \scriptsize .23--.32 & .299 & \scriptsize .22--.36 & .156 & \scriptsize .11--.20 & .206 \\
RETFound + UPerNet & DiaRetDB1 & .039 & \scriptsize .02--.06 & .459 & \scriptsize .34--.56 & .474 & \scriptsize .29--.55 & .257 & \scriptsize .07--.40 & .307 \\
\multicolumn{11}{l}{\textbf{Trained on DDR}}\\
U-Net & IDRiD & .364 & \scriptsize .29--.36 & .541 & \scriptsize .43--.57 & .800 & \scriptsize .66--.78 & .483 & \scriptsize .26--.56 & .547 \\
U-Net & DDR & .141 & \scriptsize .09--.16 & .410 & \scriptsize .28--.48 & .553 & \scriptsize .42--.59 & .410 & \scriptsize .23--.49 & .378 \\
U-Net & TJDR std & .296 & \scriptsize .16--.34 & .676 & \scriptsize .55--.75 & .614 & \scriptsize .42--.65 & .392 & \scriptsize .13--.49 & .494 \\
U-Net & TJDR UWF & .077 & \scriptsize .03--.21 & .262 & \scriptsize .11--.39 & .254 & \scriptsize .15--.32 & .067 & \scriptsize .01--.13 & .165 \\
U-Net & MAPLES & .222 & \scriptsize .13--.25 & .462 & \scriptsize .24--.58 & .509 & \scriptsize .24--.50 & .125 & \scriptsize .01--.21 & .329 \\
U-Net & FGADR & .013 & \scriptsize .01--.01 & .214 & \scriptsize .18--.24 & .292 & \scriptsize .23--.33 & .162 & \scriptsize .11--.19 & .170 \\
U-Net & Retinal-L. & .039 & \scriptsize .03--.04 & .165 & \scriptsize .10--.16 & .167 & \scriptsize .10--.16 & .197 & \scriptsize .13--.21 & .142 \\
U-Net & DiaRetDB1 & .039 & \scriptsize .02--.06 & .314 & \scriptsize .19--.34 & .268 & \scriptsize .12--.20 & .253 & \scriptsize .11--.31 & .219 \\
DeepLabv3 & IDRiD & .052 & \scriptsize .04--.07 & .640 & \scriptsize .49--.72 & .708 & \scriptsize .58--.77 & .630 & \scriptsize .40--.75 & .507 \\
DeepLabv3 & DDR & .036 & \scriptsize .02--.05 & .447 & \scriptsize .30--.61 & .614 & \scriptsize .54--.68 & .387 & \scriptsize .25--.50 & .371 \\
DeepLabv3 & TJDR std & .172 & \scriptsize .08--.28 & .649 & \scriptsize .51--.75 & .724 & \scriptsize .64--.76 & .411 & \scriptsize .15--.61 & .489 \\
DeepLabv3 & TJDR UWF & .071 & \scriptsize .03--.10 & .066 & \scriptsize .02--.20 & .296 & \scriptsize .17--.44 & .058 & \scriptsize .02--.18 & .123 \\
DeepLabv3 & MAPLES & .143 & \scriptsize .10--.19 & .452 & \scriptsize .24--.58 & .641 & \scriptsize .41--.73 & .108 & \scriptsize .01--.24 & .336 \\
DeepLabv3 & FGADR & .011 & \scriptsize .01--.01 & .339 & \scriptsize .30--.38 & .356 & \scriptsize .30--.41 & .224 & \scriptsize .18--.27 & .232 \\
DeepLabv3 & Retinal-L. & .074 & \scriptsize .06--.08 & .316 & \scriptsize .25--.37 & .440 & \scriptsize .32--.52 & .253 & \scriptsize .19--.31 & .271 \\
DeepLabv3 & DiaRetDB1 & .008 & \scriptsize .00--.01 & .495 & \scriptsize .40--.57 & .545 & \scriptsize .44--.59 & .241 & \scriptsize .03--.40 & .322 \\
HRNet-W48 & IDRiD & .372 & \scriptsize .34--.40 & .694 & \scriptsize .58--.76 & .834 & \scriptsize .75--.88 & .679 & \scriptsize .42--.79 & .645 \\
HRNet-W48 & DDR & .155 & \scriptsize .12--.20 & .477 & \scriptsize .39--.59 & .590 & \scriptsize .52--.67 & .550 & \scriptsize .42--.65 & .443 \\
HRNet-W48 & TJDR std & .412 & \scriptsize .24--.53 & .788 & \scriptsize .71--.85 & .713 & \scriptsize .63--.76 & .581 & \scriptsize .25--.76 & .624 \\
HRNet-W48 & TJDR UWF & .116 & \scriptsize .06--.24 & .120 & \scriptsize .04--.29 & .248 & \scriptsize .15--.36 & .107 & \scriptsize .02--.27 & .148 \\
HRNet-W48 & MAPLES & .383 & \scriptsize .31--.45 & .500 & \scriptsize .28--.63 & .701 & \scriptsize .48--.79 & .220 & \scriptsize .05--.37 & .451 \\
HRNet-W48 & FGADR & .026 & \scriptsize .02--.03 & .399 & \scriptsize .35--.45 & .394 & \scriptsize .34--.45 & .255 & \scriptsize .20--.32 & .269 \\
HRNet-W48 & Retinal-L. & .083 & \scriptsize .07--.09 & .251 & \scriptsize .20--.31 & .339 & \scriptsize .27--.40 & .332 & \scriptsize .28--.38 & .251 \\
HRNet-W48 & DiaRetDB1 & .034 & \scriptsize .02--.05 & .437 & \scriptsize .34--.53 & .432 & \scriptsize .32--.52 & .362 & \scriptsize .21--.46 & .316 \\
Swin-Unet & IDRiD & .420 & \scriptsize .38--.46 & .625 & \scriptsize .52--.69 & .853 & \scriptsize .76--.89 & .615 & \scriptsize .41--.73 & .628 \\
Swin-Unet & DDR & .174 & \scriptsize .13--.22 & .479 & \scriptsize .39--.58 & .650 & \scriptsize .59--.71 & .424 & \scriptsize .29--.53 & .432 \\
Swin-Unet & TJDR std & .370 & \scriptsize .24--.45 & .676 & \scriptsize .56--.76 & .696 & \scriptsize .53--.79 & .438 & \scriptsize .13--.63 & .545 \\
Swin-Unet & TJDR UWF & .107 & \scriptsize .05--.22 & .336 & \scriptsize .19--.47 & .271 & \scriptsize .18--.37 & .044 & \scriptsize .01--.14 & .189 \\
Swin-Unet & MAPLES & .395 & \scriptsize .32--.46 & .509 & \scriptsize .31--.64 & .726 & \scriptsize .50--.81 & .200 & \scriptsize .06--.35 & .457 \\
Swin-Unet & FGADR & .021 & \scriptsize .02--.02 & .349 & \scriptsize .31--.39 & .382 & \scriptsize .33--.44 & .216 & \scriptsize .17--.27 & .242 \\
Swin-Unet & Retinal-L. & .113 & \scriptsize .09--.12 & .254 & \scriptsize .20--.30 & .418 & \scriptsize .33--.48 & .242 & \scriptsize .19--.29 & .257 \\
Swin-Unet & DiaRetDB1 & .048 & \scriptsize .03--.07 & .381 & \scriptsize .28--.47 & .496 & \scriptsize .40--.55 & .255 & \scriptsize .14--.36 & .295 \\
FCT & IDRiD & .005 & \scriptsize .00--.01 & .153 & \scriptsize .10--.21 & .679 & \scriptsize .56--.74 & .046 & \scriptsize .01--.13 & .221 \\
FCT & DDR & .001 & \scriptsize .00--.00 & .163 & \scriptsize .11--.22 & .524 & \scriptsize .46--.58 & .014 & \scriptsize .01--.02 & .176 \\
FCT & TJDR std & .001 & \scriptsize .00--.00 & .099 & \scriptsize .05--.15 & .596 & \scriptsize .43--.68 & .040 & \scriptsize .01--.11 & .184 \\
FCT & TJDR UWF & .001 & \scriptsize .00--.00 & .055 & \scriptsize .02--.10 & .135 & \scriptsize .06--.24 & .009 & \scriptsize .00--.02 & .050 \\
FCT & MAPLES & .004 & \scriptsize .00--.01 & .030 & \scriptsize .01--.08 & .510 & \scriptsize .21--.66 & .007 & \scriptsize .00--.02 & .138 \\
FCT & FGADR & .004 & \scriptsize .00--.00 & .065 & \scriptsize .05--.08 & .182 & \scriptsize .14--.23 & .015 & \scriptsize .01--.02 & .067 \\
FCT & Retinal-L. & .010 & \scriptsize .01--.01 & .057 & \scriptsize .04--.07 & .210 & \scriptsize .14--.28 & .018 & \scriptsize .01--.02 & .074 \\
FCT & DiaRetDB1 & .000 & \scriptsize .00--.00 & .087 & \scriptsize .02--.17 & .270 & \scriptsize .16--.34 & .009 & \scriptsize .00--.02 & .092 \\
nnU-Net & IDRiD & .334 & \scriptsize .24--.34 & .637 & \scriptsize .52--.71 & .810 & \scriptsize .71--.86 & .623 & \scriptsize .38--.71 & .601 \\
nnU-Net & DDR & .154 & \scriptsize .11--.17 & .501 & \scriptsize .39--.57 & .614 & \scriptsize .51--.66 & .458 & \scriptsize .30--.56 & .432 \\
nnU-Net & TJDR std & .309 & \scriptsize .17--.36 & .705 & \scriptsize .61--.75 & .635 & \scriptsize .45--.71 & .425 & \scriptsize .15--.61 & .518 \\
nnU-Net & TJDR UWF & .111 & \scriptsize .03--.21 & .374 & \scriptsize .18--.54 & .157 & \scriptsize .07--.26 & .142 & \scriptsize .06--.27 & .196 \\
nnU-Net & MAPLES & .372 & \scriptsize .18--.29 & .493 & \scriptsize .27--.58 & .717 & \scriptsize .38--.75 & .196 & \scriptsize .06--.29 & .445 \\
nnU-Net & FGADR & .020 & \scriptsize .01--.02 & .396 & \scriptsize .33--.41 & .389 & \scriptsize .31--.42 & .232 & \scriptsize .18--.27 & .259 \\
nnU-Net & Retinal-L. & .067 & \scriptsize .04--.05 & .232 & \scriptsize .14--.19 & .319 & \scriptsize .17--.28 & .224 & \scriptsize .17--.23 & .211 \\
nnU-Net & DiaRetDB1 & .065 & \scriptsize .04--.09 & .462 & \scriptsize .31--.48 & .476 & \scriptsize .26--.39 & .295 & \scriptsize .10--.35 & .325 \\
nnU-Net ResEnc-M & IDRiD & .349 & \scriptsize .26--.32 & .664 & \scriptsize .54--.72 & .832 & \scriptsize .73--.87 & .603 & \scriptsize .36--.67 & .612 \\
nnU-Net ResEnc-M & DDR & .147 & \scriptsize .10--.16 & .481 & \scriptsize .37--.55 & .653 & \scriptsize .56--.70 & .432 & \scriptsize .28--.52 & .428 \\
nnU-Net ResEnc-M & TJDR std & .320 & \scriptsize .19--.34 & .751 & \scriptsize .63--.81 & .673 & \scriptsize .48--.75 & .538 & \scriptsize .19--.69 & .571 \\
nnU-Net ResEnc-M & TJDR UWF & .115 & \scriptsize .03--.21 & .317 & \scriptsize .17--.44 & .132 & \scriptsize .06--.24 & .112 & \scriptsize .02--.23 & .169 \\
nnU-Net ResEnc-M & MAPLES & .341 & \scriptsize .16--.29 & .469 & \scriptsize .26--.58 & .762 & \scriptsize .42--.78 & .293 & \scriptsize .12--.37 & .466 \\
nnU-Net ResEnc-M & FGADR & .018 & \scriptsize .01--.02 & .409 & \scriptsize .33--.41 & .401 & \scriptsize .31--.42 & .242 & \scriptsize .17--.27 & .268 \\
nnU-Net ResEnc-M & Retinal-L. & .062 & \scriptsize .03--.04 & .269 & \scriptsize .13--.21 & .355 & \scriptsize .20--.29 & .257 & \scriptsize .18--.24 & .236 \\
nnU-Net ResEnc-M & DiaRetDB1 & .062 & \scriptsize .03--.08 & .450 & \scriptsize .28--.44 & .475 & \scriptsize .27--.37 & .261 & \scriptsize .06--.32 & .312 \\
M2MRF & IDRiD & .392 & \scriptsize .34--.40 & .555 & \scriptsize .45--.56 & .776 & \scriptsize .68--.78 & .563 & \scriptsize .35--.66 & .571 \\
M2MRF & DDR & .148 & \scriptsize .10--.18 & .402 & \scriptsize .27--.48 & .547 & \scriptsize .46--.59 & .422 & \scriptsize .25--.49 & .380 \\
M2MRF & TJDR std & .299 & \scriptsize .19--.35 & .696 & \scriptsize .58--.75 & .618 & \scriptsize .50--.68 & .412 & \scriptsize .14--.55 & .506 \\
M2MRF & TJDR UWF & .093 & \scriptsize .04--.17 & .403 & \scriptsize .22--.51 & .209 & \scriptsize .11--.30 & .037 & \scriptsize .01--.10 & .186 \\
M2MRF & MAPLES & .266 & \scriptsize .17--.30 & .437 & \scriptsize .24--.55 & .454 & \scriptsize .30--.50 & .138 & \scriptsize .02--.25 & .324 \\
M2MRF & FGADR & .019 & \scriptsize .01--.02 & .310 & \scriptsize .24--.33 & .340 & \scriptsize .27--.38 & .181 & \scriptsize .13--.22 & .212 \\
M2MRF & Retinal-L. & .050 & \scriptsize .04--.05 & .183 & \scriptsize .13--.17 & .174 & \scriptsize .12--.18 & .195 & \scriptsize .13--.22 & .150 \\
M2MRF & DiaRetDB1 & .049 & \scriptsize .02--.07 & .218 & \scriptsize .15--.26 & .177 & \scriptsize .13--.18 & .247 & \scriptsize .14--.30 & .173 \\
HACDR-Net & IDRiD & .225 & \scriptsize .18--.26 & .365 & \scriptsize .28--.47 & .789 & \scriptsize .67--.83 & .199 & \scriptsize .04--.44 & .395 \\
HACDR-Net & DDR & .114 & \scriptsize .08--.16 & .309 & \scriptsize .22--.40 & .485 & \scriptsize .37--.58 & .168 & \scriptsize .08--.28 & .269 \\
HACDR-Net & TJDR std & .356 & \scriptsize .21--.46 & .632 & \scriptsize .54--.71 & .553 & \scriptsize .45--.64 & .196 & \scriptsize .03--.39 & .434 \\
HACDR-Net & TJDR UWF & .095 & \scriptsize .03--.20 & .385 & \scriptsize .20--.50 & .204 & \scriptsize .11--.32 & .130 & \scriptsize .04--.24 & .203 \\
HACDR-Net & MAPLES & .224 & \scriptsize .15--.29 & .378 & \scriptsize .17--.54 & .512 & \scriptsize .22--.63 & .093 & \scriptsize .01--.22 & .302 \\
HACDR-Net & FGADR & .021 & \scriptsize .02--.02 & .186 & \scriptsize .15--.22 & .134 & \scriptsize .10--.17 & .049 & \scriptsize .03--.07 & .098 \\
HACDR-Net & Retinal-L. & .058 & \scriptsize .04--.06 & .163 & \scriptsize .12--.19 & .166 & \scriptsize .11--.19 & .138 & \scriptsize .10--.17 & .131 \\
HACDR-Net & DiaRetDB1 & .016 & \scriptsize .01--.03 & .228 & \scriptsize .12--.32 & .319 & \scriptsize .15--.34 & .101 & \scriptsize .02--.19 & .166 \\
HRDecoder & IDRiD & .357 & \scriptsize .31--.36 & .547 & \scriptsize .41--.57 & .714 & \scriptsize .60--.73 & .539 & \scriptsize .35--.60 & .539 \\
HRDecoder & DDR & .124 & \scriptsize .08--.16 & .318 & \scriptsize .20--.40 & .496 & \scriptsize .40--.56 & .245 & \scriptsize .11--.33 & .296 \\
HRDecoder & TJDR std & .334 & \scriptsize .19--.41 & .689 & \scriptsize .59--.75 & .614 & \scriptsize .51--.67 & .358 & \scriptsize .12--.51 & .499 \\
HRDecoder & TJDR UWF & .119 & \scriptsize .05--.18 & .368 & \scriptsize .17--.53 & .180 & \scriptsize .09--.27 & .124 & \scriptsize .05--.23 & .198 \\
HRDecoder & MAPLES & .231 & \scriptsize .14--.26 & .407 & \scriptsize .23--.50 & .427 & \scriptsize .22--.48 & .127 & \scriptsize .03--.22 & .298 \\
HRDecoder & FGADR & .014 & \scriptsize .01--.02 & .267 & \scriptsize .21--.29 & .309 & \scriptsize .25--.35 & .183 & \scriptsize .14--.22 & .193 \\
HRDecoder & Retinal-L. & .043 & \scriptsize .03--.04 & .126 & \scriptsize .08--.14 & .171 & \scriptsize .11--.17 & .168 & \scriptsize .13--.19 & .127 \\
HRDecoder & DiaRetDB1 & .045 & \scriptsize .02--.06 & .275 & \scriptsize .17--.32 & .234 & \scriptsize .13--.19 & .203 & \scriptsize .08--.29 & .189 \\
PMCNet & IDRiD & .348 & \scriptsize .31--.38 & .458 & \scriptsize .37--.51 & .758 & \scriptsize .68--.81 & .524 & \scriptsize .20--.71 & .522 \\
PMCNet & DDR & .157 & \scriptsize .12--.20 & .346 & \scriptsize .23--.46 & .482 & \scriptsize .37--.58 & .262 & \scriptsize .14--.38 & .312 \\
PMCNet & TJDR std & .406 & \scriptsize .25--.50 & .639 & \scriptsize .54--.72 & .476 & \scriptsize .33--.60 & .308 & \scriptsize .06--.49 & .457 \\
PMCNet & TJDR UWF & .085 & \scriptsize .03--.19 & .427 & \scriptsize .23--.55 & .217 & \scriptsize .12--.32 & .091 & \scriptsize .03--.19 & .205 \\
PMCNet & MAPLES & .371 & \scriptsize .28--.44 & .444 & \scriptsize .24--.59 & .627 & \scriptsize .34--.74 & .131 & \scriptsize .00--.28 & .393 \\
PMCNet & FGADR & .021 & \scriptsize .02--.02 & .323 & \scriptsize .27--.37 & .339 & \scriptsize .29--.37 & .162 & \scriptsize .11--.21 & .211 \\
PMCNet & Retinal-L. & .105 & \scriptsize .08--.10 & .236 & \scriptsize .19--.25 & .272 & \scriptsize .18--.26 & .128 & \scriptsize .09--.16 & .185 \\
PMCNet & DiaRetDB1 & .036 & \scriptsize .02--.06 & .312 & \scriptsize .21--.40 & .366 & \scriptsize .23--.38 & .266 & \scriptsize .04--.42 & .245 \\
WSRFNet & IDRiD & .319 & \scriptsize .24--.32 & .489 & \scriptsize .36--.51 & .673 & \scriptsize .53--.72 & .384 & \scriptsize .17--.49 & .466 \\
WSRFNet & DDR & .088 & \scriptsize .05--.11 & .292 & \scriptsize .20--.37 & .422 & \scriptsize .31--.48 & .135 & \scriptsize .05--.21 & .234 \\
WSRFNet & TJDR std & .173 & \scriptsize .08--.26 & .588 & \scriptsize .43--.64 & .433 & \scriptsize .23--.51 & .114 & \scriptsize .01--.20 & .327 \\
WSRFNet & TJDR UWF & .038 & \scriptsize .01--.18 & .323 & \scriptsize .16--.39 & .155 & \scriptsize .06--.24 & .024 & \scriptsize .00--.08 & .135 \\
WSRFNet & MAPLES & .264 & \scriptsize .17--.29 & .428 & \scriptsize .21--.53 & .436 & \scriptsize .26--.43 & .104 & \scriptsize .01--.18 & .308 \\
WSRFNet & FGADR & .011 & \scriptsize .01--.01 & .158 & \scriptsize .12--.17 & .086 & \scriptsize .05--.09 & .020 & \scriptsize .01--.03 & .069 \\
WSRFNet & Retinal-L. & .037 & \scriptsize .03--.04 & .085 & \scriptsize .05--.10 & .107 & \scriptsize .07--.12 & .062 & \scriptsize .03--.09 & .073 \\
WSRFNet & DiaRetDB1 & .050 & \scriptsize .02--.07 & .215 & \scriptsize .15--.25 & .178 & \scriptsize .13--.19 & .158 & \scriptsize .08--.20 & .150 \\
CLC-Net & IDRiD & .358 & \scriptsize .29--.37 & .614 & \scriptsize .47--.64 & .830 & \scriptsize .71--.86 & .570 & \scriptsize .25--.67 & .593 \\
CLC-Net & DDR & .152 & \scriptsize .10--.18 & .466 & \scriptsize .34--.53 & .612 & \scriptsize .51--.64 & .428 & \scriptsize .25--.51 & .414 \\
CLC-Net & TJDR std & .352 & \scriptsize .19--.42 & .666 & \scriptsize .55--.73 & .667 & \scriptsize .49--.73 & .296 & \scriptsize .09--.45 & .495 \\
CLC-Net & TJDR UWF & .110 & \scriptsize .04--.21 & .137 & \scriptsize .02--.14 & .159 & \scriptsize .06--.28 & .095 & \scriptsize .02--.29 & .125 \\
CLC-Net & MAPLES & .287 & \scriptsize .17--.31 & .475 & \scriptsize .26--.56 & .611 & \scriptsize .33--.61 & .030 & \scriptsize .00--.11 & .351 \\
CLC-Net & FGADR & .011 & \scriptsize .01--.01 & .304 & \scriptsize .24--.32 & .305 & \scriptsize .22--.34 & .089 & \scriptsize .05--.11 & .177 \\
CLC-Net & Retinal-L. & .048 & \scriptsize .04--.05 & .164 & \scriptsize .08--.16 & .222 & \scriptsize .12--.22 & .173 & \scriptsize .11--.20 & .152 \\
CLC-Net & DiaRetDB1 & .040 & \scriptsize .02--.06 & .318 & \scriptsize .16--.38 & .260 & \scriptsize .16--.23 & .242 & \scriptsize .05--.31 & .215 \\
TC-Net & IDRiD & .225 & \scriptsize .20--.25 & .593 & \scriptsize .47--.67 & .773 & \scriptsize .67--.82 & .636 & \scriptsize .38--.76 & .557 \\
TC-Net & DDR & .112 & \scriptsize .08--.14 & .476 & \scriptsize .40--.54 & .607 & \scriptsize .55--.65 & .250 & \scriptsize .12--.39 & .361 \\
TC-Net & TJDR std & .239 & \scriptsize .12--.34 & .598 & \scriptsize .50--.67 & .670 & \scriptsize .52--.75 & .355 & \scriptsize .09--.52 & .466 \\
TC-Net & TJDR UWF & .100 & \scriptsize .04--.22 & .383 & \scriptsize .20--.53 & .234 & \scriptsize .14--.35 & .056 & \scriptsize .02--.17 & .193 \\
TC-Net & MAPLES & .241 & \scriptsize .17--.30 & .469 & \scriptsize .24--.61 & .722 & \scriptsize .48--.81 & .027 & \scriptsize .00--.09 & .365 \\
TC-Net & FGADR & .018 & \scriptsize .02--.02 & .334 & \scriptsize .30--.37 & .349 & \scriptsize .29--.40 & .156 & \scriptsize .12--.20 & .214 \\
TC-Net & Retinal-L. & .056 & \scriptsize .04--.06 & .206 & \scriptsize .16--.24 & .314 & \scriptsize .22--.35 & .145 & \scriptsize .10--.19 & .180 \\
TC-Net & DiaRetDB1 & .023 & \scriptsize .01--.04 & .475 & \scriptsize .37--.55 & .509 & \scriptsize .36--.53 & .188 & \scriptsize .05--.27 & .299 \\
SAMed & IDRiD & .373 & \scriptsize .33--.42 & .653 & \scriptsize .52--.73 & .845 & \scriptsize .76--.89 & .696 & \scriptsize .49--.80 & .642 \\
SAMed & DDR & .173 & \scriptsize .13--.22 & .574 & \scriptsize .52--.64 & .669 & \scriptsize .62--.72 & .460 & \scriptsize .30--.60 & .469 \\
SAMed & TJDR std & .428 & \scriptsize .29--.51 & .743 & \scriptsize .65--.81 & .707 & \scriptsize .57--.79 & .570 & \scriptsize .26--.73 & .612 \\
SAMed & TJDR UWF & .130 & \scriptsize .05--.28 & .450 & \scriptsize .23--.60 & .208 & \scriptsize .12--.32 & .043 & \scriptsize .01--.19 & .208 \\
SAMed & MAPLES & .365 & \scriptsize .28--.44 & .515 & \scriptsize .30--.65 & .732 & \scriptsize .51--.82 & .227 & \scriptsize .07--.36 & .460 \\
SAMed & FGADR & .027 & \scriptsize .02--.03 & .412 & \scriptsize .37--.45 & .442 & \scriptsize .39--.49 & .279 & \scriptsize .23--.34 & .290 \\
SAMed & Retinal-L. & .109 & \scriptsize .09--.12 & .316 & \scriptsize .26--.37 & .384 & \scriptsize .30--.44 & .304 & \scriptsize .25--.35 & .278 \\
SAMed & DiaRetDB1 & .047 & \scriptsize .02--.07 & .480 & \scriptsize .38--.56 & .468 & \scriptsize .37--.52 & .295 & \scriptsize .12--.43 & .322 \\
SAM-Adapter & IDRiD & .411 & \scriptsize .34--.42 & .644 & \scriptsize .51--.74 & .804 & \scriptsize .71--.85 & .649 & \scriptsize .42--.73 & .627 \\
SAM-Adapter & DDR & .147 & \scriptsize .10--.18 & .538 & \scriptsize .43--.62 & .592 & \scriptsize .51--.63 & .454 & \scriptsize .30--.56 & .433 \\
SAM-Adapter & TJDR std & .358 & \scriptsize .22--.43 & .751 & \scriptsize .67--.82 & .625 & \scriptsize .51--.67 & .432 & \scriptsize .19--.64 & .542 \\
SAM-Adapter & TJDR UWF & .111 & \scriptsize .04--.17 & .387 & \scriptsize .20--.56 & .193 & \scriptsize .09--.31 & .104 & \scriptsize .02--.23 & .199 \\
SAM-Adapter & MAPLES & .277 & \scriptsize .18--.31 & .454 & \scriptsize .25--.57 & .570 & \scriptsize .30--.61 & .187 & \scriptsize .08--.30 & .372 \\
SAM-Adapter & FGADR & .019 & \scriptsize .02--.02 & .384 & \scriptsize .32--.41 & .358 & \scriptsize .29--.41 & .219 & \scriptsize .17--.26 & .245 \\
SAM-Adapter & Retinal-L. & .054 & \scriptsize .04--.05 & .217 & \scriptsize .13--.19 & .273 & \scriptsize .15--.31 & .228 & \scriptsize .18--.25 & .193 \\
SAM-Adapter & DiaRetDB1 & .071 & \scriptsize .04--.10 & .461 & \scriptsize .28--.50 & .412 & \scriptsize .19--.33 & .221 & \scriptsize .05--.30 & .291 \\
SAM2-UNet & IDRiD & .351 & \scriptsize .29--.36 & .524 & \scriptsize .43--.56 & .782 & \scriptsize .68--.81 & .572 & \scriptsize .23--.67 & .557 \\
SAM2-UNet & DDR & .153 & \scriptsize .10--.18 & .493 & \scriptsize .36--.55 & .518 & \scriptsize .41--.56 & .278 & \scriptsize .11--.29 & .360 \\
SAM2-UNet & TJDR std & .325 & \scriptsize .18--.37 & .647 & \scriptsize .55--.71 & .640 & \scriptsize .51--.69 & .416 & \scriptsize .11--.49 & .507 \\
SAM2-UNet & TJDR UWF & .099 & \scriptsize .04--.18 & .382 & \scriptsize .20--.52 & .219 & \scriptsize .10--.30 & .118 & \scriptsize .03--.22 & .204 \\
SAM2-UNet & MAPLES & .331 & \scriptsize .22--.38 & .442 & \scriptsize .23--.55 & .715 & \scriptsize .46--.76 & .104 & \scriptsize .01--.21 & .398 \\
SAM2-UNet & FGADR & .017 & \scriptsize .01--.02 & .356 & \scriptsize .28--.37 & .374 & \scriptsize .29--.39 & .241 & \scriptsize .15--.24 & .247 \\
SAM2-UNet & Retinal-L. & .084 & \scriptsize .06--.08 & .215 & \scriptsize .13--.20 & .322 & \scriptsize .21--.29 & .301 & \scriptsize .21--.31 & .231 \\
SAM2-UNet & DiaRetDB1 & .040 & \scriptsize .02--.06 & .331 & \scriptsize .23--.35 & .409 & \scriptsize .27--.35 & .286 & \scriptsize .08--.44 & .267 \\
MedSAM (full FT) & IDRiD & .331 & \scriptsize .29--.37 & .676 & \scriptsize .55--.76 & .806 & \scriptsize .72--.86 & .522 & \scriptsize .00--.20 & .584 \\
MedSAM (full FT) & DDR & .136 & \scriptsize .10--.19 & .529 & \scriptsize .42--.64 & .626 & \scriptsize .57--.68 & .361 & \scriptsize .07--.23 & .413 \\
MedSAM (full FT) & TJDR std & .347 & \scriptsize .24--.42 & .740 & \scriptsize .65--.81 & .686 & \scriptsize .57--.77 & .440 & \scriptsize .00--.14 & .553 \\
MedSAM (full FT) & TJDR UWF & .129 & \scriptsize .05--.24 & .577 & \scriptsize .34--.70 & .225 & \scriptsize .14--.33 & .041 & \scriptsize .00--.00 & .243 \\
MedSAM (full FT) & MAPLES & .369 & \scriptsize .29--.44 & .515 & \scriptsize .29--.65 & .687 & \scriptsize .43--.81 & .097 & \scriptsize .00--.00 & .417 \\
MedSAM (full FT) & FGADR & .018 & \scriptsize .02--.02 & .419 & \scriptsize .38--.46 & .424 & \scriptsize .37--.48 & .209 & \scriptsize .03--.07 & .268 \\
MedSAM (full FT) & Retinal-L. & .097 & \scriptsize .08--.10 & .359 & \scriptsize .29--.39 & .421 & \scriptsize .33--.48 & .174 & \scriptsize .01--.05 & .262 \\
MedSAM (full FT) & DiaRetDB1 & .045 & \scriptsize .02--.07 & .462 & \scriptsize .36--.55 & .487 & \scriptsize .39--.55 & .141 & \scriptsize .00--.16 & .284 \\
RETFound + UPerNet & IDRiD & .446 & \scriptsize .41--.48 & .681 & \scriptsize .57--.74 & .845 & \scriptsize .76--.88 & .603 & \scriptsize .33--.75 & .644 \\
RETFound + UPerNet & DDR & .189 & \scriptsize .14--.24 & .491 & \scriptsize .39--.60 & .640 & \scriptsize .56--.71 & .546 & \scriptsize .38--.66 & .466 \\
RETFound + UPerNet & TJDR std & .387 & \scriptsize .28--.46 & .733 & \scriptsize .65--.79 & .724 & \scriptsize .56--.81 & .513 & \scriptsize .17--.70 & .589 \\
RETFound + UPerNet & TJDR UWF & .073 & \scriptsize .02--.17 & .064 & \scriptsize .02--.18 & .278 & \scriptsize .18--.39 & .030 & \scriptsize .01--.09 & .111 \\
RETFound + UPerNet & MAPLES & .420 & \scriptsize .35--.47 & .531 & \scriptsize .30--.67 & .722 & \scriptsize .49--.81 & .166 & \scriptsize .03--.31 & .460 \\
RETFound + UPerNet & FGADR & .023 & \scriptsize .02--.03 & .419 & \scriptsize .37--.45 & .406 & \scriptsize .34--.45 & .228 & \scriptsize .17--.29 & .269 \\
RETFound + UPerNet & Retinal-L. & .107 & \scriptsize .08--.10 & .321 & \scriptsize .22--.35 & .419 & \scriptsize .30--.43 & .243 & \scriptsize .17--.29 & .273 \\
RETFound + UPerNet & DiaRetDB1 & .043 & \scriptsize .02--.06 & .430 & \scriptsize .30--.52 & .477 & \scriptsize .33--.48 & .306 & \scriptsize .12--.42 & .314 \\
\end{longtable}
}
\vspace{-6pt}\noindent{\scriptsize TJDR UWF, ultra-widefield images predicted at native resolution with sliding windows. Retinal-L., Retinal-Lesions without its grey-level 127 regions. DiaRetDB1, pixel AUPR against the $\geq$ 0.75 expert consensus (the official image-level protocol is Supplementary Table~13).}
\end{landscape}
\clearpage

{\scriptsize\setlength{\tabcolsep}{3.2pt}\setlength{\LTcapwidth}{\textwidth}\renewcommand{\arraystretch}{1.08}
\begin{longtable}{l c c c c c c c c c c }
\caption{\textbf{Secondary metrics of every main-table model on the internal tests (mean over the four classes and over seeds).} 11-thr., the 11-threshold AUPR of the M2MRF evaluation code on the same maps; Dice and IoU at the threshold chosen on the validation split; lesion-level precision and recall by lesion size (S, M, L: small, medium, large; Methods) at the same threshold; image-level AUC of the maximum probability inside the field of view.}\label{stab:secondary}\\
\toprule
Model & Seeds & \mAUPR{} & 11-thr. & Dice & IoU & Lesion prec. & Recall S & Recall M & Recall L & Image AUC \\
\midrule
\endfirsthead
\multicolumn{11}{l}{\textit{continued}}\\
\toprule
Model & Seeds & \mAUPR{} & 11-thr. & Dice & IoU & Lesion prec. & Recall S & Recall M & Recall L & Image AUC \\
\midrule
\endhead
\midrule
\multicolumn{11}{r}{\textit{continued on the next page}}\\
\endfoot
\bottomrule
\endlastfoot
\multicolumn{11}{l}{\textbf{Trained on IDRiD (internal test; mean over seeds)}}\\
U-Net & 3 & .644\sd{.015} & .638\sd{.013} & .622\sd{.020} & .465\sd{.021} & .489\sd{.023} & .654\sd{.047} & .851\sd{.032} & .800\sd{.040} & .797\sd{.041} \\
DeepLabv3 & 3 & .537\sd{.006} & .536\sd{.004} & .515\sd{.006} & .366\sd{.005} & .551\sd{.019} & .270\sd{.049} & .533\sd{.033} & .776\sd{.026} & \na \\
HRNet-W48 & 3 & .700\sd{.004} & .691\sd{.003} & .627\sd{.004} & .468\sd{.004} & .485\sd{.010} & .584\sd{.040} & .879\sd{.013} & .862\sd{.018} & \na \\
Swin-Unet & 3 & .688\sd{.003} & .686\sd{.002} & .640\sd{.002} & .481\sd{.002} & .419\sd{.007} & .671\sd{.050} & .849\sd{.008} & .827\sd{.014} & \na \\
FCT & 3 & .298\sd{.084} & .359\sd{.071} & .330\sd{.066} & .217\sd{.048} & .267\sd{.047} & .205\sd{.032} & .523\sd{.029} & .726\sd{.027} & \na \\
nnU-Net & 3 & .665\sd{.008} & .675\sd{.006} & .656\sd{.009} & .498\sd{.010} & .552\sd{.016} & .649\sd{.044} & .828\sd{.003} & .739\sd{.004} & .786\sd{.016} \\
nnU-Net ResEnc-M & 1 & .631 & .631 & .638 & .479 & .526 & .600 & .852 & .695 & \na \\
M2MRF & 3 & .657\sd{.010} & .650\sd{.004} & .648\sd{.004} & .490\sd{.004} & .531\sd{.008} & .615\sd{.015} & .824\sd{.009} & .736\sd{.018} & \na \\
HACDR-Net & 3 & .538\sd{.031} & .559\sd{.013} & .520\sd{.025} & .367\sd{.022} & .495\sd{.040} & .411\sd{.019} & .749\sd{.046} & .764\sd{.014} & .698\sd{.163} \\
HRDecoder & 1 & .647 & .637 & .636 & .476 & .605 & .456 & .792 & .785 & .885 \\
PMCNet & 1 & .652 & .648 & .616 & .458 & .529 & .592 & .825 & .804 & .786 \\
WSRFNet & 1 & .582 & .605 & .641 & .482 & .538 & .611 & .808 & .701 & .852 \\
CLC-Net & 1 & .661 & .665 & .632 & .475 & .495 & .565 & .832 & .864 & .868 \\
TC-Net & 1 & .653 & .643 & .607 & .449 & .536 & .443 & .804 & .815 & .874 \\
SAMed & 3 & .706\sd{.015} & .703\sd{.012} & .658\sd{.007} & .502\sd{.008} & .519\sd{.007} & .553\sd{.055} & .798\sd{.033} & .838\sd{.030} & .907\sd{.033} \\
SAM-Adapter & 1 & .706 & .704 & .663 & .508 & .527 & .624 & .851 & .809 & .832 \\
SAM2-UNet & 1 & .664 & .683 & .653 & .495 & .440 & .846 & .925 & .826 & .835 \\
MedSAM (full FT) & 1 & .642 & .654 & .609 & .451 & .507 & .392 & .759 & .791 & .791 \\
RETFound + UPerNet & 3 & .699\sd{.009} & .692\sd{.008} & .650\sd{.008} & .492\sd{.009} & .469\sd{.014} & .683\sd{.015} & .873\sd{.014} & .856\sd{.028} & .821\sd{.040} \\
\multicolumn{11}{l}{\textbf{Trained on DDR (internal test; mean over seeds)}}\\
U-Net & 3 & .371\sd{.009} & .421\sd{.008} & .405\sd{.017} & .260\sd{.014} & .435\sd{.042} & .247\sd{.016} & .482\sd{.012} & .646\sd{.009} & .786\sd{.008} \\
DeepLabv3 & 3 & .377\sd{.005} & .382\sd{.003} & .380\sd{.006} & .246\sd{.004} & .497\sd{.015} & .062\sd{.005} & .287\sd{.020} & .655\sd{.010} & \na \\
HRNet-W48 & 3 & .447\sd{.009} & .461\sd{.009} & .448\sd{.004} & .297\sd{.004} & .513\sd{.006} & .217\sd{.007} & .495\sd{.006} & .680\sd{.012} & \na \\
Swin-Unet & 3 & .431\sd{.001} & .447\sd{.003} & .440\sd{.003} & .289\sd{.003} & .424\sd{.008} & .312\sd{.012} & .543\sd{.013} & .674\sd{.002} & \na \\
FCT & 1 & .176 & .344 & .197 & .125 & .224 & .096 & .351 & .608 & \na \\
nnU-Net & 3 & .431\sd{.003} & .467\sd{.001} & .458\sd{.010} & .304\sd{.008} & .515\sd{.008} & .258\sd{.006} & .470\sd{.006} & .622\sd{.007} & .796\sd{.004} \\
nnU-Net ResEnc-M & 1 & .428 & .448 & .463 & .310 & .491 & .294 & .490 & .664 & \na \\
M2MRF & 3 & .369\sd{.011} & .406\sd{.008} & .407\sd{.012} & .262\sd{.009} & .454\sd{.007} & .276\sd{.003} & .556\sd{.001} & .696\sd{.023} & \na \\
HACDR-Net & 3 & .252\sd{.017} & .289\sd{.022} & .289\sd{.022} & .175\sd{.014} & .416\sd{.019} & .106\sd{.022} & .350\sd{.028} & .562\sd{.016} & .719\sd{.012} \\
HRDecoder & 1 & .296 & .333 & .373 & .236 & .459 & .234 & .496 & .629 & .714 \\
PMCNet & 1 & .312 & .342 & .338 & .206 & .431 & .200 & .454 & .629 & .770 \\
WSRFNet & 1 & .234 & .264 & .341 & .211 & .431 & .284 & .458 & .531 & .659 \\
CLC-Net & 1 & .414 & .433 & .462 & .309 & .506 & .266 & .470 & .639 & .781 \\
TC-Net & 1 & .361 & .365 & .352 & .223 & .443 & .102 & .376 & .634 & .785 \\
SAMed & 3 & .460\sd{.027} & .472\sd{.027} & .454\sd{.025} & .301\sd{.022} & .497\sd{.008} & .242\sd{.029} & .474\sd{.031} & .674\sd{.007} & .842\sd{.005} \\
SAM-Adapter & 1 & .433 & .465 & .457 & .304 & .471 & .257 & .496 & .690 & .657 \\
SAM2-UNet & 1 & .360 & .445 & .413 & .265 & .459 & .298 & .559 & .698 & .749 \\
MedSAM (full FT) & 1 & .413 & .433 & .432 & .283 & .516 & .212 & .484 & .675 & .795 \\
RETFound + UPerNet & 3 & .450\sd{.018} & .458\sd{.016} & .447\sd{.006} & .294\sd{.005} & .472\sd{.006} & .290\sd{.006} & .538\sd{.007} & .696\sd{.009} & .844\sd{.002} \\
\end{longtable}
}
\vspace{-6pt}\noindent{\scriptsize $\pm$, standard deviation over seeds where more than one seed was run.}

\clearpage

{\scriptsize\setlength{\tabcolsep}{3.2pt}\setlength{\LTcapwidth}{\textwidth}\renewcommand{\arraystretch}{1.08}
\begin{longtable}{l l l c c c c c}
\caption{\textbf{Registered and revised recipes of the three models whose per-image loss term emptied a rare-class channel (seed 0, internal tests).} The registered recipe is the one frozen before training (per-image Dice or IoU with smoothing 1); the revised recipe pools the term over the batch (Methods, rare-class collapse). The main tables carry the revised runs; the registered runs stay in the repository under the \texttt{\_registered} suffix.}\label{stab:registered}\\
\toprule
Model & Source & Recipe & MA & HE & EX & SE & \mAUPR{} \\
\midrule
\endfirsthead
\multicolumn{8}{l}{\textit{continued}}\\
\toprule
Model & Source & Recipe & MA & HE & EX & SE & \mAUPR{} \\
\midrule
\endhead
\midrule
\multicolumn{8}{r}{\textit{continued on the next page}}\\
\endfoot
\bottomrule
\endlastfoot
U-Net & IDRiD & registered & .376 & .534 & .859 & .605 & .594 \\
U-Net & IDRiD & revised (main table) & .406 & .622 & .896 & .707 & .658 \\
U-Net & DDR & registered & .005 & .031 & .418 & .000 & .114 \\
U-Net & DDR & revised (main table) & .141 & .410 & .553 & .410 & .378 \\
SAM2-UNet & IDRiD & registered & .439 & .665 & .870 & .571 & .636 \\
SAM2-UNet & IDRiD & revised (main table) & .439 & .670 & .871 & .678 & .664 \\
SAM2-UNet & DDR & registered & .023 & .265 & .528 & .022 & .210 \\
SAM2-UNet & DDR & revised (main table) & .153 & .493 & .518 & .278 & .360 \\
MedSAM & IDRiD & registered & .382 & .731 & .867 & .357 & .584 \\
MedSAM & IDRiD & revised (main table) & .373 & .665 & .855 & .674 & .642 \\
MedSAM & DDR & registered & .000 & .487 & .573 & .001 & .265 \\
MedSAM & DDR & revised (main table) & .136 & .529 & .626 & .361 & .413 \\
\end{longtable}
}

\clearpage
\begin{landscape}
{\scriptsize\setlength{\tabcolsep}{3.2pt}\setlength{\LTcapwidth}{\textwidth}\renewcommand{\arraystretch}{1.08}
\begin{longtable}{l l l l c c c c c}
\caption{\textbf{Every factor cell, per class (seed 0).} One factor is varied at a time on the representative models at the DDR pixel budget (Methods); the variant column names the level, the training-set column the data the cell was trained on. Each cell is listed on every test set it was applied to; in-source tests of merged cells are the rows whose test set was part of the training set.}\label{stab:factors}\\
\toprule
Model & Level & Trained on & Test set & MA & HE & EX & SE & \mAUPR{} \\
\midrule
\endfirsthead
\multicolumn{9}{l}{\textit{continued}}\\
\toprule
Model & Level & Trained on & Test set & MA & HE & EX & SE & \mAUPR{} \\
\midrule
\endhead
\midrule
\multicolumn{9}{r}{\textit{continued on the next page}}\\
\endfoot
\bottomrule
\endlastfoot
\multicolumn{9}{l}{\textbf{Training source (E3b)}}\\
M2MRF & bstd & FGADR & IDRiD & .141 & .320 & .729 & .557 & .437 \\
M2MRF & bstd & FGADR & DDR & .055 & .219 & .404 & .107 & .197 \\
M2MRF & bstd & FGADR & TJDR std & .267 & .642 & .666 & .223 & .449 \\
M2MRF & bstd & FGADR & TJDR UWF & .035 & .480 & .282 & .176 & .243 \\
M2MRF & bstd & FGADR & MAPLES & .086 & .336 & .492 & .090 & .251 \\
M2MRF & bstd & FGADR & FGADR & .066 & .511 & .622 & .238 & .359 \\
M2MRF & bstd & FGADR & Retinal-L. & .031 & .179 & .177 & .167 & .139 \\
M2MRF & bstd & FGADR & DiaRetDB1 & .006 & .270 & .287 & .142 & .176 \\
M2MRF & bstd & TJDR-std & IDRiD & .139 & .186 & .719 & .317 & .340 \\
M2MRF & bstd & TJDR-std & DDR & .082 & .221 & .516 & .104 & .231 \\
M2MRF & bstd & TJDR-std & TJDR std & .325 & .746 & .654 & .490 & .554 \\
M2MRF & bstd & TJDR-std & TJDR UWF & .022 & .270 & .131 & .157 & .145 \\
M2MRF & bstd & TJDR-std & MAPLES & .117 & .340 & .375 & .022 & .214 \\
M2MRF & bstd & TJDR-std & FGADR & .014 & .161 & .270 & .099 & .136 \\
M2MRF & bstd & TJDR-std & Retinal-L. & .030 & .103 & .231 & .084 & .112 \\
M2MRF & bstd & TJDR-std & DiaRetDB1 & .018 & .201 & .295 & .134 & .162 \\
RETFound + UPerNet & bstd & FGADR & IDRiD & .243 & .563 & .775 & .582 & .541 \\
RETFound + UPerNet & bstd & FGADR & DDR & .019 & .392 & .470 & .284 & .291 \\
RETFound + UPerNet & bstd & FGADR & TJDR std & .140 & .752 & .595 & .357 & .461 \\
RETFound + UPerNet & bstd & FGADR & TJDR UWF & .073 & .615 & .336 & .113 & .284 \\
RETFound + UPerNet & bstd & FGADR & MAPLES & .146 & .492 & .722 & .118 & .370 \\
RETFound + UPerNet & bstd & FGADR & FGADR & .118 & .651 & .691 & .370 & .457 \\
RETFound + UPerNet & bstd & FGADR & Retinal-L. & .115 & .385 & .467 & .325 & .323 \\
RETFound + UPerNet & bstd & FGADR & DiaRetDB1 & .019 & .523 & .590 & .178 & .327 \\
RETFound + UPerNet & bstd & TJDR-std & IDRiD & .315 & .506 & .815 & .275 & .478 \\
RETFound + UPerNet & bstd & TJDR-std & DDR & .136 & .347 & .614 & .236 & .333 \\
RETFound + UPerNet & bstd & TJDR-std & TJDR std & .432 & .823 & .828 & .741 & .706 \\
RETFound + UPerNet & bstd & TJDR-std & TJDR UWF & .053 & .523 & .150 & .262 & .247 \\
RETFound + UPerNet & bstd & TJDR-std & MAPLES & .310 & .330 & .605 & .131 & .344 \\
RETFound + UPerNet & bstd & TJDR-std & FGADR & .023 & .327 & .368 & .185 & .226 \\
RETFound + UPerNet & bstd & TJDR-std & Retinal-L. & .085 & .226 & .336 & .125 & .193 \\
RETFound + UPerNet & bstd & TJDR-std & DiaRetDB1 & .027 & .317 & .491 & .211 & .262 \\
SAM2-UNet & bstd & FGADR & IDRiD & .138 & .560 & .859 & .099 & .414 \\
SAM2-UNet & bstd & FGADR & DDR & .062 & .385 & .537 & .092 & .269 \\
SAM2-UNet & bstd & FGADR & TJDR std & .320 & .752 & .768 & .001 & .461 \\
SAM2-UNet & bstd & FGADR & TJDR UWF & .060 & .593 & .270 & .000 & .231 \\
SAM2-UNet & bstd & FGADR & MAPLES & .150 & .477 & .744 & .000 & .343 \\
SAM2-UNet & bstd & FGADR & FGADR & .122 & .630 & .675 & .126 & .388 \\
SAM2-UNet & bstd & FGADR & Retinal-L. & .043 & .339 & .411 & .048 & .210 \\
SAM2-UNet & bstd & FGADR & DiaRetDB1 & .012 & .468 & .516 & .082 & .270 \\
SAM2-UNet & bstd & TJDR-std & IDRiD & .025 & .439 & .793 & .026 & .321 \\
SAM2-UNet & bstd & TJDR-std & DDR & .016 & .325 & .542 & .013 & .224 \\
SAM2-UNet & bstd & TJDR-std & TJDR std & .056 & .753 & .722 & .102 & .408 \\
SAM2-UNet & bstd & TJDR-std & TJDR UWF & .008 & .380 & .095 & .010 & .123 \\
SAM2-UNet & bstd & TJDR-std & MAPLES & .007 & .281 & .480 & .000 & .192 \\
SAM2-UNet & bstd & TJDR-std & FGADR & .007 & .294 & .308 & .005 & .154 \\
SAM2-UNet & bstd & TJDR-std & Retinal-L. & .010 & .137 & .201 & .006 & .088 \\
SAM2-UNet & bstd & TJDR-std & DiaRetDB1 & .000 & .233 & .230 & .006 & .117 \\
U-Net & bstd & FGADR & IDRiD & .092 & .294 & .793 & .276 & .364 \\
U-Net & bstd & FGADR & DDR & .032 & .259 & .467 & .199 & .239 \\
U-Net & bstd & FGADR & TJDR std & .135 & .718 & .603 & .339 & .449 \\
U-Net & bstd & FGADR & TJDR UWF & .028 & .482 & .208 & .273 & .248 \\
U-Net & bstd & FGADR & MAPLES & .088 & .337 & .566 & .064 & .264 \\
U-Net & bstd & FGADR & FGADR & .084 & .568 & .634 & .325 & .403 \\
U-Net & bstd & FGADR & Retinal-L. & .084 & .224 & .368 & .155 & .208 \\
U-Net & bstd & FGADR & DiaRetDB1 & .005 & .261 & .563 & .188 & .254 \\
U-Net & bstd & TJDR-std & IDRiD & .108 & .324 & .641 & .435 & .377 \\
U-Net & bstd & TJDR-std & DDR & .071 & .264 & .531 & .161 & .257 \\
U-Net & bstd & TJDR-std & TJDR std & .326 & .775 & .642 & .455 & .550 \\
U-Net & bstd & TJDR-std & TJDR UWF & .015 & .203 & .077 & .182 & .119 \\
U-Net & bstd & TJDR-std & MAPLES & .080 & .327 & .447 & .033 & .222 \\
U-Net & bstd & TJDR-std & FGADR & .015 & .237 & .266 & .076 & .148 \\
U-Net & bstd & TJDR-std & Retinal-L. & .028 & .107 & .193 & .098 & .107 \\
U-Net & bstd & TJDR-std & DiaRetDB1 & .008 & .169 & .217 & .144 & .134 \\
nnU-Net & bstd & FGADR & IDRiD & .121 & .596 & .826 & .717 & .565 \\
nnU-Net & bstd & FGADR & DDR & .063 & .378 & .532 & .306 & .320 \\
nnU-Net & bstd & FGADR & TJDR std & .147 & .737 & .696 & .520 & .525 \\
nnU-Net & bstd & FGADR & TJDR UWF & .021 & .277 & .222 & .391 & .228 \\
nnU-Net & bstd & FGADR & MAPLES & .197 & .517 & .719 & .138 & .393 \\
nnU-Net & bstd & FGADR & FGADR & .136 & .631 & .659 & .372 & .449 \\
nnU-Net & bstd & FGADR & Retinal-L. & .039 & .390 & .439 & .247 & .279 \\
nnU-Net & bstd & FGADR & DiaRetDB1 & .028 & .474 & .559 & .226 & .322 \\
nnU-Net & bstd & TJDR-std & IDRiD & .146 & .459 & .664 & .474 & .436 \\
nnU-Net & bstd & TJDR-std & DDR & .079 & .310 & .550 & .170 & .277 \\
nnU-Net & bstd & TJDR-std & TJDR std & .305 & .785 & .787 & .668 & .636 \\
nnU-Net & bstd & TJDR-std & TJDR UWF & .013 & .311 & .124 & .268 & .179 \\
nnU-Net & bstd & TJDR-std & MAPLES & .162 & .462 & .457 & .033 & .278 \\
nnU-Net & bstd & TJDR-std & FGADR & .019 & .322 & .299 & .157 & .199 \\
nnU-Net & bstd & TJDR-std & Retinal-L. & .031 & .162 & .224 & .125 & .136 \\
nnU-Net & bstd & TJDR-std & DiaRetDB1 & .027 & .395 & .338 & .157 & .229 \\
\multicolumn{9}{l}{\textbf{Annotation granularity (E4)}}\\
M2MRF & DDR & DDR & IDRiD & .392 & .555 & .776 & .563 & .571 \\
M2MRF & DDR & DDR & TJDR std & .299 & .696 & .618 & .412 & .506 \\
M2MRF & DDR & DDR & DiaRetDB1 & .049 & .218 & .177 & .247 & .173 \\
U-Net & DDR & DDR & IDRiD & .364 & .541 & .800 & .483 & .547 \\
U-Net & DDR & DDR & TJDR std & .296 & .676 & .614 & .392 & .494 \\
U-Net & DDR & DDR & DiaRetDB1 & .039 & .314 & .268 & .253 & .219 \\
M2MRF & FGADR & FGADR & IDRiD & .141 & .320 & .729 & .557 & .437 \\
M2MRF & FGADR & FGADR & TJDR std & .267 & .642 & .666 & .223 & .449 \\
M2MRF & FGADR & FGADR & DiaRetDB1 & .006 & .270 & .287 & .142 & .176 \\
U-Net & FGADR & FGADR & IDRiD & .092 & .294 & .793 & .276 & .364 \\
U-Net & FGADR & FGADR & TJDR std & .135 & .718 & .603 & .339 & .449 \\
U-Net & FGADR & FGADR & DiaRetDB1 & .005 & .261 & .563 & .188 & .254 \\
M2MRF & DDR + FGADR (383) & DDR + FGADR (383) & IDRiD & .386 & .446 & .799 & .594 & .556 \\
M2MRF & DDR + FGADR (383) & DDR + FGADR (383) & DDR & .141 & .365 & .508 & .328 & .336 \\
M2MRF & DDR + FGADR (383) & DDR + FGADR (383) & TJDR std & .285 & .669 & .661 & .321 & .484 \\
M2MRF & DDR + FGADR (383) & DDR + FGADR (383) & TJDR UWF & .110 & .461 & .172 & .180 & .231 \\
M2MRF & DDR + FGADR (383) & DDR + FGADR (383) & MAPLES & .252 & .378 & .482 & .134 & .311 \\
M2MRF & DDR + FGADR (383) & DDR + FGADR (383) & Retinal-L. & .044 & .158 & .177 & .176 & .139 \\
M2MRF & DDR + FGADR (383) & DDR + FGADR (383) & DiaRetDB1 & .043 & .197 & .222 & .182 & .161 \\
M2MRF & DDR + FGADR & DDR + FGADR & IDRiD & .403 & .466 & .782 & .565 & .554 \\
M2MRF & DDR + FGADR & DDR + FGADR & DDR & .139 & .366 & .492 & .282 & .320 \\
M2MRF & DDR + FGADR & DDR + FGADR & TJDR std & .312 & .680 & .659 & .337 & .497 \\
M2MRF & DDR + FGADR & DDR + FGADR & TJDR UWF & .109 & .474 & .186 & .264 & .258 \\
M2MRF & DDR + FGADR & DDR + FGADR & MAPLES & .256 & .391 & .431 & .119 & .299 \\
M2MRF & DDR + FGADR & DDR + FGADR & Retinal-L. & .045 & .168 & .173 & .160 & .137 \\
M2MRF & DDR + FGADR & DDR + FGADR & DiaRetDB1 & .041 & .208 & .176 & .151 & .144 \\
M2MRF & DDR + Retinal-Lesions & DDR + Retinal-Lesions & IDRiD & .369 & .399 & .784 & .508 & .515 \\
M2MRF & DDR + Retinal-Lesions & DDR + Retinal-Lesions & DDR & .108 & .296 & .449 & .291 & .286 \\
M2MRF & DDR + Retinal-Lesions & DDR + Retinal-Lesions & TJDR std & .145 & .480 & .378 & .299 & .325 \\
M2MRF & DDR + Retinal-Lesions & DDR + Retinal-Lesions & TJDR UWF & .086 & .419 & .223 & .080 & .202 \\
M2MRF & DDR + Retinal-Lesions & DDR + Retinal-Lesions & MAPLES & .131 & .270 & .408 & .113 & .231 \\
M2MRF & DDR + Retinal-Lesions & DDR + Retinal-Lesions & FGADR & .010 & .157 & .186 & .149 & .126 \\
M2MRF & DDR + Retinal-Lesions & DDR + Retinal-Lesions & DiaRetDB1 & .048 & .199 & .308 & .185 & .185 \\
M2MRF & Retinal-Lesions & Retinal-Lesions & IDRiD & .063 & .133 & .186 & .285 & .167 \\
M2MRF & Retinal-Lesions & Retinal-Lesions & DDR & .024 & .097 & .130 & .166 & .104 \\
M2MRF & Retinal-Lesions & Retinal-Lesions & TJDR std & .039 & .112 & .160 & .287 & .150 \\
M2MRF & Retinal-Lesions & Retinal-Lesions & TJDR UWF & .015 & .121 & .053 & .120 & .077 \\
M2MRF & Retinal-Lesions & Retinal-Lesions & MAPLES & .110 & .101 & .342 & .074 & .157 \\
M2MRF & Retinal-Lesions & Retinal-Lesions & FGADR & .018 & .121 & .131 & .129 & .100 \\
M2MRF & Retinal-Lesions & Retinal-Lesions & Retinal-L. & .352 & .508 & .650 & .366 & .469 \\
M2MRF & Retinal-Lesions & Retinal-Lesions & DiaRetDB1 & .022 & .328 & .774 & .355 & .370 \\
U-Net & DDR + FGADR (383) & DDR + FGADR (383) & IDRiD & .375 & .517 & .828 & .594 & .578 \\
U-Net & DDR + FGADR (383) & DDR + FGADR (383) & DDR & .137 & .414 & .548 & .376 & .369 \\
U-Net & DDR + FGADR (383) & DDR + FGADR (383) & TJDR std & .299 & .728 & .556 & .177 & .440 \\
U-Net & DDR + FGADR (383) & DDR + FGADR (383) & TJDR UWF & .101 & .473 & .151 & .242 & .242 \\
U-Net & DDR + FGADR (383) & DDR + FGADR (383) & MAPLES & .305 & .472 & .502 & .057 & .334 \\
U-Net & DDR + FGADR (383) & DDR + FGADR (383) & Retinal-L. & .068 & .244 & .258 & .177 & .187 \\
U-Net & DDR + FGADR (383) & DDR + FGADR (383) & DiaRetDB1 & .045 & .329 & .352 & .161 & .221 \\
U-Net & DDR + FGADR & DDR + FGADR & IDRiD & .365 & .551 & .858 & .613 & .597 \\
U-Net & DDR + FGADR & DDR + FGADR & DDR & .117 & .430 & .545 & .402 & .373 \\
U-Net & DDR + FGADR & DDR + FGADR & TJDR std & .325 & .762 & .622 & .342 & .513 \\
U-Net & DDR + FGADR & DDR + FGADR & TJDR UWF & .089 & .466 & .210 & .294 & .265 \\
U-Net & DDR + FGADR & DDR + FGADR & MAPLES & .257 & .468 & .575 & .133 & .358 \\
U-Net & DDR + FGADR & DDR + FGADR & Retinal-L. & .042 & .219 & .236 & .186 & .171 \\
U-Net & DDR + FGADR & DDR + FGADR & DiaRetDB1 & .034 & .320 & .302 & .173 & .207 \\
U-Net & DDR + Retinal-Lesions & DDR + Retinal-Lesions & IDRiD & .279 & .502 & .786 & .521 & .522 \\
U-Net & DDR + Retinal-Lesions & DDR + Retinal-Lesions & DDR & .090 & .366 & .518 & .389 & .341 \\
U-Net & DDR + Retinal-Lesions & DDR + Retinal-Lesions & TJDR std & .237 & .643 & .490 & .440 & .452 \\
U-Net & DDR + Retinal-Lesions & DDR + Retinal-Lesions & TJDR UWF & .057 & .480 & .231 & .210 & .245 \\
U-Net & DDR + Retinal-Lesions & DDR + Retinal-Lesions & MAPLES & .269 & .465 & .474 & .147 & .339 \\
U-Net & DDR + Retinal-Lesions & DDR + Retinal-Lesions & FGADR & .013 & .202 & .271 & .133 & .155 \\
U-Net & DDR + Retinal-Lesions & DDR + Retinal-Lesions & DiaRetDB1 & .052 & .278 & .386 & .278 & .248 \\
U-Net & Retinal-Lesions & Retinal-Lesions & IDRiD & .091 & .230 & .353 & .317 & .248 \\
U-Net & Retinal-Lesions & Retinal-Lesions & DDR & .039 & .146 & .239 & .234 & .164 \\
U-Net & Retinal-Lesions & Retinal-Lesions & TJDR std & .087 & .378 & .238 & .419 & .281 \\
U-Net & Retinal-Lesions & Retinal-Lesions & TJDR UWF & .026 & .208 & .091 & .163 & .122 \\
U-Net & Retinal-Lesions & Retinal-Lesions & MAPLES & .168 & .284 & .470 & .102 & .256 \\
U-Net & Retinal-Lesions & Retinal-Lesions & FGADR & .017 & .136 & .148 & .064 & .091 \\
U-Net & Retinal-Lesions & Retinal-Lesions & Retinal-L. & .353 & .553 & .721 & .390 & .504 \\
U-Net & Retinal-Lesions & Retinal-Lesions & DiaRetDB1 & .045 & .494 & .874 & .483 & .474 \\
\multicolumn{9}{l}{\textbf{Input resolution (E7)}}\\
M2MRF & res:1440 & DDR & DDR & .148 & .402 & .547 & .422 & .380 \\
RETFound + UPerNet & res:1440 & DDR & DDR & .189 & .491 & .640 & .546 & .466 \\
HRDecoder & res:2880 (author canvas, equivalent epochs) & IDRiD & IDRiD & .423 & .606 & .864 & .615 & .627 \\
M2MRF & res:512 & DDR & DDR & .114 & .411 & .531 & .378 & .358 \\
M2MRF & res:native patches & DDR & DDR & .140 & .360 & .510 & .237 & .312 \\
RETFound + UPerNet & res:512 & DDR & DDR & .092 & .346 & .628 & .535 & .400 \\
RETFound + UPerNet & res:native patches & DDR & DDR & .196 & .460 & .628 & .513 & .449 \\
\multicolumn{9}{l}{\textbf{Training amount (E8)}}\\
M2MRF & size:DDR 100\% & DDR & DDR & .148 & .402 & .547 & .422 & .380 \\
M2MRF & size:DDR 100\% & DDR & TJDR std & .299 & .696 & .618 & .412 & .506 \\
M2MRF & size:DDR 100\% & DDR & FGADR & .019 & .310 & .340 & .181 & .212 \\
U-Net & size:DDR 100\% & DDR & DDR & .141 & .410 & .553 & .410 & .378 \\
U-Net & size:DDR 100\% & DDR & TJDR std & .296 & .676 & .614 & .392 & .494 \\
U-Net & size:DDR 100\% & DDR & FGADR & .013 & .214 & .292 & .162 & .170 \\
M2MRF & size:DDR 10\% & DDR & IDRiD & .264 & .481 & .735 & .413 & .473 \\
M2MRF & size:DDR 10\% & DDR & DDR & .107 & .383 & .516 & .222 & .307 \\
M2MRF & size:DDR 10\% & DDR & TJDR std & .211 & .625 & .579 & .225 & .410 \\
M2MRF & size:DDR 10\% & DDR & TJDR UWF & .102 & .221 & .167 & .032 & .131 \\
M2MRF & size:DDR 10\% & DDR & MAPLES & .188 & .345 & .446 & .067 & .262 \\
M2MRF & size:DDR 10\% & DDR & FGADR & .011 & .242 & .187 & .148 & .147 \\
M2MRF & size:DDR 10\% & DDR & Retinal-L. & .036 & .139 & .164 & .138 & .119 \\
M2MRF & size:DDR 10\% & DDR & DiaRetDB1 & .030 & .190 & .186 & .205 & .153 \\
M2MRF & size:DDR 50\% & DDR & IDRiD & .364 & .500 & .768 & .566 & .549 \\
M2MRF & size:DDR 50\% & DDR & DDR & .140 & .364 & .532 & .332 & .342 \\
M2MRF & size:DDR 50\% & DDR & TJDR std & .265 & .672 & .603 & .450 & .498 \\
M2MRF & size:DDR 50\% & DDR & TJDR UWF & .103 & .365 & .234 & .045 & .187 \\
M2MRF & size:DDR 50\% & DDR & MAPLES & .249 & .427 & .478 & .127 & .320 \\
M2MRF & size:DDR 50\% & DDR & FGADR & .014 & .245 & .338 & .184 & .195 \\
M2MRF & size:DDR 50\% & DDR & Retinal-L. & .046 & .166 & .184 & .177 & .143 \\
M2MRF & size:DDR 50\% & DDR & DiaRetDB1 & .048 & .220 & .193 & .246 & .177 \\
M2MRF & size:fine merged & DDR & IDRiD & .446 & .531 & .824 & .674 & .619 \\
M2MRF & size:fine merged & DDR & DDR & .141 & .393 & .534 & .405 & .368 \\
M2MRF & size:fine merged & DDR & TJDR std & .333 & .746 & .640 & .540 & .565 \\
M2MRF & size:fine merged & DDR & TJDR UWF & .120 & .472 & .197 & .339 & .282 \\
M2MRF & size:fine merged & DDR & MAPLES & .285 & .462 & .500 & .150 & .349 \\
M2MRF & size:fine merged & DDR & FGADR & .020 & .347 & .333 & .194 & .224 \\
M2MRF & size:fine merged & DDR & Retinal-L. & .049 & .191 & .184 & .187 & .153 \\
M2MRF & size:fine merged & DDR & DiaRetDB1 & .049 & .266 & .191 & .256 & .190 \\
U-Net & size:DDR 10\% & DDR & IDRiD & .239 & .474 & .784 & .470 & .492 \\
U-Net & size:DDR 10\% & DDR & DDR & .106 & .337 & .508 & .334 & .321 \\
U-Net & size:DDR 10\% & DDR & TJDR std & .197 & .525 & .479 & .308 & .378 \\
U-Net & size:DDR 10\% & DDR & TJDR UWF & .077 & .243 & .205 & .072 & .149 \\
U-Net & size:DDR 10\% & DDR & MAPLES & .243 & .303 & .542 & .092 & .295 \\
U-Net & size:DDR 10\% & DDR & FGADR & .010 & .132 & .190 & .119 & .112 \\
U-Net & size:DDR 10\% & DDR & Retinal-L. & .037 & .119 & .173 & .165 & .123 \\
U-Net & size:DDR 10\% & DDR & DiaRetDB1 & .030 & .180 & .247 & .204 & .165 \\
U-Net & size:DDR 50\% & DDR & IDRiD & .303 & .552 & .811 & .493 & .540 \\
U-Net & size:DDR 50\% & DDR & DDR & .132 & .407 & .570 & .383 & .373 \\
U-Net & size:DDR 50\% & DDR & TJDR std & .273 & .715 & .530 & .374 & .473 \\
U-Net & size:DDR 50\% & DDR & TJDR UWF & .123 & .338 & .261 & .075 & .199 \\
U-Net & size:DDR 50\% & DDR & MAPLES & .240 & .476 & .564 & .076 & .339 \\
U-Net & size:DDR 50\% & DDR & FGADR & .008 & .230 & .355 & .155 & .187 \\
U-Net & size:DDR 50\% & DDR & Retinal-L. & .045 & .192 & .214 & .189 & .160 \\
U-Net & size:DDR 50\% & DDR & DiaRetDB1 & .035 & .307 & .262 & .234 & .210 \\
U-Net & size:fine merged & DDR & IDRiD & .432 & .563 & .859 & .574 & .607 \\
U-Net & size:fine merged & DDR & DDR & .154 & .393 & .582 & .337 & .367 \\
U-Net & size:fine merged & DDR & TJDR std & .367 & .779 & .643 & .458 & .562 \\
U-Net & size:fine merged & DDR & TJDR UWF & .118 & .489 & .179 & .288 & .268 \\
U-Net & size:fine merged & DDR & MAPLES & .295 & .491 & .596 & .107 & .372 \\
U-Net & size:fine merged & DDR & FGADR & .014 & .276 & .322 & .154 & .192 \\
U-Net & size:fine merged & DDR & Retinal-L. & .050 & .229 & .240 & .191 & .177 \\
U-Net & size:fine merged & DDR & DiaRetDB1 & .041 & .338 & .292 & .202 & .218 \\
M2MRF & size:DDR 10\%:fov & DDR & TJDR UWF & .087 & .304 & .076 & .096 & .141 \\
M2MRF & size:DDR 50\%:fov & DDR & TJDR UWF & .115 & .432 & .160 & .163 & .218 \\
U-Net & size:DDR 10\%:fov & DDR & TJDR UWF & .047 & .238 & .080 & .121 & .122 \\
U-Net & size:DDR 50\%:fov & DDR & TJDR UWF & .065 & .412 & .187 & .221 & .221 \\
\multicolumn{9}{l}{\textbf{Loss (E9)}}\\
U-Net & loss:CE + Dice & DDR & DDR & .141 & .410 & .553 & .410 & .378 \\
U-Net & loss:CE & DDR & DDR & .126 & .429 & .602 & .271 & .357 \\
U-Net & loss:Dice & DDR & DDR & .124 & .268 & .397 & .160 & .237 \\
U-Net & loss:CE + per-image Dice & DDR & DDR & .006 & .010 & .460 & .000 & .119 \\
\end{longtable}
}
\end{landscape}
\clearpage

{\scriptsize\setlength{\tabcolsep}{3.2pt}\setlength{\LTcapwidth}{\textwidth}\renewcommand{\arraystretch}{1.08}
\begin{longtable}{l c c c c c c c}
\caption{\textbf{Computational cost (E10).} Parameters and trainable parameters under the unified protocol, training time of one seed on each source (40\,GB A100, one task per card), and, where the cost audit has run, floating-point operations for one 1,440-pixel canvas, peak training memory and inference time per image.}\label{stab:cost}\\
\toprule
Model & Params (M) & Trainable (M) & GPU h IDRiD & GPU h DDR & GFLOPs & Peak mem.\ (GB) & Infer.\ (s/image) \\
\midrule
\endfirsthead
\multicolumn{8}{l}{\textit{continued}}\\
\toprule
Model & Params (M) & Trainable (M) & GPU h IDRiD & GPU h DDR & GFLOPs & Peak mem.\ (GB) & Infer.\ (s/image) \\
\midrule
\endhead
\midrule
\multicolumn{8}{r}{\textit{continued on the next page}}\\
\endfoot
\bottomrule
\endlastfoot
U-Net & 7.8 & \na & 1.5 & 5.4 & \na & \na & \na \\
DeepLabv3 & 58.6 & \na & 3.6 & 12.8 & \na & \na & \na \\
HRNet-W48 & 65.8 & \na & 2.6 & 10.1 & \na & \na & \na \\
Swin-Unet & 27.2 & \na & 2.6 & 10.4 & \na & \na & \na \\
FCT & 1.9 & \na & 19.7 & 72.0 & \na & \na & \na \\
nnU-Net & \ph{33.4} & \na & 7.2 & 7.2 & \na & \na & \na \\
nnU-Net ResEnc-M & \ph{57.5} & \na & 12.8 & 11.4 & \na & \na & \na \\
M2MRF & 70.3 & \na & 3.0 & 11.6 & \na & \na & \na \\
HACDR-Net & 24.6 & \na & 5.9 & 27.8 & \na & \na & \na \\
HRDecoder & 65.6 & \na & 1.5 & 5.9 & \na & \na & \na \\
PMCNet & 12.9 & \na & 1.8 & 7.1 & \na & \na & \na \\
WSRFNet & 9.7 & \na & 2.7 & 10.5 & \na & \na & \na \\
CLC-Net & 119.7 & \na & 9.9 & 39.2 & \na & \na & \na \\
TC-Net & 14.9 & \na & 6.0 & 24.7 & \na & \na & \na \\
SAMed & 91.6 & \na & 2.1 & 7.4 & \na & \na & \na \\
SAM-Adapter & 641.3 & 4.3 & 6.5 & 26.3 & \na & \na & \na \\
SAM2-UNet & 216.5 & 4.4 & 4.1 & 16.2 & \na & \na & \na \\
MedSAM (full FT) & 93.7 & 93.7 & 2.8 & 11.1 & \na & \na & \na \\
RETFound + UPerNet & 352.6 & 352.6 & 6.1 & 24.2 & \na & \na & \na \\
\end{longtable}
}
\vspace{-6pt}\noindent{\scriptsize nnU-Net follows its own schedule of 1,000 epochs rather than the pixel budget. Cells without a value were not measured; the cost audit (FLOPs, memory, inference) is pending.}

\clearpage
\begin{landscape}
{\scriptsize\setlength{\tabcolsep}{3.2pt}\setlength{\LTcapwidth}{\textwidth}\renewcommand{\arraystretch}{1.08}
\begin{longtable}{l c c c c c c c c c c c}
\caption{\textbf{Refined IDRiD (label revision): per-class AUPR of every IDRiD-trained model against the official and the refined labels through the same 1,024-pixel geometry.} Test maps of the main-table checkpoints (seed 0), moved through the recovered crop and scored with the protocol evaluator against both label sets; no model is retrained or re-run.}\label{stab:refined}\\
\toprule
Model & MA off. & MA ref. & HE off. & HE ref. & EX off. & EX ref. & SE off. & SE ref. & \mAUPR{} off. & \mAUPR{} ref. & $\Delta$ \\
\midrule
\endfirsthead
\multicolumn{12}{l}{\textit{continued}}\\
\toprule
Model & MA off. & MA ref. & HE off. & HE ref. & EX off. & EX ref. & SE off. & SE ref. & \mAUPR{} off. & \mAUPR{} ref. & $\Delta$ \\
\midrule
\endhead
\midrule
\multicolumn{12}{r}{\textit{continued on the next page}}\\
\endfoot
\bottomrule
\endlastfoot
CLC-Net & .423 & .420 & .686 & .685 & .890 & .889 & .641 & .640 & .660 & .659 & $-$0.001 \\
DeepLabv3 & .091 & .091 & .643 & .643 & .736 & .736 & .664 & .664 & .533 & .533 & +0.000 \\
FCT & .079 & .078 & .165 & .165 & .687 & .686 & .115 & .115 & .261 & .261 & $-$0.000 \\
HACDR-Net & .358 & .356 & .504 & .504 & .828 & .827 & .561 & .560 & .563 & .561 & $-$0.001 \\
HRDecoder & .407 & .405 & .634 & .633 & .827 & .826 & .730 & .730 & .649 & .649 & $-$0.001 \\
HRNet-W48 & .442 & .438 & .685 & .684 & .885 & .883 & .783 & .783 & .699 & .697 & $-$0.002 \\
M2MRF & .439 & .436 & .625 & .625 & .879 & .878 & .681 & .680 & .656 & .655 & $-$0.001 \\
MedSAM (full FT) & .370 & .368 & .666 & .666 & .858 & .858 & .676 & .676 & .643 & .642 & $-$0.001 \\
PMCNet & .411 & .409 & .610 & .610 & .883 & .881 & .698 & .698 & .651 & .649 & $-$0.001 \\
RETFound + UPerNet & .483 & .480 & .719 & .718 & .896 & .895 & .736 & .736 & .709 & .707 & $-$0.001 \\
SAM2-UNet & .423 & .421 & .683 & .682 & .875 & .874 & .681 & .681 & .666 & .665 & $-$0.001 \\
SAM-Adapter & .442 & .439 & .745 & .745 & .887 & .887 & .756 & .756 & .708 & .707 & $-$0.001 \\
SAMed & .445 & .442 & .737 & .737 & .893 & .892 & .791 & .790 & .717 & .715 & $-$0.001 \\
Swin-Unet & .473 & .471 & .712 & .712 & .895 & .894 & .689 & .689 & .692 & .691 & $-$0.001 \\
TC-Net & .353 & .351 & .651 & .651 & .855 & .854 & .755 & .755 & .654 & .653 & $-$0.001 \\
U-Net & .386 & .384 & .623 & .622 & .895 & .894 & .710 & .711 & .653 & .653 & $-$0.001 \\
WSRFNet & .388 & .386 & .602 & .602 & .856 & .855 & .506 & .506 & .588 & .587 & $-$0.001 \\
nnU-Net & .482 & .480 & .673 & .673 & .878 & .876 & .645 & .644 & .669 & .668 & $-$0.001 \\
nnU-Net ResEnc-M & .442 & .440 & .667 & .666 & .878 & .877 & .541 & .540 & .632 & .631 & $-$0.001 \\
\end{longtable}
}
\vspace{-6pt}\noindent{\scriptsize Label differences over the 27 test images: MA: 13 images changed, 28,847 $\to$ 28,886 labelled pixels, 1,083 $\to$ 1,083 components; HE: 13 images changed, 311,289 $\to$ 311,524 labelled pixels, 532 $\to$ 532 components; EX: 16 images changed, 320,363 $\to$ 320,160 labelled pixels, 3,895 $\to$ 3,882 components; SE: 6 images changed, 53,684 $\to$ 53,597 labelled pixels, 38 $\to$ 38 components. The crop geometry of every image (all 27: padded to a square and resized; NCC $\geq$ 0.9988) is runs/REFINED/RefinedIDRiD\_geometry.json.}
\end{landscape}
\clearpage

{\scriptsize\setlength{\tabcolsep}{3.2pt}\setlength{\LTcapwidth}{\textwidth}\renewcommand{\arraystretch}{1.08}
\begin{longtable}{l l c c c c c}
\caption{\textbf{Encoder comparison: four ViT-L encoders under one UPerNet head, recipe and pixel budget (seed 0).} Per-class AUPR on the internal test and on the three independent fine-grained sets. RETFound-MAE is the main-table RETFound run of the same source; the other three are the control cells. The pre-specified contrasts are RETFound-MAE minus MAE-IN1k and RETFound-DINOv2 minus DINOv2.}\label{stab:encoders}\\
\toprule
Encoder & Test set & MA & HE & EX & SE & \mAUPR{} \\
\midrule
\endfirsthead
\multicolumn{7}{l}{\textit{continued}}\\
\toprule
Encoder & Test set & MA & HE & EX & SE & \mAUPR{} \\
\midrule
\endhead
\midrule
\multicolumn{7}{r}{\textit{continued on the next page}}\\
\endfoot
\bottomrule
\endlastfoot
\multicolumn{7}{l}{\textbf{Trained on IDRiD}}\\
RETFound-DINOv2 & IDRiD & .502 & .759 & .900 & .683 & .711 \\
RETFound-DINOv2 & TJDR std & .375 & .626 & .714 & .520 & .559 \\
RETFound-DINOv2 & TJDR UWF & .092 & .484 & .211 & .104 & .223 \\
RETFound-DINOv2 & MAPLES & .380 & .495 & .682 & .109 & .417 \\
RETFound-MAE & IDRiD & .484 & .718 & .895 & .734 & .708 \\
RETFound-MAE & TJDR std & .393 & .582 & .776 & .070 & .455 \\
RETFound-MAE & TJDR UWF & .099 & .494 & .106 & .311 & .253 \\
RETFound-MAE & MAPLES & .375 & .503 & .680 & .092 & .413 \\
DINOv2 ViT-L & IDRiD & .515 & .765 & .905 & .772 & .739 \\
DINOv2 ViT-L & TJDR std & .341 & .658 & .676 & .427 & .525 \\
DINOv2 ViT-L & TJDR UWF & .102 & .557 & .239 & .091 & .247 \\
DINOv2 ViT-L & MAPLES & .392 & .490 & .693 & .194 & .442 \\
MAE-IN1k ViT-L & IDRiD & .493 & .744 & .900 & .778 & .729 \\
MAE-IN1k ViT-L & TJDR std & .392 & .547 & .709 & .507 & .539 \\
MAE-IN1k ViT-L & TJDR UWF & .099 & .586 & .244 & .147 & .269 \\
MAE-IN1k ViT-L & MAPLES & .370 & .469 & .659 & .203 & .425 \\
\multicolumn{7}{l}{\textbf{Trained on DDR}}\\
RETFound-DINOv2 & DDR & .184 & .508 & .657 & .611 & .490 \\
RETFound-DINOv2 & TJDR std & .354 & .786 & .716 & .607 & .616 \\
RETFound-DINOv2 & TJDR UWF & .113 & .264 & .187 & .042 & .151 \\
RETFound-DINOv2 & MAPLES & .422 & .522 & .731 & .307 & .495 \\
RETFound-MAE & DDR & .189 & .491 & .640 & .546 & .466 \\
RETFound-MAE & TJDR std & .387 & .733 & .724 & .513 & .589 \\
RETFound-MAE & TJDR UWF & .073 & .064 & .278 & .030 & .111 \\
RETFound-MAE & MAPLES & .420 & .531 & .722 & .166 & .460 \\
DINOv2 ViT-L & DDR & .183 & .572 & .668 & .618 & .510 \\
DINOv2 ViT-L & TJDR std & .390 & .791 & .698 & .596 & .619 \\
DINOv2 ViT-L & TJDR UWF & .100 & .133 & .265 & .050 & .137 \\
DINOv2 ViT-L & MAPLES & .425 & .523 & .728 & .274 & .488 \\
MAE-IN1k ViT-L & DDR & .206 & .567 & .670 & .634 & .519 \\
MAE-IN1k ViT-L & TJDR std & .422 & .788 & .717 & .611 & .634 \\
MAE-IN1k ViT-L & TJDR UWF & .103 & .100 & .244 & .061 & .127 \\
MAE-IN1k ViT-L & MAPLES & .406 & .515 & .726 & .265 & .478 \\
\end{longtable}
}

\clearpage

{\scriptsize\setlength{\tabcolsep}{3.2pt}\setlength{\LTcapwidth}{\textwidth}\renewcommand{\arraystretch}{1.08}
\begin{longtable}{l c c c c c c c c }
\caption{\textbf{Ensembles of the main-table probability maps and the reference model, \mAUPR{} on every test set.} Compositions fixed before any ensemble was scored (Methods): all 18 registered members, the five and three best by validation \mAUPR{}, the best of each family, and each family alone. The reference model's two cells are trained on the three fine-grained sources merged; their ensemble is the pixel mean of their maps.}\label{stab:ensembles}\\
\toprule
Composition & IDRiD & DDR & TJDR std & TJDR UWF & MAPLES & FGADR & Retinal-L. & DiaRetDB1 \\
\midrule
\endfirsthead
\multicolumn{9}{l}{\textit{continued}}\\
\toprule
Composition & IDRiD & DDR & TJDR std & TJDR UWF & MAPLES & FGADR & Retinal-L. & DiaRetDB1 \\
\midrule
\endhead
\midrule
\multicolumn{9}{r}{\textit{continued on the next page}}\\
\endfoot
\bottomrule
\endlastfoot
\multicolumn{9}{l}{\textbf{Members trained on IDRiD}}\\
E-all & .764 & .400 & .594 & .311 & .455 & .279$^{\rm p}$ & .216$^{\rm p}$ & .317 \\
E-top5 & .754 & .410 & .598 & .292 & .449 & .290 & .251 & .316 \\
E-top3 & .742 & .394 & .559 & .279 & .437 & .273 & .252 & .325 \\
E-1perfam & .741 & .381 & .555 & .284 & .445 & .256 & .239 & .318 \\
E-fam:general & .727 & .352 & .513 & .276 & .410 & .244 & .194 & .302 \\
E-fam:auto & .693 & .316 & .485 & .229 & .406 & .242 & .177 & .269 \\
E-fam:fundus & .736 & .371 & .568 & .292 & .436 & .250$^{\rm p}$ & .214$^{\rm p}$ & .300 \\
E-fam:foundation & .757 & .411 & .609 & .297 & .442 & .294 & .250 & .318 \\
\multicolumn{9}{l}{\textbf{Members trained on DDR}}\\
E-all & .677 & .505 & .637 & .266 & .496 & .291$^{\rm p}$ & .278$^{\rm p}$ & .335 \\
E-top5 & .680 & .511 & .652 & .242 & .506 & .298 & .290 & .335 \\
E-top3 & .672 & .510 & .655 & .226 & .497 & .300 & .291 & .339 \\
E-1perfam & .665 & .494 & .625 & .228 & .486 & .281 & .278 & .321 \\
E-fam:general & .650 & .483 & .606 & .227 & .470 & .268 & .266 & .324 \\
E-fam:auto & .641 & .463 & .572 & .210 & .473 & .275 & .227 & .334 \\
E-fam:fundus & .637 & .444 & .568 & .246 & .452 & .251$^{\rm p}$ & .218$^{\rm p}$ & .305 \\
E-fam:foundation & .676 & .509 & .657 & .248 & .494 & .302 & .294 & .340 \\
\multicolumn{9}{l}{\textbf{Reference model: trained on IDRiD + DDR + TJDR std merged (DDR budget and validation)}}\\
RETFound + UPerNet (fine merged) & \na & .451 & \na & \na & \na & \na & \na & \na \\
SAMed (fine merged) & \na & .477 & \na & \na & \na & \na & \na & \na \\
\end{longtable}
}
\vspace{-6pt}\noindent{\scriptsize $^{\rm p}$, computed without HACDR-Net, whose maps for that set could not be regenerated (members named in the run record). In-source test sets of the reference model (IDRiD, DDR, TJDR std) are the official test splits of sources that entered its training.}

\clearpage

{\scriptsize\setlength{\tabcolsep}{3.2pt}\setlength{\LTcapwidth}{\textwidth}\renewcommand{\arraystretch}{1.08}
\begin{longtable}{l l c c c c c}
\caption{\textbf{Results by DR grade (seed 0).} On the test sets that carry a DR grade per image, \mAUPR{} and image-level AUC recomputed within grades 0--1 and within grades 2--4 from the stored per-image histograms, with the thresholds of the full test; ungradable images enter neither stratum.}\label{stab:strata}\\
\toprule
Model & Trained on & \mAUPR{} 0--1 & \mAUPR{} 2--4 & Image AUC 0--1 & Image AUC 2--4 & $\Delta$ \mAUPR{} \\
\midrule
\endfirsthead
\multicolumn{7}{l}{\textit{continued}}\\
\toprule
Model & Trained on & \mAUPR{} 0--1 & \mAUPR{} 2--4 & Image AUC 0--1 & Image AUC 2--4 & $\Delta$ \mAUPR{} \\
\midrule
\endhead
\midrule
\multicolumn{7}{r}{\textit{continued on the next page}}\\
\endfoot
\bottomrule
\endlastfoot
\multicolumn{7}{l}{\textbf{DDR: 20 images of grades 0--1, 203 of grades 2--4}}\\
CLC-Net & DDR & .254 & .384 & .653 & .752 & +0.129 \\
CLC-Net & IDRiD & .147 & .295 & .486 & .667 & +0.148 \\
DeepLabv3 & DDR & .145 & .372 & \na & \na & +0.227 \\
DeepLabv3 & IDRiD & .169 & .258 & .486 & .720 & +0.090 \\
FCT & DDR & .001 & .177 & \na & \na & +0.176 \\
FCT & IDRiD & .002 & .078 & .486 & .597 & +0.076 \\
HACDR-Net & DDR & .210 & .268 & \na & \na & +0.058 \\
HACDR-Net & IDRiD & .196 & .163 & .514 & .651 & $-$0.032 \\
HRDecoder & DDR & .266 & .273 & .639 & .692 & +0.007 \\
HRDecoder & IDRiD & .128 & .276 & .542 & .658 & +0.148 \\
HRNet-W48 & DDR & .260 & .442 & \na & \na & +0.182 \\
HRNet-W48 & IDRiD & .257 & .336 & .444 & .772 & +0.079 \\
M2MRF & DDR & .209 & .349 & \na & \na & +0.140 \\
M2MRF & IDRiD & .120 & .218 & .542 & .628 & +0.099 \\
MedSAM (full FT) & DDR & .225 & .361 & .444 & .782 & +0.137 \\
MedSAM (full FT) & IDRiD & .189 & .315 & .639 & .757 & +0.126 \\
PMCNet & DDR & .213 & .302 & .611 & .750 & +0.089 \\
PMCNet & IDRiD & .193 & .262 & .167 & .701 & +0.069 \\
RETFound + UPerNet & DDR & .256 & .462 & .597 & .829 & +0.206 \\
RETFound + UPerNet & IDRiD & .183 & .355 & .292 & .797 & +0.172 \\
SAM2-UNet & DDR & .270 & .315 & .514 & .741 & +0.044 \\
SAM2-UNet & IDRiD & .156 & .283 & .604 & .615 & +0.127 \\
SAM-Adapter & DDR & .273 & .416 & .583 & .639 & +0.143 \\
SAM-Adapter & IDRiD & .189 & .292 & .625 & .590 & +0.103 \\
SAMed & DDR & .227 & .466 & .361 & .832 & +0.240 \\
SAMed & IDRiD & .189 & .381 & .576 & .778 & +0.192 \\
Swin-Unet & DDR & .251 & .431 & \na & \na & +0.179 \\
Swin-Unet & IDRiD & .157 & .329 & .340 & .744 & +0.172 \\
TC-Net & DDR & .185 & .360 & .417 & .761 & +0.176 \\
TC-Net & IDRiD & .197 & .313 & .500 & .772 & +0.116 \\
U-Net & DDR & .206 & .341 & .458 & .776 & +0.135 \\
U-Net & IDRiD & .112 & .265 & .222 & .715 & +0.153 \\
WSRFNet & DDR & .160 & .217 & .569 & .637 & +0.057 \\
WSRFNet & IDRiD & .123 & .236 & .479 & .609 & +0.113 \\
nnU-Net & DDR & .150 & .410 & \na & \na & +0.260 \\
nnU-Net & IDRiD & .124 & .274 & .514 & .702 & +0.150 \\
nnU-Net ResEnc-M & DDR & .151 & .403 & \na & \na & +0.252 \\
nnU-Net ResEnc-M & IDRiD & .150 & .249 & .542 & .750 & +0.098 \\
\multicolumn{7}{l}{\textbf{MAPLES: 56 images of grades 0--1, 4 of grades 2--4}}\\
CLC-Net & DDR & .297 & .308 & .650 & .667 & +0.011 \\
CLC-Net & IDRiD & .349 & .445 & .638 & 1.000 & +0.096 \\
DeepLabv3 & DDR & .332 & .366 & .752 & 1.000 & +0.034 \\
DeepLabv3 & IDRiD & .261 & .425 & .693 & 1.000 & +0.164 \\
FCT & DDR & .127 & .186 & .470 & .333 & +0.060 \\
FCT & IDRiD & .107 & .177 & .521 & 1.000 & +0.069 \\
HACDR-Net & DDR & .290 & .371 & .613 & 1.000 & +0.081 \\
HACDR-Net & IDRiD & .265 & .407 & .635 & 1.000 & +0.143 \\
HRDecoder & DDR & .265 & .284 & .761 & .833 & +0.019 \\
HRDecoder & IDRiD & .338 & .436 & .650 & 1.000 & +0.098 \\
HRNet-W48 & DDR & .433 & .459 & .788 & .667 & +0.027 \\
HRNet-W48 & IDRiD & .379 & .465 & .638 & 1.000 & +0.086 \\
M2MRF & DDR & .292 & .329 & .690 & 1.000 & +0.037 \\
M2MRF & IDRiD & .280 & .330 & .638 & 1.000 & +0.050 \\
MedSAM (full FT) & DDR & .373 & .425 & .743 & .667 & +0.051 \\
MedSAM (full FT) & IDRiD & .337 & .478 & .738 & 1.000 & +0.142 \\
PMCNet & DDR & .370 & .434 & .735 & .667 & +0.064 \\
PMCNet & IDRiD & .325 & .436 & .679 & 1.000 & +0.110 \\
RETFound + UPerNet & DDR & .437 & .494 & .830 & .667 & +0.056 \\
RETFound + UPerNet & IDRiD & .386 & .509 & .721 & .667 & +0.123 \\
SAM2-UNet & DDR & .354 & .403 & .675 & .833 & +0.049 \\
SAM2-UNet & IDRiD & .332 & .391 & .658 & .833 & +0.059 \\
SAM-Adapter & DDR & .325 & .358 & .657 & .833 & +0.032 \\
SAM-Adapter & IDRiD & .334 & .350 & .677 & 1.000 & +0.016 \\
SAMed & DDR & .454 & .471 & .807 & 1.000 & +0.018 \\
SAMed & IDRiD & .393 & .495 & .792 & 1.000 & +0.101 \\
Swin-Unet & DDR & .435 & .548 & .770 & 1.000 & +0.113 \\
Swin-Unet & IDRiD & .374 & .456 & .622 & .667 & +0.083 \\
TC-Net & DDR & .324 & .432 & .735 & .667 & +0.108 \\
TC-Net & IDRiD & .383 & .522 & .701 & 1.000 & +0.138 \\
U-Net & DDR & .262 & .375 & .699 & 1.000 & +0.114 \\
U-Net & IDRiD & .277 & .374 & .681 & 1.000 & +0.097 \\
WSRFNet & DDR & .257 & .327 & .669 & .667 & +0.071 \\
WSRFNet & IDRiD & .264 & .319 & .630 & .333 & +0.054 \\
nnU-Net & DDR & .343 & .427 & .818 & 1.000 & +0.084 \\
nnU-Net & IDRiD & .338 & .420 & .700 & 1.000 & +0.082 \\
nnU-Net ResEnc-M & DDR & .375 & .423 & .797 & 1.000 & +0.048 \\
nnU-Net ResEnc-M & IDRiD & .301 & .350 & .696 & .333 & +0.050 \\
\multicolumn{7}{l}{\textbf{FGADR: 94 images of grades 0--1, 459 of grades 2--4}}\\
CLC-Net & DDR & .059 & .166 & .560 & .747 & +0.107 \\
CLC-Net & IDRiD & .016 & .172 & .633 & .698 & +0.156 \\
DeepLabv3 & DDR & .056 & .236 & .762 & .753 & +0.179 \\
DeepLabv3 & IDRiD & .055 & .222 & .661 & .762 & +0.167 \\
FCT & DDR & .003 & .071 & .571 & .602 & +0.068 \\
FCT & IDRiD & .002 & .052 & .577 & .634 & +0.050 \\
HACDR-Net & DDR & .012 & .106 & .443 & .674 & +0.095 \\
HACDR-Net & IDRiD & .005 & .038 & .501 & .573 & +0.032 \\
HRDecoder & DDR & .021 & .188 & .707 & .727 & +0.168 \\
HRDecoder & IDRiD & .024 & .199 & .650 & .715 & +0.176 \\
HRNet-W48 & DDR & .046 & .274 & .775 & .814 & +0.229 \\
HRNet-W48 & IDRiD & .036 & .256 & .759 & .802 & +0.220 \\
M2MRF & DDR & .037 & .202 & .714 & .731 & +0.165 \\
M2MRF & IDRiD & .002 & .084 & .581 & .599 & +0.082 \\
MedSAM (full FT) & DDR & .074 & .231 & .768 & .803 & +0.157 \\
MedSAM (full FT) & IDRiD & .031 & .270 & .635 & .759 & +0.238 \\
PMCNet & DDR & .036 & .211 & .610 & .793 & +0.175 \\
PMCNet & IDRiD & .027 & .216 & .730 & .752 & +0.189 \\
RETFound + UPerNet & DDR & .061 & .268 & .700 & .814 & +0.207 \\
RETFound + UPerNet & IDRiD & .017 & .240 & .719 & .773 & +0.222 \\
SAM2-UNet & DDR & .028 & .222 & .712 & .754 & +0.195 \\
SAM2-UNet & IDRiD & .037 & .235 & .734 & .759 & +0.198 \\
SAM-Adapter & DDR & .035 & .242 & .687 & .669 & +0.207 \\
SAM-Adapter & IDRiD & .014 & .221 & .564 & .620 & +0.207 \\
SAMed & DDR & .067 & .295 & .770 & .832 & +0.228 \\
SAMed & IDRiD & .053 & .289 & .745 & .800 & +0.236 \\
Swin-Unet & DDR & .043 & .248 & .700 & .802 & +0.204 \\
Swin-Unet & IDRiD & .017 & .204 & .590 & .726 & +0.187 \\
TC-Net & DDR & .039 & .219 & .735 & .795 & +0.180 \\
TC-Net & IDRiD & .020 & .221 & .640 & .782 & +0.201 \\
U-Net & DDR & .019 & .170 & .690 & .761 & +0.151 \\
U-Net & IDRiD & .003 & .144 & .576 & .700 & +0.141 \\
WSRFNet & DDR & .016 & .062 & .652 & .639 & +0.046 \\
WSRFNet & IDRiD & .016 & .187 & .642 & .651 & +0.170 \\
nnU-Net & DDR & .046 & .249 & .728 & .773 & +0.203 \\
nnU-Net & IDRiD & .036 & .222 & .673 & .784 & +0.186 \\
nnU-Net ResEnc-M & DDR & .056 & .245 & .794 & .810 & +0.190 \\
nnU-Net ResEnc-M & IDRiD & .027 & .189 & .615 & .773 & +0.162 \\
\multicolumn{7}{l}{\textbf{Retinal-L.: 150 images of grades 0--1, 328 of grades 2--4}}\\
CLC-Net & DDR & .057 & .125 & .762 & .807 & +0.068 \\
CLC-Net & IDRiD & .029 & .168 & .807 & .768 & +0.139 \\
DeepLabv3 & DDR & .097 & .273 & .730 & .798 & +0.176 \\
DeepLabv3 & IDRiD & .051 & .286 & .695 & .746 & +0.236 \\
FCT & DDR & .003 & .085 & .612 & .605 & +0.081 \\
FCT & IDRiD & .002 & .035 & .529 & .570 & +0.033 \\
HACDR-Net & DDR & .038 & .140 & .670 & .782 & +0.103 \\
HACDR-Net & IDRiD & .012 & .086 & .684 & .662 & +0.074 \\
HRDecoder & DDR & .050 & .116 & .849 & .793 & +0.066 \\
HRDecoder & IDRiD & .030 & .130 & .767 & .760 & +0.100 \\
HRNet-W48 & DDR & .112 & .264 & .792 & .846 & +0.152 \\
HRNet-W48 & IDRiD & .061 & .217 & .780 & .814 & +0.156 \\
M2MRF & DDR & .068 & .135 & .872 & .831 & +0.067 \\
M2MRF & IDRiD & .023 & .115 & .744 & .727 & +0.092 \\
MedSAM (full FT) & DDR & .058 & .225 & .834 & .839 & +0.167 \\
MedSAM (full FT) & IDRiD & .042 & .240 & .790 & .786 & +0.197 \\
PMCNet & DDR & .055 & .179 & .808 & .833 & +0.125 \\
PMCNet & IDRiD & .039 & .185 & .740 & .754 & +0.146 \\
RETFound + UPerNet & DDR & .121 & .248 & .826 & .856 & +0.127 \\
RETFound + UPerNet & IDRiD & .040 & .221 & .774 & .833 & +0.182 \\
SAM2-UNet & DDR & .081 & .194 & .807 & .777 & +0.113 \\
SAM2-UNet & IDRiD & .034 & .168 & .754 & .730 & +0.134 \\
SAM-Adapter & DDR & .074 & .169 & .789 & .730 & +0.094 \\
SAM-Adapter & IDRiD & .020 & .170 & .612 & .694 & +0.151 \\
SAMed & DDR & .125 & .281 & .828 & .864 & +0.155 \\
SAMed & IDRiD & .055 & .268 & .795 & .806 & +0.213 \\
Swin-Unet & DDR & .100 & .263 & .846 & .870 & +0.163 \\
Swin-Unet & IDRiD & .052 & .247 & .810 & .822 & +0.195 \\
TC-Net & DDR & .034 & .182 & .752 & .808 & +0.148 \\
TC-Net & IDRiD & .026 & .193 & .768 & .805 & +0.167 \\
U-Net & DDR & .069 & .123 & .850 & .840 & +0.054 \\
U-Net & IDRiD & .028 & .126 & .778 & .761 & +0.098 \\
WSRFNet & DDR & .032 & .070 & .705 & .708 & +0.037 \\
WSRFNet & IDRiD & .036 & .124 & .756 & .695 & +0.087 \\
nnU-Net & DDR & .073 & .165 & .858 & .837 & +0.092 \\
nnU-Net & IDRiD & .028 & .161 & .835 & .797 & +0.132 \\
nnU-Net ResEnc-M & DDR & .089 & .172 & .781 & .873 & +0.082 \\
nnU-Net ResEnc-M & IDRiD & .020 & .126 & .798 & .821 & +0.106 \\
\end{longtable}
}
\vspace{-6pt}\noindent{\scriptsize Image-level AUC within a stratum is undefined when every image of the stratum carries the class or none does; such cells read 0 or 1 in the record and should be read with the stratum size.}

\clearpage

{\scriptsize\setlength{\tabcolsep}{3.2pt}\setlength{\LTcapwidth}{\textwidth}\renewcommand{\arraystretch}{1.08}
\begin{longtable}{l c c c c c c c}
\caption{\textbf{Human reference on DiaRetDB1: each expert against the consensus of the other three.} The expert's own markings, with their confidence levels mapped to scores, are the probability map and the consensus of the other three experts ($\geq$ 0.75) is the ground truth; thresholds come from the official training list. A map with four confidence levels gives a four-point precision--recall curve, so the pixel-level AUPR is a lower bound.}\label{stab:experts}\\
\toprule
Reader & MA & HE & EX & SE & \mAUPR{} & Image AUC & Dice \\
\midrule
\endfirsthead
\multicolumn{8}{l}{\textit{continued}}\\
\toprule
Reader & MA & HE & EX & SE & \mAUPR{} & Image AUC & Dice \\
\midrule
\endhead
\midrule
\multicolumn{8}{r}{\textit{continued on the next page}}\\
\endfoot
\bottomrule
\endlastfoot
Expert 01 & .008 & .246 & .414 & .102 & .193 & .870 & .349 \\
Expert 02 & .001 & .166 & .273 & .162 & .150 & .931 & .261 \\
Expert 03 & .015 & .341 & .516 & .244 & .279 & .931 & .448 \\
Expert 04 & .022 & .271 & .366 & .220 & .219 & .913 & .360 \\
\multicolumn{8}{l}{\textit{Models trained on IDRiD, against the four-expert consensus (median over 19 models; range in brackets)}}\\
models & \na & \na & \na & \na & 0.245 [0.10--0.32] & 0.833 [0.71--0.92] & \na \\
\multicolumn{8}{l}{\textit{Models trained on DDR, against the four-expert consensus (median over 19 models; range in brackets)}}\\
models & \na & \na & \na & \na & 0.284 [0.09--0.32] & 0.840 [0.74--0.94] & \na \\
\end{longtable}
}

\clearpage

{\scriptsize\setlength{\tabcolsep}{3.2pt}\setlength{\LTcapwidth}{\textwidth}\renewcommand{\arraystretch}{1.08}
\begin{longtable}{l l r r c c r c c c r}
\caption{\textbf{Annotation style of every source: connected components per image and their size.} All images of each source, every id once; 8-connected components of the prepared masks with ignore regions removed; the relative diameter is the component's equivalent diameter divided by the long side of the field-of-view box, in percent (Methods, annotation style).}\label{stab:labels}\\
\toprule
Source & Class & Images & Components & Per image & Per positive image & Positive images & Diam.\ p10 (\%) & p50 (\%) & p90 (\%) & Area px (median) \\
\midrule
\endfirsthead
\multicolumn{11}{l}{\textit{continued}}\\
\toprule
Source & Class & Images & Components & Per image & Per positive image & Positive images & Diam.\ p10 (\%) & p50 (\%) & p90 (\%) & Area px (median) \\
\midrule
\endhead
\midrule
\multicolumn{11}{r}{\textit{continued on the next page}}\\
\endfoot
\bottomrule
\endlastfoot
IDRiD & MA & 81 & 3,497 & 43.2 & 43.2 & 81 & 0.33 & 0.50 & 0.78 & 234 \\
IDRiD & HE & 81 & 1,900 & 23.5 & 23.8 & 80 & 0.70 & 1.32 & 3.29 & 1608 \\
IDRiD & EX & 81 & 11,642 & 143.7 & 143.7 & 81 & 0.26 & 0.48 & 1.20 & 217 \\
IDRiD & SE & 81 & 150 & 1.9 & 3.8 & 40 & 1.86 & 3.07 & 5.45 & 8653 \\
DDR & MA & 757 & 10,380 & 13.7 & 18.2 & 570 & 0.21 & 0.43 & 0.77 & 48 \\
DDR & HE & 757 & 12,538 & 16.6 & 20.9 & 601 & 0.27 & 0.87 & 2.36 & 209 \\
DDR & EX & 757 & 23,659 & 31.3 & 48.7 & 486 & 0.16 & 0.45 & 1.31 & 56 \\
DDR & SE & 757 & 1,293 & 1.7 & 5.4 & 239 & 0.24 & 1.91 & 3.83 & 855 \\
TJDR\_std & MA & 257 & 742 & 2.9 & 7.5 & 99 & 0.41 & 0.80 & 1.16 & 210 \\
TJDR\_std & HE & 257 & 1,192 & 4.6 & 8.1 & 148 & 0.94 & 1.78 & 3.85 & 1048 \\
TJDR\_std & EX & 257 & 1,130 & 4.4 & 11.0 & 103 & 0.62 & 1.26 & 2.76 & 525 \\
TJDR\_std & SE & 257 & 216 & 0.8 & 3.3 & 65 & 1.36 & 3.11 & 6.19 & 3178 \\
TJDR\_uwf & MA & 304 & 765 & 2.5 & 10.2 & 75 & 0.21 & 0.35 & 0.59 & 148 \\
TJDR\_uwf & HE & 304 & 2,012 & 6.6 & 11.8 & 170 & 0.35 & 0.84 & 2.14 & 841 \\
TJDR\_uwf & EX & 304 & 2,222 & 7.3 & 10.1 & 219 & 0.15 & 0.40 & 1.19 & 188 \\
TJDR\_uwf & SE & 304 & 287 & 0.9 & 2.3 & 127 & 0.28 & 1.53 & 2.64 & 2826 \\
FGADR & MA & 1,842 & 11,235 & 6.1 & 7.9 & 1424 & 0.46 & 0.78 & 3.05 & 79 \\
FGADR & HE & 1,842 & 37,958 & 20.6 & 26.1 & 1456 & 0.59 & 1.17 & 3.45 & 176 \\
FGADR & EX & 1,842 & 52,896 & 28.7 & 41.4 & 1279 & 0.37 & 0.76 & 1.77 & 75 \\
FGADR & SE & 1,842 & 2,087 & 1.1 & 3.3 & 627 & 1.09 & 3.21 & 6.40 & 1323 \\
Retinal-Lesions & MA & 1,593 & 14,873 & 9.3 & 11.0 & 1346 & 1.37 & 2.31 & 4.10 & 336 \\
Retinal-Lesions & HE & 1,593 & 13,305 & 8.4 & 12.5 & 1064 & 1.82 & 3.45 & 7.12 & 749 \\
Retinal-Lesions & EX & 1,593 & 3,416 & 2.1 & 4.2 & 807 & 1.87 & 4.16 & 12.49 & 1089 \\
Retinal-Lesions & SE & 1,593 & 1,267 & 0.8 & 2.7 & 471 & 2.94 & 5.08 & 8.81 & 1624 \\
DiaRetDB1 & MA & 89 & 183 & 2.1 & 4.6 & 40 & 0.48 & 0.62 & 0.72 & 61 \\
DiaRetDB1 & HE & 89 & 402 & 4.5 & 10.1 & 40 & 0.71 & 2.93 & 8.61 & 1363 \\
DiaRetDB1 & EX & 89 & 137 & 1.5 & 3.6 & 38 & 0.71 & 4.06 & 16.01 & 2624 \\
DiaRetDB1 & SE & 89 & 27 & 0.3 & 1.8 & 15 & 1.90 & 4.31 & 8.02 & 2945 \\
MAPLES-DR & MA & 198 & 5,412 & 27.3 & 28.3 & 191 & 0.63 & 0.84 & 1.29 & 74 \\
MAPLES-DR & HE & 198 & 946 & 4.8 & 7.8 & 122 & 1.06 & 1.89 & 3.19 & 446 \\
MAPLES-DR & EX & 198 & 3,080 & 15.6 & 29.6 & 104 & 0.65 & 0.99 & 2.07 & 122 \\
MAPLES-DR & SE & 198 & 199 & 1.0 & 3.3 & 61 & 1.29 & 2.48 & 4.29 & 620 \\
\end{longtable}
}

\clearpage

{\scriptsize\setlength{\tabcolsep}{3.2pt}\setlength{\LTcapwidth}{\textwidth}\renewcommand{\arraystretch}{1.08}
\begin{longtable}{l l r}
\caption{\textbf{Near-duplicate audit between training sources and test sets, and the images excluded.} Vessel signatures (green-channel black top-hat at 192 $\times$ 192 inside the field of view) are correlated between every train or validation image of a source and every test image; a test image at or above the pair's threshold is left out of that source's independent tests (default 0.5; 0.33 within FGADR, whose images are near-duplicates of each other at scale).}\label{stab:leaks}\\
\toprule
Training source & Test set & Excluded images \\
\midrule
\endfirsthead
\multicolumn{3}{l}{\textit{continued}}\\
\toprule
Training source & Test set & Excluded images \\
\midrule
\endhead
\midrule
\multicolumn{3}{r}{\textit{continued on the next page}}\\
\endfoot
\bottomrule
\endlastfoot
TJDR\_uwf & TJDR\_uwf & 2 \\
FGADR & FGADR & 106 \\
FGADR & Retinal-Lesions & 19 \\
Retinal-Lesions & FGADR & 27 \\
Retinal-Lesions & Retinal-Lesions & 1 \\
IDRiD & every test set & 0 \\
DDR & every test set & 0 \\
TJDR\_std & every test set & 0 \\
\end{longtable}
}
\vspace{-6pt}\noindent{\scriptsize Audit: bench/scripts/leak\_audit.py --min-vessel 0.3 (all seven prepared sources), 4923 images. Strongest pairs (vessel NCC): FGADR train:0224\_1 $\leftrightarrow$ FGADR 0271\_1 (1.00); FGADR train:0322\_1 $\leftrightarrow$ FGADR 0741\_1 (1.00); FGADR train:0526\_1 $\leftrightarrow$ FGADR 1597\_1 (1.00); FGADR train:0878\_1 $\leftrightarrow$ FGADR 0988\_1 (1.00); FGADR val:0553\_1 $\leftrightarrow$ FGADR 0995\_1 (1.00); TJDR\_uwf train:TJDR\_train\_092 $\leftrightarrow$ TJDR\_uwf TJDR\_test\_021 (1.00); TJDR\_uwf val:TJDR\_train\_101 $\leftrightarrow$ TJDR\_uwf TJDR\_test\_023 (1.00); Retinal-Lesions train:4428\_left $\leftrightarrow$ FGADR 0706\_3 (0.99); Retinal-Lesions val:31746\_right $\leftrightarrow$ FGADR 0466\_3 (0.98); Retinal-Lesions train:24235\_right $\leftrightarrow$ FGADR 0785\_3 (0.98); Retinal-Lesions train:658\_right $\leftrightarrow$ FGADR 0236\_3 (0.98); Retinal-Lesions train:15684\_right $\leftrightarrow$ FGADR 0627\_3 (0.98). The full lists are bench/data/splits/leak\_exclusions.json.}

\clearpage

{\scriptsize\setlength{\tabcolsep}{3.2pt}\setlength{\LTcapwidth}{\textwidth}\renewcommand{\arraystretch}{1.08}
\begin{longtable}{l l l c c c c c}
\caption{\textbf{Released checkpoints on every test set, per class, under this protocol.} Each checkpoint is run through its authors' own test pipeline (resize, normalisation, whole image or sliding window as configured) and the maps are scored with the protocol evaluator; a source the checkpoint was trained on is its internal test.}\label{stab:authors}\\
\toprule
Checkpoint & Trained on & Test set & MA & HE & EX & SE & \mAUPR{} \\
\midrule
\endfirsthead
\multicolumn{8}{l}{\textit{continued}}\\
\toprule
Checkpoint & Trained on & Test set & MA & HE & EX & SE & \mAUPR{} \\
\midrule
\endhead
\midrule
\multicolumn{8}{r}{\textit{continued on the next page}}\\
\endfoot
\bottomrule
\endlastfoot
M2MRF & IDRiD & IDRiD & .408 & .654 & .832 & .664 & .639 \\
M2MRF & IDRiD & DDR & .122 & .289 & .506 & .061 & .244 \\
M2MRF & IDRiD & TJDR std & .294 & .387 & .566 & .241 & .372 \\
M2MRF & IDRiD & TJDR UWF & .047 & .066 & .044 & .096 & .063 \\
M2MRF & IDRiD & MAPLES & .220 & .363 & .457 & .048 & .272 \\
M2MRF & IDRiD & FGADR & .010 & .243 & .210 & .113 & .144 \\
M2MRF & IDRiD & Retinal-L. & .033 & .136 & .179 & .068 & .104 \\
M2MRF & IDRiD & DiaRetDB1 & .043 & .247 & .265 & .107 & .165 \\
M2MRF & DDR & IDRiD & .276 & .549 & .678 & .573 & .519 \\
M2MRF & DDR & DDR & .165 & .426 & .558 & .456 & .401 \\
M2MRF & DDR & TJDR std & .300 & .628 & .664 & .479 & .518 \\
M2MRF & DDR & TJDR UWF & .073 & .512 & .164 & .075 & .206 \\
M2MRF & DDR & MAPLES & .204 & .455 & .429 & .102 & .297 \\
M2MRF & DDR & FGADR & .015 & .263 & .312 & .163 & .188 \\
M2MRF & DDR & Retinal-L. & .044 & .148 & .195 & .188 & .144 \\
M2MRF & DDR & DiaRetDB1 & .055 & .234 & .178 & .247 & .179 \\
HRDecoder & IDRiD & IDRiD & .475 & .699 & .892 & .763 & .707 \\
HRDecoder & IDRiD & DDR & .110 & .355 & .556 & .148 & .292 \\
HRDecoder & IDRiD & TJDR std & .298 & .578 & .610 & .426 & .478 \\
HRDecoder & IDRiD & TJDR UWF & .075 & .413 & .161 & .359 & .252 \\
HRDecoder & IDRiD & MAPLES & .287 & .444 & .692 & .083 & .377 \\
HRDecoder & IDRiD & FGADR & .016 & .285 & .338 & .102 & .186 \\
HRDecoder & IDRiD & Retinal-L. & .038 & .157 & .338 & .125 & .165 \\
HRDecoder & IDRiD & DiaRetDB1 & .047 & .329 & .452 & .175 & .251 \\
HRDecoder & DDR & IDRiD & .291 & .639 & .799 & .633 & .590 \\
HRDecoder & DDR & DDR & .172 & .557 & .629 & .486 & .461 \\
HRDecoder & DDR & TJDR std & .298 & .707 & .693 & .517 & .554 \\
HRDecoder & DDR & TJDR UWF & .077 & .543 & .148 & .102 & .218 \\
HRDecoder & DDR & MAPLES & .279 & .471 & .617 & .159 & .382 \\
HRDecoder & DDR & FGADR & .013 & .302 & .291 & .211 & .204 \\
HRDecoder & DDR & Retinal-L. & .064 & .216 & .283 & .232 & .199 \\
HRDecoder & DDR & DiaRetDB1 & .060 & .409 & .323 & .306 & .275 \\
\end{longtable}
}

\clearpage
\begin{landscape}
{\scriptsize\setlength{\tabcolsep}{3.2pt}\setlength{\LTcapwidth}{\textwidth}\renewcommand{\arraystretch}{1.08}
\begin{longtable}{l l c c c c c c c c }
\caption{\textbf{Foundation models without training, on every test set.} Zero-shot automatic masks (density and area cap chosen on the validation pool), SAM 3 text prompts with the wording chosen on the validation pool, and prompts derived from the ground truth (an upper bound that reads the test labels). Values are \mAUPR{} where the output carries classes and the any-lesion AUPR$^{\rm a}$ where it does not.}\label{stab:inference}\\
\toprule
Model & Setting & IDRiD & DDR & TJDR std & TJDR UWF & MAPLES & FGADR & Retinal-L. & DiaRetDB1 \\
\midrule
\endfirsthead
\multicolumn{10}{l}{\textit{continued}}\\
\toprule
Model & Setting & IDRiD & DDR & TJDR std & TJDR UWF & MAPLES & FGADR & Retinal-L. & DiaRetDB1 \\
\midrule
\endhead
\midrule
\multicolumn{10}{r}{\textit{continued on the next page}}\\
\endfoot
\bottomrule
\endlastfoot
\multicolumn{10}{l}{\textbf{Zero-shot (any-lesion AUPR of the class-agnostic map; text prompts per class)}}\\
SAM 2 & zero\_shot:auto\_any\_lesion & .023 & .008 & .012 & .007 & .003 & .010 & .014 & .010 \\
SAM 2 & zero\_shot:auto\_oracle\_class & .029 & .009 & .012 & .007 & .003 & .010 & .014 & .010 \\
SAM 2 & zero\_shot:auto\_stability & .026 & .011 & .008 & .018 & .004 & .011 & .014 & .008 \\
SAM & zero\_shot:auto\_any\_lesion & .020 & .008 & .008 & .004 & .003 & .012 & .014 & .007 \\
SAM & zero\_shot:auto\_oracle\_class & .022 & .008 & .009 & .005 & .003 & .013 & .014 & .007 \\
SAM & zero\_shot:auto\_stability & .031 & .014 & .028 & .015 & .004 & .020 & .015 & .009 \\
\multicolumn{10}{l}{\textbf{Oracle boxes, every component}}\\
MedSAM & oracle:box:all & .665 & .648 & .840 & .860 & .738 & .827 & .864 & .714 \\
SAM 2 & oracle:box:all & .848 & .756 & .932 & .922 & .826 & .825 & .835 & .771 \\
SAM & oracle:box:all & .860 & .772 & .870 & .869 & .824 & .759 & .747 & .777 \\
\multicolumn{10}{l}{\textbf{Oracle points, one per component}}\\
SAM 2 & oracle:point:all & .116 & .019 & .081 & .291 & .009 & .050 & .058 & .074 \\
\multicolumn{10}{l}{\textbf{SAM 3, selected wording}}\\
SAM 3 & text:selected & .228 & .125 & .111 & .185 & .079 & .108 & .066 & .088 \\
\end{longtable}
}
\vspace{-6pt}\noindent{\scriptsize $^{\rm a}$ any-lesion AUPR of a class-agnostic map. Variant tokens as in the run records (density, cap, canvas).}
\end{landscape}
\clearpage

{\scriptsize\setlength{\tabcolsep}{3.2pt}\setlength{\LTcapwidth}{\textwidth}\renewcommand{\arraystretch}{1.08}
\begin{longtable}{l p{40mm} p{66mm} p{38mm}}
\caption{\textbf{Experiment names used in the text and the identifiers used in the repository.} The run directories under runs/, the experiment column of data/results.csv and the planning documents in docs/ carry the identifiers; the text names each experiment by what it does.}\label{stab:identifiers}\\
\toprule
Identifier & Name in the text & What it is & Where it is reported \\
\midrule
\endfirsthead
\multicolumn{4}{l}{\textit{continued}}\\
\toprule
Identifier & Name in the text & What it is & Where it is reported \\
\midrule
\endhead
\midrule
\multicolumn{4}{r}{\textit{continued on the next page}}\\
\endfoot
\bottomrule
\endlastfoot
\texttt{L1} & Paper audit & Extraction of the setup and the comparison tables of every method paper in the ledger & Fig.~2b,c; Supplementary Tables~1 and 2 \\
\texttt{L2} & Reproducibility funnel & Registered fundus-specific methods by code, training code, weights and reproduced number & Fig.~2a \\
\texttt{E1} & Released checkpoints (author weights) & Released checkpoints scored with the paper's number, the authors' function and this protocol, on all test sets & Fig.~3a,b,d; Table~3a; Supplementary Table~16 \\
\texttt{E1r} & Author-recipe reruns & The authors' own pipelines retrained from scratch on IDRiD (M2MRF, HRDecoder) & Fig.~3c; Table~3b \\
\texttt{E2 (B-std)} & Unified retraining; the unified protocol & Every tier-1 model with complete training code retrained on IDRiD and on DDR under the shared protocol & Fig.~4; Tables~4 and 5 \\
\texttt{E2b (B1)} & Source-recipe track & Each model trained with its source recipe on the shared data; stopped before any run finished & Methods only \\
\texttt{E2c} & Learning-rate sweep & HRNet-W48 and M2MRF on DDR with the recipe learning rate $\times$1/3 and $\times$3, judged on validation only & Extended Data Fig.~7 \\
\texttt{E3a} & Independent tests & Every main-table checkpoint applied to the six test sets of sources never used for training & Fig.~4b--e; Table~5; Extended Data Fig.~9 \\
\texttt{E3b} & Training-source experiment & Five representative models trained on the standard-field TJDR images and on FGADR & Fig.~6b--d; Supplementary Table~7 \\
\texttt{E4} & Annotation-granularity experiment & U-Net and M2MRF on fine, mixed and coarse training labels, alone and combined & Fig.~6e; Table~7; Supplementary Table~7 \\
\texttt{E5} & Foundation models under four kinds of use & Zero-shot, oracle prompts, parameter-efficient fine-tuning, encoder + decoder; the encoder comparison & Fig.~5; Table~6; Supplementary Tables~10 and 17 \\
\texttt{E7} & Input-resolution experiment & M2MRF and RETFound + UPerNet at a 512- or 1,440-pixel canvas and on native-resolution patches & Fig.~7d; Table~7 \\
\texttt{E8} & Training-amount experiment & U-Net and M2MRF on 10\%, 50\% and 100\% of DDR and on the merged fine-grained sources & Fig.~7e; Table~7 \\
\texttt{E9} & Loss experiment & U-Net on DDR with four losses, architecture fixed & Extended Data Fig.~8; Table~7 \\
\texttt{E10} & Computational cost & Parameters, trainable parameters and GPU hours of every run & Extended Data Fig.~10; Supplementary Table~8 \\
\texttt{ENS} & Ensembles & Pixelwise means of the probability maps in the registered compositions & Table~5; Supplementary Table~11 \\
\texttt{REF} & Label revision (Refined IDRiD) & IDRiD-trained checkpoints re-scored against the refined and the official labels & Supplementary Table~9 \\
\texttt{HUMAN} & Human reference & Each DiaRetDB1 expert against the consensus of the other three & Table~5; Supplementary Table~13 \\
\texttt{CLIN} & Results by DR grade & Every score recomputed within grades 0--1 and 2--4 & Supplementary Table~12 \\
\texttt{N8} & Reference model & Two models trained on the merged fine-grained sources and averaged; in training & Methods \\
\end{longtable}
}


\clearpage
\section*{Supplementary Data}
\begin{itemize}[leftmargin=*,itemsep=2pt]
\item \textbf{Supplementary Data 1.} Split lists of every source (\path{bench/data/splits/} in the repository): the image identifiers of each train, validation and test split, the subset lists of the training-amount experiment, and the near-duplicate exclusions.
\item \textbf{Supplementary Data 2.} The literature extraction (\path{literature/meta/reported_results.tsv}): one row per paper, baseline, dataset, class and metric, with the stated provenance of every number, and the per-paper audit (\path{paper_audit.tsv}).
\item \textbf{Supplementary Data 3.} The results table behind every figure and table of the main text and of this document (\path{docs/paper/retiseg_paper/paper/data/results.csv}), with the recipe, environment and deviation record of every run; its experiment column uses the identifiers of Supplementary Table~18.
\item \textbf{Supplementary Data 4.} Probability maps of every main-table model on every test set whose licence allows derived data to be released (IDRiD, DDR, TJDR, MAPLES-DR, DiaRetDB1), as 16-bit images, so that any metric can be recomputed.
\end{itemize}

\end{document}
